\documentclass[11pt]{article}
\usepackage{amsmath,amssymb,amsthm}
\usepackage[letterpaper,margin=1in,headheight=15pt,headsep=20pt]{geometry}
\usepackage[T1]{fontenc}
\usepackage{lmodern}
\usepackage{mathtools}
\usepackage{enumitem}
\usepackage{caption}
\usepackage{fancyhdr}
\usepackage{needspace}
\usepackage{xcolor}
\usepackage{placeins}
\usepackage{etoolbox}
\usepackage{microtype}
\usepackage{booktabs}
\usepackage{array}
\usepackage{bm}
\usepackage[unicode,colorlinks=true,linkcolor=black,citecolor=black,urlcolor=blue!45!black]{hyperref}
\usepackage{graphicx}

% Working-manuscript typography: restrained and journal-neutral.
\setlength{\parindent}{1.25em}
\setlength{\parskip}{0pt}
\linespread{1.055}
\setlength{\emergencystretch}{2em}
\setcounter{tocdepth}{2}
\setcounter{secnumdepth}{3}
\setlist{topsep=4pt,itemsep=2pt,parsep=0pt,partopsep=0pt}
\captionsetup{font=small,labelfont=bf,skip=7pt}
\numberwithin{equation}{section}
\pagestyle{fancy}
\fancyhf{}
\fancyhead[L]{\small\nouppercase{\leftmark}}
\fancyhead[R]{\small Quaternionic Spin Fluid}
\fancyfoot[C]{\thepage}
\renewcommand{\headrulewidth}{0.3pt}
\fancypagestyle{plain}{\fancyhf{}\fancyfoot[C]{\thepage}\renewcommand{\headrulewidth}{0pt}}
\makeatletter
\renewcommand\section{\@startsection{section}{1}{\z@}{-3.0ex plus -1ex minus -.2ex}{1.6ex plus .2ex}{\normalfont\Large\bfseries}}
\renewcommand\subsection{\@startsection{subsection}{2}{\z@}{-2.5ex plus -1ex minus -.2ex}{1.0ex plus .2ex}{\normalfont\large\bfseries}}
\renewcommand\subsubsection{\@startsection{subsubsection}{3}{\z@}{-2ex plus -.8ex minus -.2ex}{.8ex plus .2ex}{\normalfont\normalsize\bfseries}}
\makeatother
\AtBeginEnvironment{figure}{\centering}
\AtBeginEnvironment{table}{\centering}
% Figures remain buildable until the numerical plots are supplied.
\newcommand{\manuscriptfigure}[2]{%
  \IfFileExists{#1}{\includegraphics[width=#2\linewidth]{#1}}{%
    \fbox{\begin{minipage}[c][1.55in][c]{.88\linewidth}\centering\small
    Figure pending: \texttt{\detokenize{#1}}\\[5pt]
    \emph{Insert generated numerical graphic here.}
    \end{minipage}}}}
\overfullrule=0pt

% ============================================================
% MACROS
% ============================================================
\providecommand{\dd}{\mathrm{d}}
\providecommand{\vx}{\mathbf{x}}
\providecommand{\vp}{\mathbf{p}}
\providecommand{\vv}{\mathbf{v}}
\providecommand{\vu}{\mathbf{u}}
\providecommand{\vnabla}{\bm{\nabla}}
\providecommand{\vomega}{\bm{\omega}}
\providecommand{\cS}{\mathcal{S}}
\providecommand{\cQ}{\mathcal{Q}}
\providecommand{\cP}{\mathcal{P}}
\providecommand{\cA}{\mathcal{A}}
\providecommand{\cC}{\mathcal{C}}
\providecommand{\cH}{\mathcal{H}}
\providecommand{\cM}{\mathcal{M}}
\providecommand{\Oh}{\mathbb{O}}
\providecommand{\Hh}{\mathbb{H}}
\providecommand{\Rr}{\mathbb{R}}
\providecommand{\Cc}{\mathbb{C}}
\providecommand{\Spin}{\mathrm{Spin}}
\providecommand{\SO}{\mathrm{SO}}
\providecommand{\SU}{\mathrm{SU}}
\providecommand{\Uu}{\mathrm{U}}
\providecommand{\Panti}{P^{\mathrm{anti}}}

\title{Helical Breathing Solitons on the Madelung Surface\\
	of a Quaternionic Spin Fluid}
\author{}
\date{\today}

\begin{document}
	\maketitle
	
	\begin{abstract}
		We develop a quaternionic kinetic theory for a spin fluid of gyrons
		--- axisymmetric rigid bodies carrying orientation, spin, and
		position. The orientation is described by a unit quaternion, and the
		moment hierarchy is closed with relaxation-time approximations for
		the translational, spin, and nematic sectors. We identify a
		constraint surface --- the Madelung surface --- on which the spin
		amplitude and spin phase gradient satisfy equations of the Madelung
		form, including a quantum-potential term whose coefficient is
		calibrated to the physical value of $\hbar$. On this surface, we
		construct a family of localized, oscillatory, helical breather
		solutions parameterized by four free parameters. The breathers have
		the far-field structure of an electric monopole, a magnetic dipole,
		and a gravitational monopole. The mass spectrum is quantized by the
		vacuum-matching boundary condition, giving $M_n = n M_1$. We
		calibrate the model parameters against the electron mass, charge,
		spin, magnetic moment, and the speed of light. The model reproduces
		the observed values of these quantities by construction. The
		gravitational constant $G$ remains a fundamental input. We discuss
		what the model explains, what it does not, and the open problems.
	\end{abstract}
	
	\tableofcontents
\clearpage
	
	% ============================================================
	\clearpage
\part{Foundations}
	% ============================================================
	
	% ------------------------------------------------------------
	\section{Introduction}
	\label{sec:intro}
	% ------------------------------------------------------------
	
	\subsection{The gyron spin fluid}
	
	Continuum and kinetic theories with internal rotational degrees of
	freedom have a long history, ranging from the Cosserat theory of
	continua \cite{cosserat1909} and the kinetic theory of nonspherical
	molecules \cite{curtiss1956} to micropolar \cite{eringen1966} and
	spin hydrodynamics \cite{florkowski2018}. Such systems provide a
	natural setting in which translational motion, intrinsic rotation,
	vorticity, and orientational order coexist and exchange energy. They
	also provide a useful laboratory for asking how field-like structures
	can emerge from comparatively simple microscopic dynamics.
	
The gyron picture also has a distinct speculative lineage. Meno
proposed a kinetic gyron medium \cite{meno1991,meno2000}; Saul
studied spin waves and field analogies in kinetic space-time
\cite{saul2003,saul2015}, and Kyriazi developed a related
mechanical-vortex framework \cite{kyriazi2013}. These works motivate
the analogy but do not establish the present closure or soliton
results. In a different algebraic direction, Weng developed a
trigintaduonion treatment of four interactions \cite{weng2025}.
That construction employs multiple quaternionic subspaces; our
approach instead builds quaternionic structure into the orientation
variable of a kinetic spin fluid.

In a companion paper \cite{gyron}, we developed a classical kinetic
	theory for a spin fluid of \emph{gyrons}: point particles carrying
	position, velocity, and intrinsic angular momentum. Collisions
	conserve mass, momentum, total angular momentum (spin plus orbital),
	and energy. The moment hierarchy was closed with a consistent
	relaxation-time approximation, and the antisymmetric stress was
	identified as the couple stress that transfers angular momentum
	between spin and orbital sectors. We constructed a dictionary
	mapping the spin-fluid fields onto electromagnetic and gravitational
	analogues: spin pressure to electric potential, vorticity to
	magnetic field, spin density to magnetization, and total energy
	density to gravitational source. We found an exact travelling
	helical vortex soliton with a bump (anticyclone) and hole (cyclone)
	pair that serve as classical electron and positron analogues, with
	equal active gravitational mass and opposite charge-like signatures.
	In the marginal regime where the gyron thermal speed equals the
	vortex-spin wave speed, the collective transverse dynamics acquire
	an emergent Lorentz invariance, an emergent mass-shell condition,
	and an emergent $CPT$ symmetry, with the speed of light identified
	as the vortex-spin wave speed.
	
	\subsection{The wall: non-abelian gauge structure and generations}
	
	The gyron spin fluid succeeds in reproducing the kinematics of
	electromagnetism and gravity. It stalls, however, at the non-abelian
	gauge structure and the three-generation structure of the Standard
	Model. The spin variable $\bm{\sigma} \in \Rr^3$ is a real three-vector,
	and its moments are symmetric tensors. Non-abelian gauge fields
	require antisymmetric structure and internal indices, which are not
	present in the real three-vector formulation. Moreover, the model
	has no mechanism for producing three generations of particles.
	
	\subsection{The proposal: quaternionic orientation}
	
	In this paper we propose an extension in which the spin variable of
	the gyron is promoted from a real three-vector to a unit quaternion
	$q \in S^3$. The unit quaternions form the three-sphere $S^3 \cong
	\SU(2)$, which is the double cover of the rotation group $\SO(3)$.
	A unit quaternion naturally carries:
	\begin{itemize}
		\item the orientation of the gyron in space,
		\item the spin state (up to the double-cover ambiguity),
		\item the phase (the overall sign, which is the $\Uu(1)$ part).
	\end{itemize}
	
	The quaternionic formulation is strictly richer than the real
	three-vector formulation. It includes the real three-vector case as
	a special limit (when the quaternion is purely imaginary) but adds
	the phase, the $\SU(2)$ structure, and the chirality. The quaternionic
	moment hierarchy naturally carries $\SU(2)$ representations, which is
	the first step toward non-abelian gauge structure.
	
	\subsection{The Madelung surface}
	
	We identify a constraint surface in the space of gyron configurations
	--- the \emph{Madelung surface} --- on which the spin amplitude
	$\cA = |\cQ|$ and the spin phase gradient $\mathbf{v} =
	(\hbar/M) \mathrm{Im}(\bar{U} \nabla U)$ satisfy equations of the
	Madelung form:
	\begin{align}
		&\partial_t \cA + \nabla \cdot (\cA \mathbf{v}) = 0, \\
		&\partial_t \mathbf{v} + (\mathbf{v} \cdot \nabla) \mathbf{v}
		= -\nabla \Phi + D_{\mathrm{eff}} \nabla \left(
		\frac{\nabla^2 \sqrt{\cA}}{\sqrt{\cA}} \right).
	\end{align}
	The quantum-potential term has coefficient $D_{\mathrm{eff}} =
	D_s + c_2 D_a \cQ_\Hh$, and the calibration condition
	$|D_{\mathrm{eff}}| = \hbar/(2M)$ ties the model to the physical
	value of Planck's constant.
	
	The Madelung surface is defined by three closure conditions:
	an amplitude condition, a momentum condition (which requires a curl
	cancellation), and a scalar-potential condition. On this surface,
	the gyron dynamics reduces to the Madelung-transformed Schr\"odinger
	equation for the spatial part of the wavefunction, with the spin
	part carried by the orientation quaternion.
	
	\subsection{The helical breathing soliton}
	
	On the Madelung surface, we construct a family of localized,
	oscillatory, helical breather solutions. The breathers are
	characterized by:
	\begin{itemize}
		\item a radial profile in the Gaussian-Laguerre basis,
		\item a helical profile in the coordinate $\xi = z - l\theta$,
		\item an oscillatory time dependence $e^{\omega_{\mathrm{osc}} t
			\hat{n}}$.
	\end{itemize}
	
	The breathers have four free parameters: the amplitude scale $A_0$,
	the radial width $R$, the oscillation frequency
	$\omega_{\mathrm{osc}}$, and the breather speed $c$. The numerical
	solution of the spectral ODEs shows that the breather is dominated
	by the lowest radial mode and is exponentially localized in the
	helical coordinate.
	
	The far-field structure of the breather is characterized by an
	electric monopole, a magnetic dipole, and a gravitational monopole.
	The mass spectrum is quantized by the vacuum-matching boundary
	condition:
	\begin{equation}
		M_n = n M_1, \qquad M_1 = \frac{2\pi \hbar}{v_0 W}.
	\end{equation}
	This is a linear mass spectrum, which matches the hadron spectrum
	rather than the exponential lepton hierarchy.
	
	\subsection{Calibration}
	
	We calibrate the model parameters against the measured electron
	mass, charge, spin, magnetic moment, and the speed of light. The
	calibration conditions are:
	\begin{align}
		&c_{\mathrm{vortex}} = \frac{\sqrt{\chi_\rho}}{\tau_s}, \\
		&M = m_e, \\
		&Q = -e, \\
		&S = \frac{\hbar}{2}, \\
		&|D_{\mathrm{eff}}| = \frac{\hbar}{2M}, \\
		&|\mathbf{m}| = \mu_B.
	\end{align}
	These six conditions fix six combinations of the gyron parameters,
	leaving a four-dimensional residual freedom (after gauge fixing).
	The gravitational constant $G$ remains a fundamental input, not
	derived from the gyron parameters.
	
	\subsection{What the paper does and does not do}
	
	The paper does:
	\begin{itemize}
		\item develop a quaternionic kinetic theory for a spin fluid,
		\item identify the Madelung surface and its closure conditions,
		\item construct a family of localized helical breather solutions,
		\item compute the far-field multipole structure,
		\item calibrate the model against the electron properties.
	\end{itemize}
	
	The paper does not:
	\begin{itemize}
		\item derive the Standard Model gauge group,
		\item explain the three generations,
		\item derive quantum mechanics,
		\item derive general relativity,
		\item fix the gravitational constant $G$.
	\end{itemize}
	
	The model is a toy model in the analogue-gravity tradition
	\cite{unruh1981,visser1998,barcelo2011}: a classical medium may
	produce effective field kinematics without reproducing the
	fundamental dynamics of the field it resembles. In conventional
	acoustic analogues, a wave equation can be recast in terms of an
	effective Lorentzian metric under specified fluid assumptions
	\cite{visser1998,vissermolina2010}. The present quaternionic
	spin-fluid associations are distinct: an effective spacetime metric
	and Einstein-equation limit have not been established here.
	
	\subsection{Organization}
	
	The paper is organized as follows. Section~\ref{sec:kinetic}
	develops the quaternionic kinetic theory. Section~\ref{sec:closure}
	presents the closure and the closed system. Section~\ref{sec:spectrum}
	analyzes the linear spectrum. Section~\ref{sec:madelung} defines the
	Madelung surface and its closure conditions.
	Section~\ref{sec:sign} discusses the sign problem and its resolution
	by oscillation. Section~\ref{sec:helical} presents the helical
	breathing soliton and its numerical solution.
	Section~\ref{sec:farfield} computes the far-field multipole structure.
	Section~\ref{sec:calibration} presents the calibration.
	Section~\ref{sec:explains} discusses what the model explains and
	what remains open.
% ------------------------------------------------------------
\section{The Quaternionic Kinetic Theory}
\label{sec:kinetic}
% ------------------------------------------------------------

\subsection{Phase Space}
\label{sec:phase-space}

We consider a classical fluid of gyrons: axisymmetric rigid bodies
carrying position, velocity, intrinsic spin, and orientation. Each
gyron is characterized by a fixed body-frame moment-of-inertia tensor
$I_{ab}$ ($a,b$ body indices), with the symmetry axis defining a unit
director $\hat{\mathbf{n}} \in S^2$ in the lab frame. In the
quaternionic formulation, the director is replaced by a unit quaternion
$q \in S^3$, which encodes the full orientation (including the phase
about the symmetry axis).

The single-gyron phase-space variables are
\begin{equation}
	(\vx, \vv, q, \boldsymbol{\omega}),
	\label{eq:phase-space}
\end{equation}
where $\vx$ is position, $\vv$ is velocity, $q \in S^3$ is the
orientation quaternion, and $\boldsymbol{\omega}$ is the body angular
velocity. The lab-frame angular momentum is
\begin{equation}
	\boldsymbol{\sigma} = \mathbb{I}(q) \boldsymbol{\omega},
	\label{eq:angular-momentum}
\end{equation}
where $\mathbb{I}(q)$ is the lab-frame inertia tensor, which depends on
the orientation $q$.

Assuming angular velocity relaxes rapidly to quasi-equilibrium with the
orientation, we integrate out $\boldsymbol{\omega}$ and work with the
reduced distribution
\begin{equation}
	f(\vx, \vv, q, t).
	\label{eq:distribution}
\end{equation}
The measure on $S^3$ is the Haar measure $\dd\Omega_3$, normalized so
that $\int_{S^3} \dd\Omega_3 = 1$. The full phase-space measure is
$\dd\Gamma = \dd^3 v \, \dd\Omega_3$.

\subsection{Quaternion Algebra}
\label{sec:quaternion-algebra}

A quaternion is written as
\begin{equation}
	q = q_0 + q_1 i + q_2 j + q_3 k,
	\label{eq:quaternion}
\end{equation}
with multiplication rules
\begin{equation}
	i^2 = j^2 = k^2 = ijk = -1,
\end{equation}
\begin{equation}
	ij = k, \quad jk = i, \quad ki = j, \quad
	ji = -k, \quad kj = -i, \quad ik = -j.
\end{equation}
The conjugate is
\begin{equation}
	\bar{q} = q_0 - q_1 i - q_2 j - q_3 k,
\end{equation}
and the norm is
\begin{equation}
	|q|^2 = q\bar{q} = q_0^2 + q_1^2 + q_2^2 + q_3^2.
\end{equation}
For unit quaternions, $|q| = 1$, so $q \in S^3$.

A vector $\mathbf{a} \in \Rr^3$ is rotated by $q$ to
\begin{equation}
	\mathbf{a}' = q \, \mathbf{a} \, \bar{q},
\end{equation}
where $\mathbf{a}$ is embedded as the pure imaginary quaternion
$a_1 i + a_2 j + a_3 k$.

\subsection{The Kinetic Equation}
\label{sec:kinetic-equation}

The kinetic equation is
\begin{equation}
	\underbrace{\frac{\partial f}{\partial t}}_{\text{time evolution}}
	+ \underbrace{v_i \frac{\partial f}{\partial x_i}}_{\text{spatial advection}}
	+ \underbrace{\dot{q}^a \frac{\partial f}{\partial q^a}}_{\text{orientation advection}}
	= \underbrace{C[f]}_{\text{collision operator}},
	\label{eq:kinetic}
\end{equation}
where $\dot{q}$ is the orientation evolution.

The orientation evolves by the quaternionic Jeffery equation for a rigid
axisymmetric body in a flow:
\begin{equation}
	\underbrace{\dot{q}}_{\text{orientation rate}}
	= \underbrace{\frac{1}{2} \Omega \, q}_{\text{rigid-body rotation}}
	+ \underbrace{\lambda \, \mathcal{S} \, q}_{\text{Jeffery alignment}},
	\label{eq:jeffery}
\end{equation}
where
\begin{equation}
	\Omega = \Omega_{ij} e_i e_j, \qquad
	\mathcal{S} = S_{ij} e_i e_j,
\end{equation}
are the vorticity and strain-rate tensors embedded as pure imaginary
quaternions, and
\begin{equation}
	\lambda = \frac{(a/b)^2 - 1}{(a/b)^2 + 1}
	\label{eq:jeffery-parameter}
\end{equation}
is the Jeffery shape parameter, with $a/b$ the aspect ratio. For a
prolate body $\lambda > 0$; for an oblate body $\lambda < 0$; for a
sphere $\lambda = 0$.

The vorticity and strain-rate tensors are
\begin{equation}
	\Omega_{ij} = \frac{1}{2}(\partial_i u_j - \partial_j u_i), \qquad
	S_{ij} = \frac{1}{2}(\partial_i u_j + \partial_j u_i),
	\label{eq:vorticity-strain}
\end{equation}
where $u_i$ is the fluid velocity.

\subsection{Collision Invariants}
\label{sec:collision-invariants}

The collision operator conserves, at the microscopic level,
\begin{align}
	&\underbrace{\int C[f] \, \dd\Gamma = 0}_{\text{mass}}, \\
	&\underbrace{\int v_i \, C[f] \, \dd\Gamma = 0}_{\text{momentum}}, \\
	&\underbrace{\int \left( \sigma_i(q) + \epsilon_{ijk} x_j v_k \right)
		C[f] \, \dd\Gamma = 0}_{\text{total angular momentum}}, \\
	&\underbrace{\int \left( \tfrac{1}{2} v^2 + \tfrac{1}{2} I_\rho
		\omega^2 \right) C[f] \, \dd\Gamma = 0}_{\text{energy}},
	\label{eq:invariants}
\end{align}
where $I_\rho$ is the specific moment of inertia and $\sigma_i(q)$ is
the lab-frame spin angular momentum.

The spin angular momentum in quaternion form is
\begin{equation}
	\sigma_i(q) = \mathrm{Im}\left( q \, I_\rho \boldsymbol{\omega} \,
	\bar{q} \right).
	\label{eq:spin-angular-momentum}
\end{equation}

\subsection{Hydrodynamic Moments}
\label{sec:moments}

We define the mass-weighted moments:
\begin{align}
	\underbrace{\rho = \int m f \, \dd\Gamma}_{\text{mass density}}, \\
	\underbrace{\rho u_i = \int m v_i f \, \dd\Gamma}_{\text{momentum density}}, \\
	\underbrace{\cQ = \int m q f \, \dd\Gamma}_{\text{quaternionic spin density}}, \\
	\underbrace{\cP = \int m p f \, \dd\Gamma}_{\text{conjugate momentum density}}, \\
	\underbrace{P_{ij} = \int m (v_i - u_i)(v_j - u_j) f \, \dd\Gamma}_{\text{translational pressure}}, \\
	\underbrace{\Pi_{ab} = \int m (q_a - \cQ_a)(\bar{q}_b - \bar{\cQ}_b)
		f \, \dd\Gamma}_{\text{quaternionic spin pressure}}, \\
	\underbrace{\mathcal{K}_{ab} = \int m (p_a - \cP_a)(p_b - \cP_b)
		f \, \dd\Gamma}_{\text{conjugate pressure}}, \\
	\underbrace{Q_{ia} = \int m (v_i - u_i)(q_a - \cQ_a) f \, \dd\Gamma}_{\text{velocity-spin flux}}, \\
	\underbrace{R_{ia} = \int m (v_i - u_i)(p_a - \cP_a) f \, \dd\Gamma}_{\text{velocity-conjugate flux}}, \\
	\underbrace{\mathcal{M}_{ab} = \int m (q_a - \cQ_a)(p_b - \cP_b)
		f \, \dd\Gamma}_{\text{spin-conjugate flux}},
	\label{eq:moments}
\end{align}
where $p_a$ is the conjugate momentum to $q_a$, and
$\Sigma_i = \cQ_i/\rho$ is the specific spin.

\subsection{The Exact Moment Hierarchy to Third Order}
\label{sec:hierarchy}

Taking moments of the kinetic equation against each weight gives the
exact hierarchy. We present the equations to third order.

\subsubsection{Zeroth Moment (Mass)}

\begin{equation}
	\underbrace{\partial_t \rho}_{\text{density rate}}
	+ \underbrace{\partial_j (\rho u_j)}_{\text{mass flux}}
	= \underbrace{0}_{\text{mass conservation}}.
	\label{eq:moment-0}
\end{equation}

\subsubsection{First Moment (Momentum)}

\begin{equation}
	\underbrace{\partial_t (\rho u_i)}_{\text{momentum rate}}
	+ \underbrace{\partial_j (\rho u_i u_j)}_{\text{advection}}
	+ \underbrace{\partial_j P_{ij}}_{\text{pressure gradient}}
	= \underbrace{0}_{\text{momentum conservation}}.
	\label{eq:moment-1}
\end{equation}

\subsubsection{First Quaternionic Moment (Spin)}

\begin{equation}
	\underbrace{\partial_t \cQ_a}_{\text{spin rate}}
	+ \underbrace{\partial_j (u_j \cQ_a)}_{\text{advection}}
	+ \underbrace{\Omega_{ab} \cQ_b}_{\text{vorticity coupling}}
	+ \underbrace{\lambda \mathcal{S}_{ab} \cQ_b}_{\text{Jeffery coupling}}
	+ \underbrace{\partial_j Q_{ja}}_{\text{spin flux divergence}}
	= \underbrace{\mathcal{C}^{(\cQ)}_a}_{\text{collision source}}.
	\label{eq:moment-spin}
\end{equation}

\subsubsection{First Conjugate Moment}

\begin{equation}
	\underbrace{\partial_t \cP_a}_{\text{conjugate rate}}
	+ \underbrace{\partial_j (u_j \cP_a)}_{\text{advection}}
	+ \underbrace{\Omega_{ab} \cP_b}_{\text{vorticity coupling}}
	+ \underbrace{\lambda \mathcal{S}_{ab} \cP_b}_{\text{Jeffery coupling}}
	+ \underbrace{\partial_j R_{ja}}_{\text{conjugate flux divergence}}
	= \underbrace{\mathcal{C}^{(\cP)}_a}_{\text{collision source}}.
	\label{eq:moment-conjugate}
\end{equation}

\subsubsection{Second Translational Moment (Stress)}

\begin{equation}
	\begin{aligned}
		&\underbrace{\partial_t P_{ij}}_{\text{stress rate}}
		+ \underbrace{\partial_k (u_k P_{ij})}_{\text{advection}}
		+ \underbrace{P_{ik} \partial_k u_j + P_{jk} \partial_k u_i}_{\text{stretching}} \\
		&\quad + \underbrace{\partial_k Q_{ijk}}_{\text{third moment}}
		= \underbrace{\mathcal{C}^{(P)}_{ij}}_{\text{collision source}}.
	\end{aligned}
	\label{eq:moment-P}
\end{equation}

\subsubsection{Second Quaternionic Moment (Spin Pressure)}

\begin{equation}
	\begin{aligned}
		&\underbrace{\partial_t \Pi_{ab}}_{\text{spin pressure rate}}
		+ \underbrace{\partial_k (u_k \Pi_{ab})}_{\text{advection}}
		+ \underbrace{\Pi_{ak} \partial_k u_b + \Pi_{bk} \partial_k u_a}_{\text{stretching}} \\
		&\quad + \underbrace{\partial_k \Theta_{kab}}_{\text{third moment}}
		= \underbrace{\mathcal{C}^{(\Pi)}_{ab}}_{\text{collision source}}.
	\end{aligned}
	\label{eq:moment-Pi}
\end{equation}

\subsubsection{Second Conjugate Moment}

\begin{equation}
	\begin{aligned}
		&\underbrace{\partial_t \mathcal{K}_{ab}}_{\text{conjugate pressure rate}}
		+ \underbrace{\partial_k (u_k \mathcal{K}_{ab})}_{\text{advection}}
		+ \underbrace{\mathcal{K}_{ak} \partial_k u_b + \mathcal{K}_{bk} \partial_k u_a}_{\text{stretching}} \\
		&\quad + \underbrace{\partial_k \Omega_{kab}}_{\text{third moment}}
		= \underbrace{\mathcal{C}^{(\mathcal{K})}_{ab}}_{\text{collision source}}.
	\end{aligned}
	\label{eq:moment-K}
\end{equation}

\subsubsection{Mixed Velocity-Spin Moment}

\begin{equation}
	\begin{aligned}
		&\underbrace{\partial_t Q_{ia}}_{\text{velocity-spin rate}}
		+ \underbrace{\partial_k (u_k Q_{ia})}_{\text{advection}}
		+ \underbrace{Q_{ka} \partial_k u_i + Q_{ik} \partial_k u_a}_{\text{stretching}} \\
		&\quad + \underbrace{\partial_k \Theta_{kia}}_{\text{third moment}}
		= \underbrace{\mathcal{C}^{(Q)}_{ia}}_{\text{collision source}}.
	\end{aligned}
	\label{eq:moment-Q}
\end{equation}

\subsubsection{Mixed Velocity-Conjugate Moment}

\begin{equation}
	\begin{aligned}
		&\underbrace{\partial_t R_{ia}}_{\text{velocity-conjugate rate}}
		+ \underbrace{\partial_k (u_k R_{ia})}_{\text{advection}}
		+ \underbrace{R_{ka} \partial_k u_i + R_{ik} \partial_k u_a}_{\text{stretching}} \\
		&\quad + \underbrace{\partial_k \Omega_{kia}}_{\text{third moment}}
		= \underbrace{\mathcal{C}^{(R)}_{ia}}_{\text{collision source}}.
	\end{aligned}
	\label{eq:moment-R}
\end{equation}

\subsubsection{Mixed Spin-Conjugate Moment}

\begin{equation}
	\begin{aligned}
		&\underbrace{\partial_t \mathcal{M}_{ab}}_{\text{spin-conjugate rate}}
		+ \underbrace{\partial_k (u_k \mathcal{M}_{ab})}_{\text{advection}}
		+ \underbrace{\mathcal{M}_{ak} \partial_k u_b + \mathcal{M}_{bk} \partial_k u_a}_{\text{stretching}} \\
		&\quad + \underbrace{\partial_k \Xi_{kab}}_{\text{third moment}}
		= \underbrace{\mathcal{C}^{(\mathcal{M})}_{ab}}_{\text{collision source}}.
	\end{aligned}
	\label{eq:moment-M}
\end{equation}

\subsubsection{Third Moments (Definition)}

The third moments appearing in the hierarchy are
\begin{align}
	\underbrace{Q_{ijk} = \int m (v_i - u_i)(v_j - u_j)(v_k - u_k) f \, \dd\Gamma}_{\text{translational third moment}}, \\
	\underbrace{\Theta_{kab} = \int m (v_k - u_k)(q_a - \cQ_a)(\bar{q}_b - \bar{\cQ}_b) f \, \dd\Gamma}_{\text{velocity-spin-spin}}, \\
	\underbrace{\Omega_{kab} = \int m (v_k - u_k)(p_a - \cP_a)(p_b - \cP_b) f \, \dd\Gamma}_{\text{velocity-conjugate-conjugate}}, \\
	\underbrace{\Xi_{kab} = \int m (q_a - \cQ_a)(\bar{q}_b - \bar{\cQ}_b)(q_c - \cQ_c) f \, \dd\Gamma}_{\text{spin-spin-spin}}.
	\label{eq:third-moments}
\end{align}

\subsection{Structure of the Hierarchy}
\label{sec:hierarchy-structure}

The exact hierarchy is summarized in Table~\ref{tab:hierarchy}.

\begin{table}[htbp]
	\centering
	\begin{tabular}{@{}llll@{}}
		\toprule
		Equation & Moments & Third moments & Collision \\
		\midrule
		Mass & $\rho$ & --- & 0 \\
		Momentum & $\rho u_i$ & --- & 0 \\
		Spin & $\cQ_a$ & $Q_{ja}$ & $\mathcal{C}^{(\cQ)}_a$ \\
		Conjugate & $\cP_a$ & $R_{ja}$ & $\mathcal{C}^{(\cP)}_a$ \\
		Stress & $P_{ij}$ & $Q_{ijk}$ & $\mathcal{C}^{(P)}_{ij}$ \\
		Spin pressure & $\Pi_{ab}$ & $\Theta_{kab}$ & $\mathcal{C}^{(\Pi)}_{ab}$ \\
		Conjugate pressure & $\mathcal{K}_{ab}$ & $\Omega_{kab}$ & $\mathcal{C}^{(\mathcal{K})}_{ab}$ \\
		Velocity-spin & $Q_{ia}$ & $\Theta_{kia}$ & $\mathcal{C}^{(Q)}_{ia}$ \\
		Velocity-conjugate & $R_{ia}$ & $\Omega_{kia}$ & $\mathcal{C}^{(R)}_{ia}$ \\
		Spin-conjugate & $\mathcal{M}_{ab}$ & $\Xi_{kab}$ & $\mathcal{C}^{(\mathcal{M})}_{ab}$ \\
		\bottomrule
	\end{tabular}
	\caption{The exact moment hierarchy to third order.}
	\label{tab:hierarchy}
\end{table}

The hierarchy is exact. Closure is required at the level of the third
moments and the collision terms.

\subsection{Angular Momentum Conservation Constraint}
\label{sec:angular-momentum}

The total angular momentum is
\begin{equation}
	\underbrace{\mathcal{L}_i}_{\text{total angular momentum}}
	= \underbrace{\epsilon_{ijk} x_j \rho u_k}_{\text{orbital}}
	+ \underbrace{\mathrm{Im}(\cQ \cP^* e_i)}_{\text{spin}}.
	\label{eq:total-angular-momentum}
\end{equation}
Its conservation law is
\begin{equation}
	\partial_t \mathcal{L}_i + \partial_j \mathcal{J}_{ij} = 0,
	\label{eq:angular-momentum-conservation}
\end{equation}
where $\mathcal{J}_{ij}$ is the angular momentum flux. This constraint
ties the spin collision term to the antisymmetric stress:
\begin{equation}
	\underbrace{\mathcal{C}^{(\cQ)}_a}_{\text{spin source}}
	= \underbrace{\mathrm{Im}(\cQ \cP^* e_k)}_{\text{spin projection}}
	\underbrace{P^{\mathrm{anti}}_{jk}}_{\text{antisymmetric stress}}.
	\label{eq:spin-source}
\end{equation}
This is the quaternionic analogue of the vector relation
$\mathcal{C}^{(\cS)}_i = \epsilon_{ijk} P^{\mathrm{anti}}_{jk}$.

% Concise-edit: redundant recap removed; detailed derivation above retained.

\section{Closure and the Closed System}
\label{sec:closure}
% ------------------------------------------------------------

The exact moment hierarchy of Section~\ref{sec:hierarchy} is infinite.
To obtain a closed system, we must express the third moments and the
collision terms in terms of the lower moments. In this section we
present the closure approximations and the resulting closed system.

\subsection{Relaxation-Time Closure}
\label{sec:relaxation-closure}

For the translational, spin, and spin-flux sectors we adopt the
relaxation-time (BGK) approximation \cite{bgk1954}, valid in the
hydrodynamic regime $\mathrm{Kn} \ll 1$:
\begin{align}
	\underbrace{\mathcal{C}^{(P)}_{ij}}_{\text{stress collision}}
	&= \underbrace{-\frac{P_{ij} - P \delta_{ij}}{\tau}}_{\text{isotropic relaxation}}, \\
	\underbrace{\mathcal{C}^{(\Pi)}_{ab}}_{\text{spin pressure collision}}
	&= \underbrace{-\frac{\Pi_{ab} - \Pi \delta_{ab}}{\tau_s}}_{\text{spin relaxation}}, \\
	\underbrace{\mathcal{C}^{(Q)}_{ia}}_{\text{spin flux collision}}
	&= \underbrace{-\frac{Q_{ia}}{\tau_Q}}_{\text{flux relaxation}}, \\
	\underbrace{\mathcal{C}^{(R)}_{ia}}_{\text{conjugate flux collision}}
	&= \underbrace{-\frac{R_{ia}}{\tau_R}}_{\text{conjugate flux relaxation}}, \\
	\underbrace{\mathcal{C}^{(\mathcal{K})}_{ab}}_{\text{conjugate pressure collision}}
	&= \underbrace{-\frac{\mathcal{K}_{ab} - \mathcal{K} \delta_{ab}}{\tau_{\mathcal{K}}}}_{\text{conjugate relaxation}}, \\
	\underbrace{\mathcal{C}^{(\mathcal{M})}_{ab}}_{\text{spin-conjugate collision}}
	&= \underbrace{-\frac{\mathcal{M}_{ab}}{\tau_{\mathcal{M}}}}_{\text{spin-conjugate relaxation}},
	\label{eq:relaxation-closure}
\end{align}
where
\begin{equation}
	P = \frac{1}{3} P_{ii}, \qquad
	\Pi = \frac{1}{4} \Pi_{aa}, \qquad
	\mathcal{K} = \frac{1}{4} \mathcal{K}_{aa}
\end{equation}
are the isotropic pressures. The relaxation times
$\tau, \tau_s, \tau_Q, \tau_R, \tau_{\mathcal{K}}, \tau_{\mathcal{M}}$
are determined by the collision integral.

\subsection{Anisotropic Nematic Closure}
\label{sec:nematic-closure}

For the nematic sector, the isotropic relaxation
$-\hat{\mathcal{A}}_{ab}/\tau_a$ is insufficient, because the collision
integral distinguishes between different angular modes. Alignment-tensor
and orientational Fokker-Planck closures provide the relevant kinetic
precedent \cite{hess1976,doi1986}.

We define the quaternionic nematic order parameter as the traceless part
of the spin pressure:
\begin{equation}
	\underbrace{\hat{\mathcal{A}}_{ab}}_{\text{quaternionic nematic}}
	= \underbrace{\Pi_{ab} - \tfrac{1}{4} \Pi_{cc} \delta_{ab}}_{\text{traceless projection}},
	\label{eq:nematic-order}
\end{equation}
which is dimensionless after normalization by the density.

The minimal anisotropic closure is
\begin{equation}
	\underbrace{\mathcal{R}_{abcd} \hat{\mathcal{A}}_{cd}}_{\text{nematic relaxation}}
	= \underbrace{\frac{1}{\tau_a} \left[ \hat{\mathcal{A}}_{ab}
		- c_2 \mathcal{P}^{(2)}_{abcd} \hat{\mathcal{A}}_{cd}
		- c_4 \mathcal{P}^{(4)}_{abcd} \hat{\mathcal{A}}_{cd}
		- \ldots \right]}_{\text{anisotropic relaxation}},
	\label{eq:nematic-closure}
\end{equation}
where $\mathcal{P}^{(2)}, \mathcal{P}^{(4)}$ are projectors onto the
$l = 2$ and $l = 4$ angular modes. The Legendre moments $c_2, c_4$ are
those of the hard-body cross section
\begin{equation}
	\underbrace{\sigma(\hat{\mathbf{n}}, \hat{\mathbf{g}})}_{\text{cross section}}
	= \underbrace{\sigma_0 \left[ 1 + c_2 P_2(\hat{\mathbf{n}} \cdot
		\hat{\mathbf{g}}) + c_4 P_4(\hat{\mathbf{n}} \cdot \hat{\mathbf{g}})
		+ \ldots \right]}_{\text{Legendre expansion}},
	\label{eq:cross-section}
\end{equation}
where $\hat{\mathbf{n}}$ is the gyron symmetry axis, $\hat{\mathbf{g}}$
is the collision direction, and $P_l$ are Legendre polynomials. The
shape parameters $c_2, c_4$ are determined by the gyron profile.

At leading order in the anisotropy, we retain only $c_2$; $c_4$ enters
at higher order in the nematic expansion.

\subsection{Constitutive Relations}
\label{sec:constitutive}

The leading-order constitutive relations for the fluxes are
\begin{align}
	\underbrace{Q_{ia}}_{\text{spin flux}}
	&= \underbrace{-D_s \partial_i \cQ_a}_{\text{isotropic diffusion}}
	\underbrace{- c_2 D_a \partial_i (\cQ_b \hat{\mathcal{A}}_{ba})}_{\text{nematic correction}}
	+ \ldots, \\
	\underbrace{P^{\mathrm{anti}}_{ij}}_{\text{antisymmetric stress}}
	&= \underbrace{\eta_r \omega_{ij}}_{\text{viscous coupling}}
	\underbrace{- \frac{1}{2\tau_s} \mathrm{Im}(\cQ \cP^* e_k)}_{\text{spin relaxation}}
	\underbrace{+ c_2 \kappa_2 \mathrm{Im}(\cQ \hat{\mathcal{A}} \cP^* e_k)
	}_{\text{nematic correction}}
	+ \ldots, \\
	\underbrace{P_{ij}}_{\text{pressure tensor}}
	&= \underbrace{P \delta_{ij}}_{\text{isotropic}}
	\underbrace{- c_2 \mu_a \hat{\mathcal{A}}_{ij}}_{\text{nematic stress}}
	+ \ldots,
	\label{eq:constitutive}
\end{align}
where $D_s, D_a, \eta_r, \kappa_2, \mu_a$ are transport coefficients
determined by the collision integral and the gyron shape. The appearance
of intrinsic rotation, antisymmetric stress, and couple-stress-type
response has classical precedents in Cosserat and micropolar continuum
mechanics \cite{cosserat1909,eringen1966} and in kinetic theory for
nonspherical molecules \cite{curtiss1956}.

All transport coefficients depend on the Legendre moments $c_2, c_4,
\ldots$ through the hard-body cross section. For the leading-order
theory we retain only $c_2$.

For the higher moments we adopt the factorized (quasi-equilibrium)
closure,
\begin{equation}
	\underbrace{Q_{ijk}}_{\text{third moment}}
	\approx \underbrace{\frac{1}{\rho} (P_{ij} \rho u_k +
		\text{perms})}_{\text{factorized}},
	\qquad
	\underbrace{\Theta_{kab}}_{\text{velocity-spin-spin}}
	\approx \underbrace{0}_{\text{negligible}},
	\qquad
	\underbrace{\Xi_{kab}}_{\text{spin-spin-spin}}
	\approx \underbrace{0}_{\text{negligible}},
	\label{eq:factorized}
\end{equation}
which is exact for near-equilibrium distributions and approximate
otherwise.

\subsection{The Closed System}
\label{sec:closed-system}

With the closures of Sections~\ref{sec:relaxation-closure}--
\ref{sec:constitutive}, the closed system is
\begin{align}
	&\underbrace{\partial_t \rho + \partial_j (\rho u_j) = 0}_{\text{mass}}, \\
	&\underbrace{\partial_t (\rho u_i) + \partial_j (\rho u_i u_j + P \delta_{ij}
		+ P_{ij}^{(a)} + P^{\mathrm{anti}}_{ij}) = 0}_{\text{momentum}}, \\
	&\underbrace{\partial_t \cQ_a + \partial_j (u_j \cQ_a) + \Omega_{ab} \cQ_b
		+ \lambda \mathcal{S}_{ab} \cQ_b + \partial_j Q_{ja}
		= \mathcal{C}^{(\cQ)}_a}_{\text{spin}}, \\
	&\underbrace{\partial_t \cP_a + \partial_j (u_j \cP_a) + \Omega_{ab} \cP_b
		+ \lambda \mathcal{S}_{ab} \cP_b + \partial_j R_{ja}
		= \mathcal{C}^{(\cP)}_a}_{\text{conjugate}}, \\
	&\underbrace{\partial_t \hat{\mathcal{A}}_{ab} + \partial_k (u_k
		\hat{\mathcal{A}}_{ab}) + \hat{\mathcal{A}}_{ak} \Omega_{kb}
		+ \hat{\mathcal{A}}_{bk} \Omega_{ka} + \lambda \mathcal{J}_{ab}
		= -\mathcal{R}_{abcd} \hat{\mathcal{A}}_{cd}}_{\text{nematic}}, \\
	&\underbrace{\partial_t P_{ij} + \partial_k (u_k P_{ij}) + P_{ik} \partial_k u_j
		+ P_{jk} \partial_k u_i = -\frac{P_{ij} - P \delta_{ij}}{\tau}}_{\text{stress}}, \\
	&\underbrace{\partial_t \Pi_{ab} + \partial_k (u_k \Pi_{ab})
		+ \Pi_{ak} \partial_k u_b + \Pi_{bk} \partial_k u_a
		= -\frac{\Pi_{ab} - \Pi \delta_{ab}}{\tau_s}}_{\text{spin pressure}}, \\
	&\underbrace{\partial_t \mathcal{K}_{ab} + \partial_k (u_k \mathcal{K}_{ab})
		+ \mathcal{K}_{ak} \partial_k u_b + \mathcal{K}_{bk} \partial_k u_a
		= -\frac{\mathcal{K}_{ab} - \mathcal{K} \delta_{ab}}{\tau_{\mathcal{K}}}}_{\text{conjugate pressure}}, \\
	&\underbrace{\partial_t Q_{ia} + \partial_k (u_k Q_{ia}) + Q_{ka} \partial_k u_i
		+ Q_{ik} \partial_k u_a = -\frac{Q_{ia}}{\tau_Q}}_{\text{spin flux}}, \\
	&\underbrace{\partial_t R_{ia} + \partial_k (u_k R_{ia}) + R_{ka} \partial_k u_i
		+ R_{ik} \partial_k u_a = -\frac{R_{ia}}{\tau_R}}_{\text{conjugate flux}}, \\
	&\underbrace{\partial_t \mathcal{M}_{ab} + \partial_k (u_k \mathcal{M}_{ab})
		+ \mathcal{M}_{ak} \partial_k u_b + \mathcal{M}_{bk} \partial_k u_a
		= -\frac{\mathcal{M}_{ab}}{\tau_{\mathcal{M}}}}_{\text{spin-conjugate flux}},
	\label{eq:closed-system}
\end{align}
with the constitutive relations \eqref{eq:constitutive}, the
quaternionic Jeffery coupling
\begin{equation}
	\underbrace{\mathcal{J}_{ab}}_{\text{Jeffery coupling}}
	= \underbrace{\hat{\mathcal{A}}_{ak} \mathcal{S}_{kb}
		+ \hat{\mathcal{A}}_{bk} \mathcal{S}_{ka}}_{\text{strain coupling}}
	\underbrace{- \tfrac{1}{2} \delta_{ab} \hat{\mathcal{A}}_{kl}
		\mathcal{S}_{kl}}_{\text{trace subtraction}},
	\label{eq:jeffery-coupling}
\end{equation}
and the anisotropic nematic relaxation \eqref{eq:nematic-closure}.

The shape parameters $c_2, c_4$ and the Jeffery parameter $\lambda$
appear explicitly. This is the system on which all subsequent analysis
is based.

\subsection{The Quaternionic Jeffery Coupling}
\label{sec:jeffery-coupling}

The Jeffery coupling $\mathcal{J}_{ab}$ enters the nematic balance
equation. It couples the nematic order parameter to the strain rate. In
the quaternionic formulation, it takes the form
\begin{equation}
	\underbrace{\mathcal{J}_{ab}}_{\text{Jeffery coupling}}
	= \underbrace{\hat{\mathcal{A}}_{ak} \mathcal{S}_{kb}
		+ \hat{\mathcal{A}}_{bk} \mathcal{S}_{ka}}_{\text{strain alignment}}
	\underbrace{- \tfrac{1}{2} \delta_{ab} \hat{\mathcal{A}}_{kl}
		\mathcal{S}_{kl}}_{\text{traceless projection}},
	\label{eq:jeffery-coupling-explicit}
\end{equation}
where $\mathcal{S}_{ab}$ is the strain-rate tensor embedded in
quaternion indices. This is the quaternionic generalization of the
vector Jeffery coupling
\begin{equation}
	\mathcal{J}_{ij} = \hat{a}_{ik} S_{kj} + \hat{a}_{jk} S_{ki}
	- \tfrac{2}{3} \delta_{ij} \hat{a}_{kl} S_{kl}.
\end{equation}

\subsection{Transport Coefficients}
\label{sec:transport-coefficients}

The transport coefficients are determined by the collision integral and
the gyron shape. In the relaxation-time approximation they are
\begin{align}
	\underbrace{D_s}_{\text{spin diffusion}}
	&= \underbrace{\frac{k_B T \tau_s}{m}}_{\text{kinetic theory}}, \\
	\underbrace{D_a}_{\text{nematic spin transport}}
	&= \underbrace{\frac{k_B T \tau_a}{m}}_{\text{kinetic theory}}, \\
	\underbrace{\eta_r}_{\text{rotational viscosity}}
	&= \underbrace{\frac{\rho \chi_\rho}{\tau_s}}_{\text{slaved regime}}, \\
	\underbrace{\kappa_2}_{\text{spin-nematic coupling}}
	&= \underbrace{\frac{\rho \chi_\rho^2}{\tau_s}}_{\text{leading order}}, \\
	\underbrace{\mu_a}_{\text{nematic stress}}
	&= \underbrace{\frac{\rho k_B T}{m}}_{\text{ideal gas}},
	\label{eq:transport-coefficients}
\end{align}
where $T$ is the temperature and $\chi_\rho$ is the spin susceptibility.

These relations are valid in the hydrodynamic regime and at leading
order in the anisotropy. Higher-order corrections involve $c_4$ and the
full collision integral.

\subsection{The Slaved Regime}
\label{sec:slaved}

In the limit of strong spin relaxation ($\tau_s \to 0$), the spin is
slaved to the vorticity:
\begin{equation}
	\underbrace{\cQ_i}_{\text{spin density}}
	\to \underbrace{\rho \chi_\rho \omega_i}_{\text{slaved value}}.
	\label{eq:slaved}
\end{equation}
In this regime, the spin equation becomes algebraic, and the spin
density is determined by the local vorticity. This simplification is
used in the analysis of the linear spectrum and the Madelung surface.

% Concise-edit: redundant recap removed; detailed derivation above retained.

\section{Linear Spectrum}
\label{sec:spectrum}
% ------------------------------------------------------------

The closed system \eqref{eq:closed-system} admits a rich spectrum of
linear modes when perturbed around a uniform background. We catalogue
these modes here, as they determine the propagation speeds, relaxation
rates, and coupling structure that underlie the subsequent analysis.

\subsection{Setup}
\label{sec:spectrum-setup}

Linearize around the uniform background
\begin{equation}
	\underbrace{\rho = \rho_0}_{\text{uniform density}},
	\qquad
	\underbrace{\mathbf{u} = 0}_{\text{rest frame}},
	\qquad
	\underbrace{\cQ = 0}_{\text{no spin}},
	\qquad
	\underbrace{\hat{\mathcal{A}}_{ab} = 0}_{\text{no nematic}},
	\label{eq:background}
\end{equation}
and Fourier transform in space and time:
\begin{equation}
	\underbrace{\delta X \sim \hat{X} e^{i(\mathbf{k} \cdot \mathbf{x}
			- \omega t)}}_{\text{plane-wave perturbation}}.
	\label{eq:fourier}
\end{equation}
Decompose each perturbation into longitudinal
($\parallel \mathbf{k}$) and transverse ($\perp \mathbf{k}$) parts.

\subsection{Translational Sector}
\label{sec:translational-sector}

\subsubsection{Sound Mode (Longitudinal)}

For longitudinal perturbations, the spin and nematic couplings cancel,
and the density and velocity perturbations satisfy the standard
acoustic dispersion relation
\begin{equation}
	\underbrace{\omega^2 = c_s^2 k^2}_{\text{sound dispersion}},
	\qquad
	\underbrace{c_s^2 = \left. \frac{\partial P}{\partial \rho}
		\right|_S}_{\text{sound speed}},
	\label{eq:sound}
\end{equation}
where the derivative is taken at constant entropy. This is the acoustic
branch, unaffected by the spin or nematic sectors at leading order.

\subsubsection{Shear Mode (Transverse)}

For transverse perturbations, the momentum equation is diffusive:
\begin{equation}
	\underbrace{\omega = -i \nu k^2}_{\text{shear dispersion}},
	\qquad
	\underbrace{\nu = \frac{\eta}{\rho_0}}_{\text{kinematic viscosity}},
	\label{eq:shear}
\end{equation}
where $\eta$ is the shear viscosity. The shear mode is dissipative, not
propagating. In the spherical-top limit ($\lambda = 0$), the
translational transverse sector has no propagating wave; any propagating
transverse mode must come from the spin or nematic sector.

\subsection{Spin Sector}
\label{sec:spin-sector}

\subsubsection{Vortex-Spin Wave (Transverse)}

Perturbing the spin density transversally, the spin and momentum
equations couple through the antisymmetric stress. In the
short-wavelength limit $k \to \infty$ the spin mode decouples from the
diffusive velocity mode and
\begin{equation}
	\underbrace{\omega^2 \simeq c_{\mathrm{vortex}}^2 k^2}_{\text{vortex-spin dispersion}},
	\qquad
	\underbrace{c_{\mathrm{vortex}} = \frac{\sqrt{\chi_\rho}}{\tau_s}}_{\text{vortex-spin speed}}.
	\label{eq:vortex-spin}
\end{equation}
This is the vortex-spin wave, with two circular polarizations $(\pm)$.
It is the propagating transverse mode of the spin sector, and its speed
is identified with the emergent speed of light.

\subsubsection{Spin Relaxation (Longitudinal)}

For longitudinal spin perturbations, the spin equation is purely
relaxational:
\begin{equation}
	\underbrace{\omega = -\frac{i}{\tau_s}}_{\text{spin relaxation}}.
	\label{eq:spin-relaxation}
\end{equation}
The longitudinal spin has no propagating branch; it decays on the spin
relaxation timescale.

\subsection{Nematic Sector}
\label{sec:nematic-sector}

The nematic perturbation decomposes into trace, traceless symmetric, and
antisymmetric parts. The trace is slaved to the density
($\hat{\mathcal{A}}_{aa} = 0$ by construction), so only the traceless
symmetric part is dynamical.

\subsubsection{Nematic Relaxation (Transverse)}

For transverse nematic perturbations, the linearized nematic equation is
\begin{equation}
	\underbrace{-i\omega \, \delta \hat{\mathcal{A}}_{ab}}_{\text{rate}}
	= \underbrace{-\frac{\delta \hat{\mathcal{A}}_{ab}}{\tau_a}}_{\text{relaxation}}
	+ \underbrace{i \mathcal{I}_0 k \, \delta u_{ab}^+}_{\text{Jeffery driving}}
	+ \ldots,
	\label{eq:nematic-linear}
\end{equation}
where $\mathcal{I}_0$ is the isotropic nematic moment and the Jeffery
coupling at linear order reduces to the strain-rate driving. The
dispersion relation is
\begin{equation}
	\underbrace{\omega = -\frac{i}{\tau_a} + \mathcal{O}(k^2)}_{\text{nematic dispersion}}.
	\label{eq:nematic-dispersion}
\end{equation}
The nematic mode is relaxational on the uniform background. The Jeffery
coupling contributes at $\mathcal{O}(k^2)$, giving a propagating
correction at longer wavelengths, but the leading behavior is
relaxation.

\subsubsection{Nematic Longitudinal}

For longitudinal nematic perturbations, the Jeffery term is proportional
to $k \, \delta u_\parallel$, which is the sound mode. The longitudinal
nematic mode is slaved to sound and does not constitute an independent
branch.

\subsection{Hybrid Mode on the Helical Background}
\label{sec:hybrid-mode}

On the uniform background, the nematic-spin coupling
$c_2 \kappa_2 \mathrm{Im}(\cQ \hat{\mathcal{A}} \cP^* e_k)$ is quadratic
in perturbations, so it does not contribute at linear order. The nematic
and spin sectors are decoupled at the linear level.

On the helical soliton background, where $\cQ_0 \neq 0$ and
$\hat{\mathcal{A}}_{ab}^{(0)} \neq 0$, the coupling becomes linear and
the transverse spin, nematic, and velocity modes hybridize. The
resulting $3 \times 3$ system has three branches:

\begin{enumerate}
	\item \textbf{The vortex-spin wave} (spin-dominated),
	with speed $\omega^2 \simeq c_{\mathrm{vortex}}^2 k^2$ at short
	wavelengths.
	\item \textbf{The nematic relaxation mode} (nematic-dominated),
	with $\omega \simeq -i/\tau_a$.
	\item \textbf{A hybrid mode} (mixed),
	with mixed spin and nematic character, whose speed depends on the
	background coupling $c_2 \kappa_2 \cQ_0 \hat{\mathcal{A}}_0$.
\end{enumerate}


The hybrid mode is a soliton-specific excitation. It does not exist in
the vacuum, and its properties depend on the soliton background. Its
dispersion relation is the eigenvalue of the $3 \times 3$ matrix derived
from the linearized spin, nematic, and momentum equations.

\subsection{Summary of the Linear Spectrum}
\label{sec:spectrum-summary}

The linear spectrum is summarized in Table~\ref{tab:spectrum}.

\begin{table}[htbp]
	\centering
	\begin{tabular}{@{}llll@{}}
		\toprule
		Branch & Sector & Dispersion & Character \\
		\midrule
		Sound & Translational long. & $\omega^2 = c_s^2 k^2$ & Propagating \\
		Shear & Translational trans. & $\omega = -i \nu k^2$ & Diffusive \\
		Vortex-spin ($\pm$) & Spin trans. & $\omega^2 = c_{\mathrm{vortex}}^2 k^2$ & Propagating \\
		Spin relaxation & Spin long. & $\omega = -i/\tau_s$ & Relaxational \\
		Nematic relaxation & Nematic trans. & $\omega = -i/\tau_a$ & Relaxational \\
		Nematic long. & Nematic long. & slaved to sound & Not independent \\
		Hybrid & Spin + nematic & mixed & Propagating (soliton only) \\
		\bottomrule
	\end{tabular}
	\caption{The linear spectrum of the quaternionic gyron spin fluid.}
	\label{tab:spectrum}
\end{table}

The census gives four propagating branches (sound, two vortex-spin
polarizations, hybrid), three relaxational (shear, spin-longitudinal,
nematic), and one slaved (nematic trace). The structure is determined
by the moment hierarchy and the closure, and is independent of the
detailed soliton profile. The hybrid mode is the only branch that
requires the soliton background; all others exist in the vacuum.

\subsection{The Vortex-Spin Wave Speed}
\label{sec:vortex-spin-speed}

The vortex-spin wave speed is
\begin{equation}
	\underbrace{c_{\mathrm{vortex}} = \frac{\sqrt{\chi_\rho}}{\tau_s}}_{\text{vortex-spin speed}}.
	\label{eq:vortex-spin-speed}
\end{equation}
This is the speed of the propagating transverse mode of the spin sector.
It is determined by the spin susceptibility $\chi_\rho$ and the spin
relaxation time $\tau_s$.

In the marginal regime, where the gyron thermal speed equals the
vortex-spin wave speed,
\begin{equation}
	\underbrace{v_{\mathrm{th}} = c_{\mathrm{vortex}}}_{\text{marginal condition}},
	\qquad
	\underbrace{v_{\mathrm{th}} = \sqrt{\frac{k_B T}{m}}}_{\text{thermal speed}},
	\label{eq:marginal}
\end{equation}
the collective transverse dynamics acquire an emergent Lorentz
invariance. This is the subject of Section~\ref{sec:explains-relativity}.

\subsection{The Sound Speed}
\label{sec:sound-speed}

The sound speed is
\begin{equation}
	\underbrace{c_s^2 = \left. \frac{\partial P}{\partial \rho}
		\right|_S}_{\text{sound speed}}
	+ \underbrace{\left. \frac{\partial \Pi^{(s)}}{\partial \rho}
		\right|_S}_{\text{spin pressure}}
	+ \underbrace{(\text{nematic contribution})}_{\text{shape-dependent}},
	\label{eq:sound-speed}
\end{equation}
where the derivatives are taken at constant entropy. The sound speed
includes contributions from the translational pressure, the spin
pressure, and the nematic order parameter.

\subsection{The Shear Viscosity}
\label{sec:shear-viscosity}

The shear viscosity is
\begin{equation}
	\underbrace{\eta = \rho_0 \nu}_{\text{shear viscosity}},
	\label{eq:shear-viscosity}
\end{equation}
where $\nu$ is the kinematic viscosity. In the relaxation-time
approximation,
\begin{equation}
	\underbrace{\nu = \frac{k_B T \tau}{m}}_{\text{kinematic viscosity}}.
	\label{eq:kinematic-viscosity}
\end{equation}
The shear mode is diffusive, with diffusion coefficient $\nu$.

\subsection{The Spin Relaxation Time}
\label{sec:spin-relaxation-time}

The spin relaxation time $\tau_s$ determines the vortex-spin wave speed
and the spin relaxation rate. It is a fundamental parameter of the
model. In the slaved regime, it is related to the rotational viscosity
by
\begin{equation}
	\underbrace{\eta_r = \frac{\rho \chi_\rho}{\tau_s}}_{\text{slaved regime}}.
	\label{eq:rotational-viscosity}
\end{equation}

\subsection{The Nematic Relaxation Time}
\label{sec:nematic-relaxation-time}

The nematic relaxation time $\tau_a$ determines the nematic relaxation
rate. It is related to the nematic spin transport coefficient by
\begin{equation}
	\underbrace{D_a = \frac{k_B T \tau_a}{m}}_{\text{nematic transport}}.
	\label{eq:nematic-transport}
\end{equation}
The nematic mode is relaxational, with rate $1/\tau_a$.

\subsection{The Marginal Regime}
\label{sec:marginal-regime}

The marginal regime is defined by the condition
\begin{equation}
	\underbrace{v_{\mathrm{th}} = c_{\mathrm{vortex}}}_{\text{marginal}},
	\label{eq:marginal-condition}
\end{equation}
where $v_{\mathrm{th}}$ is the gyron thermal speed and
$c_{\mathrm{vortex}}$ is the vortex-spin wave speed. Solving for the
temperature,
\begin{equation}
	\underbrace{k_B T = \frac{m \chi_\rho}{\tau_s^2}}_{\text{marginal temperature}}
	= \underbrace{m c_{\mathrm{vortex}}^2}_{\text{rest mass energy}}.
	\label{eq:marginal-temperature}
\end{equation}
The gyrons are at the relativistic temperature: their thermal energy
equals their rest mass energy. This is the Galilean analogue of a
relativistic gas.

% Concise-edit: redundant recap removed; detailed derivation above retained.

\part{The Madelung Surface}
% ============================================================

% ------------------------------------------------------------
\section{The Madelung Variables}
\label{sec:madelung}
% ------------------------------------------------------------

The closed system of Section~\ref{sec:closure} admits a distinguished
class of configurations on which the spin amplitude and spin phase
gradient satisfy equations of the Madelung form. In this section we
define the Madelung variables and derive the form of the equations they
satisfy.

\subsection{The Spin Amplitude}
\label{sec:spin-amplitude}

The quaternionic spin density is
\begin{equation}
	\underbrace{\cQ = \int m q f \, \dd\Gamma}_{\text{quaternionic spin density}}
	\in \underbrace{\Hh}_{\text{quaternion algebra}}.
	\label{eq:spin-density}
\end{equation}
We decompose it into amplitude and orientation:
\begin{equation}
	\underbrace{\cQ = \cA \, U}_{\text{polar decomposition}},
	\qquad
	\underbrace{\cA = |\cQ| = \sqrt{\cQ \bar{\cQ}}}_{\text{spin amplitude}},
	\qquad
	\underbrace{U = \cQ/|\cQ| \in S^3}_{\text{orientation quaternion}}.
	\label{eq:polar-decomposition}
\end{equation}
The spin amplitude $\cA$ is a scalar field. The orientation $U$ is a
unit quaternion field.

The spin amplitude plays the role of the Madelung density:
\begin{equation}
	\underbrace{\cA \leftrightarrow \rho_{\mathrm{Madelung}}}_{\text{density association}}.
	\label{eq:amplitude-association}
\end{equation}

\subsection{The Spin Phase Gradient}
\label{sec:spin-phase-gradient}

The orientation quaternion defines a Maurer-Cartan form:
\begin{equation}
	\underbrace{\omega_i = \bar{U} \partial_i U}_{\text{Maurer-Cartan form}}
	\in \underbrace{\mathfrak{su}(2)}_{\text{Lie algebra}}.
	\label{eq:maurer-cartan}
\end{equation}
The spin phase gradient is the imaginary part of the Maurer-Cartan form:
\begin{equation}
	\underbrace{v_i = \mathrm{Im}(\bar{U} \partial_i U)}_{\text{spin phase gradient}}.
	\label{eq:spin-phase-gradient}
\end{equation}
The Madelung velocity is defined as
\begin{equation}
	\underbrace{\mathbf{v}_{\mathrm{Madelung}} = \frac{\hbar}{M}
		\mathrm{Im}(\bar{U} \nabla U)}_{\text{Madelung velocity}},
	\label{eq:madelung-velocity}
\end{equation}
where $M$ is the soliton effective mass.

The Madelung velocity plays the role of the fluid velocity:
\begin{equation}
	\underbrace{\mathbf{v}_{\mathrm{Madelung}} \leftrightarrow
		\mathbf{v}_{\mathrm{fluid}}}_{\text{velocity association}}.
	\label{eq:velocity-association}
\end{equation}

\subsection{The Curl-Free Condition}
\label{sec:curl-free}

The Maurer-Cartan form has a curvature:
\begin{equation}
	\underbrace{F_{ij} = \partial_i \omega_j - \partial_j \omega_i
		+ [\omega_i, \omega_j]}_{\text{curvature}},
	\label{eq:curvature}
\end{equation}
where $[\cdot, \cdot]$ is the commutator in $\mathfrak{su}(2)$. For a
pure gauge configuration $U$, the curvature vanishes:
\begin{equation}
	\underbrace{F_{ij} = 0}_{\text{curl-free condition}}.
	\label{eq:curl-free}
\end{equation}
This is the non-abelian curl-free condition. In the abelian limit
(complex phase), it reduces to the standard condition
$\nabla \times \mathbf{v} = 0$.

The curl-free condition ensures that the Madelung velocity is a pure
gradient:
\begin{equation}
	\underbrace{\mathbf{v}_{\mathrm{Madelung}} = \nabla \phi}_{\text{gradient}},
	\label{eq:gradient}
\end{equation}
where $\phi$ is a scalar phase. In the quaternionic case, the phase is
the orientation angle.

\subsection{The Amplitude Equation}
\label{sec:amplitude-equation}

Starting from the spin balance equation
\begin{equation}
	\underbrace{\partial_t \cQ_a + \partial_j (u_j \cQ_a) +
		\Omega_{ab} \cQ_b + \lambda \mathcal{S}_{ab} \cQ_b
		+ \partial_j Q_{ja} = \mathcal{C}^{(\cQ)}_a}_{\text{spin balance}},
	\label{eq:spin-balance}
\end{equation}
and projecting onto $\bar{U}$, we obtain the amplitude equation:
\begin{equation}
	\underbrace{\partial_t \cA + \partial_j (u_j \cA)}_{\text{amplitude transport}}
	= \underbrace{\mathrm{Re}(\bar{U} \mathcal{C}^{(\cQ)})}_{\text{collision source}}
	\underbrace{- \mathrm{Re}(\bar{U} \partial_j Q_j)}_{\text{flux divergence}}.
	\label{eq:amplitude-equation}
\end{equation}
With the spin flux closure
\begin{equation}
	\underbrace{Q_{ia} = -D_s \partial_i \cQ_a
		- c_2 D_a \partial_i (\cQ_b \hat{\mathcal{A}}_{ba})
		+ \ldots}_{\text{spin flux}},
	\label{eq:spin-flux}
\end{equation}
the amplitude equation becomes
\begin{equation}
	\underbrace{\partial_t \cA + \partial_j (u_j \cA)}_{\text{transport}}
	= \underbrace{D_{\mathrm{eff}} \nabla^2 \cA}_{\text{diffusion}}
	+ \underbrace{\cA \, \mathrm{Re}(\bar{U} \nabla^2 U)}_{\text{orientation curvature}}
	+ \underbrace{\mathrm{Re}(\bar{U} \mathcal{C}^{(\cQ)})}_{\text{source}},
	\label{eq:amplitude-equation-explicit}
\end{equation}
where the effective diffusion coefficient is
\begin{equation}
	\underbrace{D_{\mathrm{eff}} = D_s + c_2 D_a
		\mathcal{Q}_{\Hh}}_{\text{effective diffusion}},
	\qquad
	\underbrace{\mathcal{Q}_{\Hh} = \mathrm{Re}(\bar{U}
		\hat{\mathcal{A}} U)}_{\text{alignment scalar}}.
	\label{eq:effective-diffusion}
\end{equation}

\subsection{The Momentum Equation}
\label{sec:momentum-equation}

Differentiating the orientation equation and projecting onto the
orientation gradient, we obtain the momentum equation for the spin
phase gradient:
\begin{equation}
	\begin{aligned}
		&\underbrace{\partial_t v_i + u_j \partial_j v_i}_{\text{material derivative}}
		+ \underbrace{v_j \partial_i u_j}_{\text{stretching}} \\
		&\quad = \underbrace{-\partial_i \Phi_{\mathrm{spin}}}_{\text{scalar potential}}
		+ \underbrace{\mathcal{Q}_i}_{\text{quantum potential}}
		+ \underbrace{\mathcal{S}_i}_{\text{collision source}}
		+ \underbrace{\mathcal{N}_i}_{\text{non-abelian curvature}},
	\end{aligned}
	\label{eq:momentum-equation}
\end{equation}
where the individual terms are defined below.

\subsubsection{Scalar Potential}

\begin{equation}
	\underbrace{\Phi_{\mathrm{spin}}}_{\text{scalar potential}}
	= \underbrace{\tfrac{1}{2} |\mathbf{v}|^2}_{\text{kinetic energy}}
	+ \underbrace{\Omega_{jk} \mathrm{Im}(\bar{U} e_j e_k U)}_{\text{vorticity coupling}}
	+ \underbrace{\lambda S_{jk} \mathrm{Im}(\bar{U} e_j e_k U)}_{\text{strain coupling}}.
	\label{eq:scalar-potential}
\end{equation}

\subsubsection{Quantum Potential}

\begin{equation}
	\underbrace{\mathcal{Q}_i}_{\text{quantum potential}}
	= \underbrace{D_{\mathrm{eff}} \partial_i \left(
		\frac{\nabla^2 \sqrt{\cA}}{\sqrt{\cA}} \right)}_{\text{Madelung term}}
	+ \underbrace{\mathcal{Q}_i^{\mathrm{corr}}}_{\text{quaternionic correction}},
	\label{eq:quantum-potential}
\end{equation}
where the quaternionic correction involves the orientation curvature:
\begin{equation}
	\underbrace{\mathcal{Q}_i^{\mathrm{corr}}}_{\text{correction}}
	= \underbrace{D_{\mathrm{eff}} \, \mathrm{Im}\left( \bar{U} \partial_i
		\left( \frac{\nabla^2 U}{U} \right) \right)}_{\text{non-abelian Laplacian}}.
	\label{eq:quantum-potential-correction}
\end{equation}

\subsubsection{Collision Source}

\begin{equation}
	\underbrace{\mathcal{S}_i}_{\text{collision source}}
	= \underbrace{\mathrm{Im}\left( \bar{U} \partial_i \left(
		\frac{1}{\cA} \mathrm{Re}(\bar{U} \mathcal{C}^{(\cQ)}) \right)
		\right)}_{\text{projection of collision integral}}.
	\label{eq:collision-source}
\end{equation}

\subsubsection{Non-Abelian Curvature}

\begin{equation}
	\underbrace{\mathcal{N}_i}_{\text{non-abelian curvature}}
	= \underbrace{\mathrm{Im}\left( \bar{U} [\omega_i, \omega_j]
		\partial_j U \right)}_{\text{commutator of Maurer-Cartan forms}}.
	\label{eq:non-abelian-curvature}
\end{equation}
This term vanishes when the curvature $F_{ij} = 0$, which is the
curl-free condition on the Madelung surface.

\subsection{The Madelung Form}
\label{sec:madelung-form}

On the Madelung surface, the collision source and the non-abelian
curvature vanish, and the equations reduce to the Madelung form:
\begin{align}
	&\underbrace{\partial_t \cA + \nabla \cdot (\cA \mathbf{v}) = 0}_{\text{continuity}}, \\
	&\underbrace{\partial_t \mathbf{v} + (\mathbf{v} \cdot \nabla) \mathbf{v}
		= -\nabla \Phi_{\mathrm{spin}} + D_{\mathrm{eff}} \nabla \left(
		\frac{\nabla^2 \sqrt{\cA}}{\sqrt{\cA}} \right)}_{\text{momentum}}, \\
	&\underbrace{\nabla \times \mathbf{v} = 0}_{\text{curl-free}}, \\
	&\underbrace{\Phi_{\mathrm{spin}} = \Phi_{\mathrm{spin}}(\cA, \mathbf{v})}_{\text{closure}}.
	\label{eq:madelung-form}
\end{align}
These are the Madelung-transformed Schr\"odinger equations for the
spatial part of the wavefunction, with the spin part carried by the
orientation quaternion.

\subsection{The Calibration of \texorpdfstring{$\hbar$}{hbar}}
\label{sec:calibration-hbar}

The quantum potential coefficient $D_{\mathrm{eff}}$ is calibrated to
the physical value of Planck's constant:
\begin{equation}
	\underbrace{|D_{\mathrm{eff}}| = \frac{\hbar}{2M}}_{\text{calibration}},
	\label{eq:calibration-hbar}
\end{equation}
where $M$ is the soliton effective mass. This is the calibration
condition that ties the toy model to the physical value of $\hbar$.

With the measured values $\hbar = 1.055 \times 10^{-34}$ J$\cdot$s and
$M = m_e = 9.109 \times 10^{-31}$ kg:
\begin{equation}
	\underbrace{|D_{\mathrm{eff}}| = 5.79 \times 10^{-5} \,
		\mathrm{m}^2/\mathrm{s}}_{\text{calibrated value}}.
	\label{eq:calibrated-diffusion}
\end{equation}

The sign condition for the Madelung form is
\begin{equation}
	\underbrace{D_{\mathrm{eff}} < 0}_{\text{Madelung sign}},
	\label{eq:madelung-sign}
\end{equation}
which forces a strongly prolate gyron. This is the subject of
Section~\ref{sec:sign}.

\subsection{The Madelung Surface}
\label{sec:madelung-surface}

The Madelung surface is the set of gyron configurations satisfying:

\begin{enumerate}
	\item $\cA = |\cQ|$ \hfill (amplitude definition),
	\item $\mathbf{v} = (\hbar/M) \mathrm{Im}(\bar{U} \nabla U)$
	\hfill (velocity definition),
	\item $F_{ij} = 0$ \hfill (curl-free),
	\item $\partial_t \cA + \nabla \cdot (\cA \mathbf{v}) = 0$
	\hfill (continuity),
	\item $\partial_t \mathbf{v} + (\mathbf{v} \cdot \nabla)
	\mathbf{v} = -\nabla \Phi_{\mathrm{spin}} + D_{\mathrm{eff}}
	\nabla (\nabla^2 \sqrt{\cA}/\sqrt{\cA})$ \hfill (momentum),
	\item $\Phi_{\mathrm{spin}} = \Phi_{\mathrm{spin}}(\cA,
	\mathbf{v})$ \hfill (closure).
\end{enumerate}

This is the constraint surface on which the gyron dynamics reduces to
the Madelung form.

% Concise-edit: redundant recap removed; detailed derivation above retained.

\section{Closure Conditions}
\label{sec:closure-conditions}
% ------------------------------------------------------------

The Madelung form \eqref{eq:madelung-form} is not a general property of
the gyron dynamics. It holds only on a constraint surface, defined by
three closure conditions. In this section we state these conditions
exactly and derive the constraints they impose on the gyron fields.

\subsection{Closure Condition 1: Amplitude}
\label{sec:closure-1}

The amplitude equation is
\begin{equation}
	\underbrace{\partial_t \cA + \partial_j (u_j \cA)}_{\text{amplitude transport}}
	= \underbrace{-D_s \nabla^2 \cA}_{\text{isotropic diffusion}}
	\underbrace{- D_s \cA \, \mathrm{Re}(\bar{U} \nabla^2 U)}_{\text{orientation curvature}}
	\underbrace{- c_2 D_a \, \mathrm{Re}(\bar{U} \nabla^2 (\cA U
		\hat{\mathcal{A}}))}_{\text{nematic correction}}
	+ \underbrace{\mathrm{Re}(\bar{U} \mathcal{C}^{(\cQ)})}_{\text{collision source}}.
	\label{eq:closure-1-exact}
\end{equation}

For the Madelung continuity equation
\begin{equation}
	\underbrace{\partial_t \cA + \nabla \cdot (\cA \mathbf{v}) = 0}_{\text{Madelung continuity}}
	\label{eq:madelung-continuity}
\end{equation}
to hold, we require
\begin{equation}
	\boxed{
		\underbrace{-D_s \nabla^2 \cA - D_s \cA \, \mathrm{Re}(\bar{U}
			\nabla^2 U) - c_2 D_a \, \mathrm{Re}(\bar{U} \nabla^2 (\cA U
			\hat{\mathcal{A}})) + \mathrm{Re}(\bar{U} \mathcal{C}^{(\cQ)})}
		_{\text{right-hand side of \eqref{eq:closure-1-exact}}}
		= \underbrace{-\partial_j ((u_j - v_j) \cA)}_{\text{advection mismatch}}.
	}
	\label{eq:closure-1}
\end{equation}

This is the first closure condition. It relates the orientation
dynamics, the spin flux, the collision integral, and the difference
between the fluid velocity and the Madelung velocity.

In the case where the fluid velocity equals the Madelung velocity
($u_j = v_j$) and the collision integral is orthogonal to the
orientation ($\mathrm{Re}(\bar{U} \mathcal{C}^{(\cQ)}) = 0$), the
condition simplifies to
\begin{equation}
	\underbrace{-D_s \nabla^2 \cA - D_s \cA \, \mathrm{Re}(\bar{U}
		\nabla^2 U) - c_2 D_a \, \mathrm{Re}(\bar{U} \nabla^2 (\cA U
		\hat{\mathcal{A}})) = 0}_{\text{simplified closure 1}}.
	\label{eq:closure-1-simplified}
\end{equation}

\subsection{Closure Condition 2: Momentum}
\label{sec:closure-2}

The momentum equation for the spin phase gradient is
\begin{equation}
	\begin{aligned}
		&\underbrace{\partial_t v_i + u_j \partial_j v_i}_{\text{material derivative}}
		+ \underbrace{v_j \partial_i u_j}_{\text{stretching}}
		+ \underbrace{\mathrm{Im}(\bar{U} [\omega_i, \omega_j] \partial_j U)}
		_{\text{non-abelian curvature}} \\
		&\quad = \underbrace{-\partial_i \Phi_{\mathrm{spin}}}_{\text{scalar potential}}
		+ \underbrace{D_{\mathrm{eff}} \partial_i \left(
			\frac{\nabla^2 \sqrt{\cA}}{\sqrt{\cA}} \right)}_{\text{quantum potential}}
		+ \underbrace{\mathrm{Im}\left( \bar{U} \partial_i \left(
			\frac{1}{\cA} \mathrm{Re}(\bar{U} \mathcal{C}^{(\cQ)}) \right)
			\right)}_{\text{collision source}}.
	\end{aligned}
	\label{eq:closure-2-exact}
\end{equation}

For the Madelung momentum equation
\begin{equation}
	\underbrace{\partial_t \mathbf{v} + (\mathbf{v} \cdot \nabla) \mathbf{v}
		= -\nabla \Phi_{\mathrm{spin}} + D_{\mathrm{eff}} \nabla \left(
		\frac{\nabla^2 \sqrt{\cA}}{\sqrt{\cA}} \right)}_{\text{Madelung momentum}}
	\label{eq:madelung-momentum}
\end{equation}
to hold, we require
\begin{equation}
	\boxed{
		\underbrace{\mathrm{curl}\left( \mathcal{L}^{(A)} + \mathcal{L}^{(B)}
			+ \mathcal{L}^{(D)} + \mathcal{L}^{(F)} \right) = 0}_{\text{curl cancellation}},
	}
	\label{eq:closure-2}
\end{equation}
where the terms are
\begin{align}
	\underbrace{\mathcal{L}_i^{(A)}}_{\text{orientation time derivative}}
	&= \mathrm{Im}(\bar{U} \partial_i (\cA \partial_t U)), \\
	\underbrace{\mathcal{L}_i^{(B)}}_{\text{orientation advection}}
	&= \mathrm{Im}(\bar{U} \partial_i (\cA u_j \partial_j U)), \\
	\underbrace{\mathcal{L}_i^{(D)}}_{\text{collision source}}
	&= \mathrm{Im}\left( \bar{U} \partial_i \left( \frac{1}{\cA}
	\mathrm{Re}(\bar{U} \mathcal{C}^{(\cQ)}) \right) \right), \\
	\underbrace{\mathcal{L}_i^{(F)}}_{\text{stretching}}
	&= v_{0,j} \partial_i u_j.
	\label{eq:closure-2-terms}
\end{align}

The terms $\mathcal{L}^{(C)}$ (spin flux divergence) and
$\mathcal{L}^{(E)}$ (non-abelian curvature) are already gradients or
vanish on the Madelung surface, so they do not contribute to the curl
condition.

This is the second closure condition. It is the most restrictive.

The curl condition can be satisfied in two ways:
\begin{enumerate}
	\item \textbf{Individual cancellation:} each term is curl-free
	separately.
	\item \textbf{Pairwise cancellation:} the non-gradient parts of
	different terms cancel each other.
\end{enumerate}

The pairwise cancellation is the more general mechanism and is the one
relevant for the localized breather.

\subsection{Closure Condition 3: Scalar Potential}
\label{sec:closure-3}

The scalar potential is
\begin{equation}
	\underbrace{\Phi_{\mathrm{spin}}}_{\text{scalar potential}}
	= \underbrace{\tfrac{1}{2} |\mathbf{v}|^2}_{\text{kinetic energy}}
	+ \underbrace{\Omega_{jk} \mathrm{Im}(\bar{U} e_j e_k U)}_{\text{vorticity coupling}}
	+ \underbrace{\lambda S_{jk} \mathrm{Im}(\bar{U} e_j e_k U)}_{\text{strain coupling}}.
	\label{eq:closure-3-exact}
\end{equation}

For the closure condition
\begin{equation}
	\underbrace{\Phi_{\mathrm{spin}} = \Phi_{\mathrm{spin}}(\cA, \mathbf{v})
	}_{\text{closure}}.
	\label{eq:closure-3}
\end{equation}
to hold, we require
\begin{equation}
	\boxed{
		\underbrace{\Omega_{jk} \mathrm{Im}(\bar{U} e_j e_k U)
			+ \lambda S_{jk} \mathrm{Im}(\bar{U} e_j e_k U)}_{\text{orientation-fluid coupling}}
		= \underbrace{\Phi_{\mathrm{couple}}(\cA, \mathbf{v})}_{\text{function of } \cA, \mathbf{v}}.
	}
	\label{eq:closure-3-explicit}
\end{equation}

This is the third closure condition. It says that the coupling of the
orientation to the fluid vorticity and strain must be expressible as a
function of $\cA$ and $\mathbf{v}$ alone. No independent dependence on
$U$, $\hat{\mathcal{A}}$, $\Omega_{jk}$, or $S_{jk}$ is allowed.

The orientation coupling is related to the Jeffery coupling:
\begin{equation}
	\underbrace{\Omega_{jk} \mathrm{Im}(\bar{U} e_j e_k U)
		+ \lambda S_{jk} \mathrm{Im}(\bar{U} e_j e_k U)}_{\text{orientation coupling}}
	= \underbrace{2 \, \mathrm{Im}(\bar{U} \dot{U})}_{\text{time component of Maurer-Cartan}}.
	\label{eq:closure-3-jeffery}
\end{equation}

This is the time component of the Maurer-Cartan form. It is the rate of
change of the spin phase.

For the closure to hold, the rate of change of the spin phase must be a
function of $\cA$ and $\mathbf{v}$ only.

\subsection{The Calibration of \texorpdfstring{$\hbar$}{hbar}}
\label{sec:calibration}

On the Madelung surface, the quantum potential coefficient is
calibrated to the physical value of Planck's constant:
\begin{equation}
	\boxed{
		\underbrace{D_{\mathrm{eff}} = D_s + c_2 D_a \mathcal{Q}_{\Hh}}
		_{\text{effective diffusion}}
		= \underbrace{-\frac{\hbar}{2M}}_{\text{calibration}}.
	}
	\label{eq:calibration}
\end{equation}

With the measured values $\hbar = 1.055 \times 10^{-34}$ J$\cdot$s and
$M = m_e = 9.109 \times 10^{-31}$ kg:
\begin{equation}
	\underbrace{|D_{\mathrm{eff}}| = 5.79 \times 10^{-5} \,
		\mathrm{m}^2/\mathrm{s}}_{\text{calibrated value}}.
	\label{eq:calibrated-value}
\end{equation}

This is a calibration, not a derivation. The model uses the measured
value of $\hbar$ to fix one combination of the kinetic parameters.

\subsection{The Sign Condition}
\label{sec:sign-condition}

The Madelung form requires
\begin{equation}
	\underbrace{D_{\mathrm{eff}} < 0}_{\text{Madelung sign}}.
	\label{eq:sign-condition}
\end{equation}

With $D_{\mathrm{eff}} = D_s + c_2 D_a \mathcal{Q}_{\Hh}$ and
$D_s > 0$ (by the second law), this requires
\begin{equation}
	\underbrace{c_2 D_a \mathcal{Q}_{\Hh} < -D_s}_{\text{sign condition}}.
	\label{eq:sign-condition-explicit}
\end{equation}

Since $\mathcal{Q}_{\Hh} \in [-2/3, 2/3]$ and $D_s \sim D_a$, this
forces
\begin{equation}
	\boxed{
		\underbrace{|c_2| > \frac{3}{2}}_{\text{shape constraint}},
	}
	\label{eq:shape-constraint}
\end{equation}
which requires a strongly prolate gyron. Intermediate
(near-spherical) shapes are excluded, and concave profiles are not
required.

For a prolate gyron in a shear flow, the nematic principal axis aligns
with the local stretching direction while the spin is perpendicular to
it, giving $\mathcal{Q}_{\Hh} < 0$ and satisfying the sign condition
for a range of helical pitches.

\subsection{Summary of the Closure Conditions}
\label{sec:closure-summary-later-2}

The three closure conditions are summarized in
Table~\ref{tab:closure-conditions}.

\begin{table}[htbp]
	\centering
	\begin{tabular}{@{}lll@{}}
		\toprule
		Condition & Equation & Constraint \\
		\midrule
		C1 (Amplitude) & Eq.~\eqref{eq:closure-1} & Advection alignment \\
		C2 (Momentum) & Eq.~\eqref{eq:closure-2} & Curl cancellation \\
		C3 (Scalar potential) & Eq.~\eqref{eq:closure-3-explicit} & Jeffery coupling \\
		Calibration & Eq.~\eqref{eq:calibration} & $\hbar$ \\
		Sign condition & Eq.~\eqref{eq:sign-condition} & Prolate gyron \\
		\bottomrule
	\end{tabular}
	\caption{The closure conditions of the Madelung surface.}
	\label{tab:closure-conditions}
\end{table}

The closure conditions are exact, checkable, and falsifiable. They
define the Madelung surface on which the gyron dynamics reduces to the
Madelung form.

\subsection{Physical Interpretation}
\label{sec:closure-interpretation}

The three closure conditions have clear physical interpretations:

\begin{enumerate}
	\item \textbf{Closure Condition 1 (Amplitude):} The spin amplitude
	is advected by the fluid velocity and diffused by the spin flux.
	For the Madelung continuity equation to hold, the diffusion and
	source terms must cancel.
	
	\item \textbf{Closure Condition 2 (Momentum):} The spin phase
	gradient evolves by the Madelung momentum equation. For this to
	hold, the non-gradient terms in the orientation equation must
	cancel.
	
	\item \textbf{Closure Condition 3 (Scalar potential):} The scalar
	potential must close at the Madelung level. This requires the
	orientation-fluid coupling to be expressible as a function of
	$\cA$ and $\mathbf{v}$ alone.
\end{enumerate}

The closure conditions are not automatically satisfied. They define a
constraint surface in the space of gyron configurations.

\subsection{The Madelung Surface as a Constraint Surface}
\label{sec:madelung-surface-constraint}

The Madelung surface is defined by the closure conditions. It is a
constraint surface in the space of gyron configurations.

Configurations on the surface satisfy the Madelung equations.
Configurations off the surface do not.

The surface is not invariant under the gyron dynamics in general. A
configuration on the surface may evolve off it. The surface is
invariant only if the closure conditions are preserved by the
evolution, which requires additional consistency conditions.

For the localized breather, the surface is invariant by construction:
the breather is a solution of the Madelung equations, so it remains on
the surface.

% Concise-edit: redundant recap removed; detailed derivation above retained.

\section{The Sign Problem and Its Resolution}
\label{sec:sign}
% ------------------------------------------------------------

The Madelung form \eqref{eq:madelung-form} requires a specific sign of
the effective diffusion coefficient. In this section we show that the
naive sign is the wrong one for localization, and that the resolution
comes from the oscillatory character of the breathing soliton.

\subsection{The Sign Problem}
\label{sec:sign-problem}

The amplitude equation on the Madelung surface is
\begin{equation}
	\underbrace{\partial_t \cA + \partial_j (v_j \cA)}_{\text{Madelung transport}}
	= \underbrace{-D_{\mathrm{eff}} \nabla^2 \cA}_{\text{diffusion}}
	+ \underbrace{D_{\mathrm{eff}} \cA (\theta')^2}_{\text{phase gradient}}
	+ \ldots
	\label{eq:amplitude-sign}
\end{equation}

The two terms on the right-hand side have opposite signs. The diffusion
term
\begin{equation}
	\underbrace{-D_{\mathrm{eff}} \nabla^2 \cA}_{\text{diffusion}}
	\label{eq:diffusion-term}
\end{equation}
is restoring if $D_{\mathrm{eff}} > 0$ and anti-restoring if
$D_{\mathrm{eff}} < 0$.

The phase gradient term
\begin{equation}
	\underbrace{D_{\mathrm{eff}} \cA (\theta')^2}_{\text{phase gradient}}
	\label{eq:phase-gradient-term}
\end{equation}
is anti-restoring if $D_{\mathrm{eff}} > 0$ and restoring if
$D_{\mathrm{eff}} < 0$.

For a localized solution, we need the total right-hand side to be
restoring. This requires the phase gradient term to be subdominant, or
to have the opposite sign.

\subsection{The Naive Sign}
\label{sec:naive-sign}

Consider the amplitude equation for a steady, one-dimensional breather
with $\partial_t = 0$ and $\cA = \cA(\xi)$:
\begin{equation}
	\underbrace{(v_\xi - c) \cA'}_{\text{advection}}
	= \underbrace{-D_{\mathrm{eff}} \cA''}_{\text{diffusion}}
	+ \underbrace{D_{\mathrm{eff}} \cA (\theta')^2}_{\text{phase gradient}}.
	\label{eq:amplitude-1d}
\end{equation}

For a localized solution, we need $\cA \to 0$ as $\xi \to \pm\infty$.
The dominant balance at the core is between the diffusion and the
phase gradient:
\begin{equation}
	\underbrace{-D_{\mathrm{eff}} \cA''}_{\text{diffusion}}
	\approx \underbrace{-D_{\mathrm{eff}} \cA (\theta')^2}_{\text{phase gradient}}.
	\label{eq:balance}
\end{equation}
This gives
\begin{equation}
	\underbrace{\cA'' = \cA (\theta')^2}_{\text{amplitude equation}}.
	\label{eq:amplitude-ode}
\end{equation}

For $\theta' > 0$ (real phase gradient), the solution is
\begin{equation}
	\underbrace{\cA \sim e^{\pm \int \theta' d\xi}}_{\text{exponential}},
	\label{eq:exponential-solution}
\end{equation}
which is **not localized**. The amplitude grows or decays exponentially
without bound.

**This is the sign problem.** The amplitude equation has the wrong sign
for a localized solution.

\subsection{The Sign Condition}
\label{sec:sign-condition-later-2}

The resolution is that the effective diffusion coefficient must be
negative:
\begin{equation}
	\boxed{
		\underbrace{D_{\mathrm{eff}} < 0}_{\text{Madelung sign}}.
	}
	\label{eq:sign-condition-later-2}
\end{equation}

With $D_{\mathrm{eff}} < 0$, the amplitude equation becomes
\begin{equation}
	\underbrace{-|D_{\mathrm{eff}}| \cA''}_{\text{diffusion}}
	= \underbrace{-|D_{\mathrm{eff}}| \cA (\theta')^2}_{\text{phase gradient}}.
	\label{eq:amplitude-ode-negative}
\end{equation}
Simplifying:
\begin{equation}
	\underbrace{\cA'' = \cA (\theta')^2}_{\text{same equation}}.
	\label{eq:amplitude-ode-same}
\end{equation}

**The sign of $D_{\mathrm{eff}}$ does not change the sign of the
amplitude equation.** The equation is homogeneous in $D_{\mathrm{eff}}$,
so dividing by $D_{\mathrm{eff}}$ removes the sign.

This is the heart of the sign problem: the sign of $D_{\mathrm{eff}}$
does not, by itself, resolve the localization problem.

\subsection{The Shape Constraint}
\label{sec:shape-constraint}

The sign condition $D_{\mathrm{eff}} < 0$ requires
\begin{equation}
	\underbrace{D_s + c_2 D_a \mathcal{Q}_{\Hh} < 0}_{\text{sign condition}},
	\label{eq:sign-condition-explicit-later-2}
\end{equation}
where $D_s > 0$ (by the second law of thermodynamics). This requires
\begin{equation}
	\underbrace{c_2 D_a \mathcal{Q}_{\Hh} < -D_s}_{\text{negative nematic}},
	\label{eq:negative-nematic}
\end{equation}
or equivalently
\begin{equation}
	\boxed{
		\underbrace{c_2 \mathcal{Q}_{\Hh} < -\frac{D_s}{D_a}}_{\text{shape constraint}}.
	}
	\label{eq:shape-constraint-later-2}
\end{equation}

Since $\mathcal{Q}_{\Hh} \in [-2/3, 2/3]$ and $D_s \sim D_a$, this
forces
\begin{equation}
	\boxed{
		\underbrace{|c_2| > \frac{3}{2}}_{\text{strongly prolate}},
	}
	\label{eq:prolate-constraint}
\end{equation}
which requires a strongly prolate gyron. Intermediate
(near-spherical) shapes are excluded.

This is a **geometric prediction** of the model: the gyron must be a
strongly elongated, axisymmetric rigid body.

\subsection{The Oscillatory Resolution}
\label{sec:oscillatory-resolution}

The sign problem is resolved by the **oscillatory character** of the
breathing soliton. The amplitude is not static:
\begin{equation}
	\underbrace{\cA(\xi,t) = \cA_0(\xi) + \delta \cA(\xi)
		e^{i \omega_{\mathrm{osc}} t}}_{\text{oscillatory amplitude}}.
	\label{eq:oscillatory-amplitude}
\end{equation}

The oscillatory correction satisfies a driven oscillator equation:
\begin{equation}
	\underbrace{i \omega_{\mathrm{osc}} \delta \cA + (v_\xi \delta \cA)'}
	_{\text{oscillator}}
	= \underbrace{-D_s \delta \cA'' - \frac{1}{\tau_s} \delta \cA}
	_{\text{restoring}}
	+ \underbrace{\mathcal{F}_0 \mathrm{sech}(\xi/L)}_{\text{driving}}.
	\label{eq:oscillatory-equation}
\end{equation}

The driving term comes from the oscillatory vorticity:
\begin{equation}
	\underbrace{\mathcal{F}_0 \sim \frac{\rho \chi_\rho}{\tau_s}
		\mathrm{Re}(\bar{U}_0 \delta \omega_\phi)}_{\text{driving}}.
	\label{eq:driving}
\end{equation}

The oscillatory driving term can have **either sign**, depending on the
phase of the oscillation. At the soliton core, if the driving is
negative, it can balance the positive phase gradient term and allow a
localized solution.

**This is the resolution of the sign problem.** The oscillation provides
the missing degree of freedom that allows localization.

\subsection{The Breather Condition}
\label{sec:breather-condition}

For the oscillatory correction to provide a restoring force, the
driving term must be out of phase with the amplitude. This requires
\begin{equation}
	\underbrace{\omega_{\mathrm{osc}} = \frac{D_s}{L^2} -
		\frac{1}{\tau_s}}_{\text{breather condition}}.
	\label{eq:breather-condition}
\end{equation}

If this condition holds, the oscillatory correction is purely imaginary
and the time-averaged correction provides a negative contribution to the
amplitude equation:
\begin{equation}
	\underbrace{\langle \mathcal{F}_{\mathrm{osc}} \delta \cA \rangle
		\sim -|\mathcal{F}_0|^2 \kappa \, \mathrm{sech}^2(\xi/L)}
	_{\text{time-averaged correction}},
	\label{eq:time-averaged}
\end{equation}
where $\kappa > 0$.

The negative contribution balances the positive phase gradient term,
allowing a localized solution.

\subsection{The Localization Mechanism}
\label{sec:localization-mechanism}

The localization mechanism is now clear:

\begin{enumerate}
	\item The **phase gradient** term $D_{\mathrm{eff}} \cA (\theta')^2$
	is anti-restoring.
	\item The **diffusion** term $-D_{\mathrm{eff}} \cA''$ is restoring
	but has the same sign as the phase gradient after division by
	$D_{\mathrm{eff}}$.
	\item The **oscillatory driving** term provides an additional
	restoring contribution.
	\item The balance between the phase gradient, the diffusion, and
	the oscillatory driving allows a localized solution.
\end{enumerate}

**The oscillation is the key.** Without it, the amplitude equation has
no localized solution. With it, the equation admits a localized
breather.

\subsection{The Breather Frequency}
\label{sec:breather-frequency}

The breather frequency is determined by the oscillation of the
orientation:
\begin{equation}
	\underbrace{\omega_{\mathrm{osc}}}_{\text{breather frequency}}
	= \underbrace{\frac{\partial \theta}{\partial t}}_{\text{phase rate}}.
	\label{eq:breather-frequency}
\end{equation}

For a traveling breather, the phase is a function of $\xi - ct$, so
\begin{equation}
	\underbrace{\omega_{\mathrm{osc}} = c \theta'}_{\text{breather frequency}}
	= \underbrace{c \frac{M}{\hbar} v_\xi}_{\text{Madelung condition}}.
	\label{eq:breather-frequency-explicit}
\end{equation}

With $v_\xi = c$ (Madelung condition):
\begin{equation}
	\underbrace{\omega_{\mathrm{osc}} = \frac{M c^2}{\hbar}}_{\text{breather frequency}}.
	\label{eq:breather-frequency-final}
\end{equation}

**The breather frequency is proportional to the kinetic energy.**

% Concise-edit: redundant recap removed; detailed derivation above retained.

\part{The Helical Breathing Soliton}
% ============================================================

% ------------------------------------------------------------
\section{The Helical Ansatz}
\label{sec:helical}
% ------------------------------------------------------------

On the Madelung surface, the gyron dynamics reduces to the Madelung
equations. In this part we construct a family of localized, oscillatory,
helical solutions --- the \emph{helical breathing solitons} --- and
solve the resulting equations numerically.

\subsection{Helical Coordinates}
\label{sec:helical-coordinates}

We work in cylindrical coordinates $(r, \theta, z)$. The helical
coordinate is
\begin{equation}
	\underbrace{\xi = z - l\theta}_{\text{helical coordinate}},
	\label{eq:helical-coordinate}
\end{equation}
where $l$ is the \emph{helical pitch}. The helical coordinate is
constant along helical streamlines.

The derivatives with respect to $z$ and $\theta$ are related by
\begin{equation}
	\underbrace{\partial_z = \frac{1}{l}\partial_\theta =
		\partial_\xi}_{\text{helical derivative}}.
	\label{eq:helical-derivative}
\end{equation}

The helical ansatz assumes that all fields depend on $z$ and $\theta$
only through the combination $\xi$.

\subsection{The Velocity Field}
\label{sec:velocity-field}

The velocity field takes the form
\begin{equation}
	\underbrace{\mathbf{u} = u_r(r,\xi) \hat{r} + u_\theta(r,\xi)
		\hat{\theta} + u_z(r,\xi) \hat{z}}_{\text{helical velocity}},
	\label{eq:helical-velocity}
\end{equation}
with the helical constraint
\begin{equation}
	\underbrace{u_z = l u_\theta + w(r,\xi)}_{\text{helical constraint}},
	\label{eq:helical-constraint}
\end{equation}
where $w$ is the axial velocity relative to the helical frame.

The helical constraint ensures that the velocity field is tangent to
the helical streamlines. It is the defining property of the helical
ansatz.

\subsection{The Spin Density}
\label{sec:spin-density-helical}

The spin density is assumed to be azimuthal:
\begin{equation}
	\underbrace{\boldsymbol{\mathcal{S}} = \mathcal{S}_\phi(r,\xi)
		\hat{\phi}}_{\text{azimuthal spin}},
	\label{eq:azimuthal-spin}
\end{equation}
with radial and axial components vanishing by symmetry.

The azimuthal spin is the natural configuration for a helical vortex.
The spin axis is perpendicular to the helical direction.

\subsection{The Orientation Quaternion}
\label{sec:orientation-helical}

The orientation quaternion is assumed to be a phase rotation about the
azimuthal axis:
\begin{equation}
	\underbrace{U(r,\xi,t) = U_0(r,\xi) e^{\omega_{\mathrm{osc}} t
			\hat{n}}}_{\text{helical orientation}},
	\label{eq:helical-orientation}
\end{equation}
where $U_0(r,\xi)$ is the slow (envelope) orientation,
$\omega_{\mathrm{osc}}$ is the oscillation frequency, and $\hat{n}$ is
the oscillation axis.

For the helical breather, the oscillation axis is aligned with the
azimuthal direction:
\begin{equation}
	\underbrace{\hat{n} = \hat{\phi}}_{\text{oscillation axis}}.
	\label{eq:oscillation-axis}
\end{equation}

This ensures that the oscillation is consistent with the helical
symmetry.

\subsection{The Breather Ansatz}
\label{sec:breather-ansatz}

The full breather ansatz combines the helical symmetry with the
oscillatory time dependence:
\begin{align}
	\underbrace{\cA(r,\xi,t) = \cA_0(r,\xi) + \delta \cA(r,\xi)
		e^{i \omega_{\mathrm{osc}} t}}_{\text{oscillatory amplitude}}, \\
	\underbrace{U(r,\xi,t) = U_0(r,\xi) e^{\omega_{\mathrm{osc}} t
			\hat{\phi}}}_{\text{oscillatory orientation}}, \\
	\underbrace{u_\xi(r,\xi) = u_0(r,\xi)}_{\text{steady velocity}}, \\
	\underbrace{\omega_\phi(r,\xi) = \omega_0(r,\xi)}_{\text{steady vorticity}}.
	\label{eq:breather-ansatz}
\end{align}

The amplitude and orientation oscillate, while the velocity and
vorticity are steady.

\subsection{The Radial Eigenfunction Expansion}
\label{sec:radial-expansion}

The radial dependence is expanded in Gaussian-Laguerre functions:
\begin{equation}
	\underbrace{\cA_0(r,\xi) = \sum_{n=0}^{N} A_n(\xi) \phi_n(r)}
	_{\text{radial expansion}},
	\label{eq:radial-expansion}
\end{equation}
where
\begin{equation}
	\underbrace{\phi_n(r) = e^{-r^2/(2R^2)} L_n(r^2/R^2)}
	_{\text{Gaussian-Laguerre}},
	\label{eq:gaussian-laguerre}
\end{equation}
$R$ is the radial width, and $L_n$ are Laguerre polynomials.

The Gaussian-Laguerre functions are complete in the space of localized
functions and provide a natural basis for the radial structure of the
breather.

\subsection{Properties of the Gaussian-Laguerre Functions}
\label{sec:gaussian-laguerre-properties}

The Gaussian-Laguerre functions have the following properties:

\begin{enumerate}
	\item \textbf{Localization:} $\phi_n(r) \to 0$ as $r \to \infty$.
	
	\item \textbf{Regularity:} $\partial_r \phi_n(0) = 0$.
	
	\item \textbf{Orthogonality:} with respect to the weight
	$r e^{-r^2/R^2}$.
	
	\item \textbf{Completeness:} the set $\{\phi_n\}$ is complete in
	$L^2(\mathbb{R}_+, r \, dr)$.
\end{enumerate}

The first few functions are
\begin{align}
	\underbrace{\phi_0(r) = e^{-r^2/(2R^2)}}_{\text{ground state}}, \\
	\underbrace{\phi_1(r) = e^{-r^2/(2R^2)} \left( 1 -
		\frac{r^2}{R^2} \right)}_{\text{first excited}}, \\
	\underbrace{\phi_2(r) = e^{-r^2/(2R^2)} \left( 1 -
		\frac{2r^2}{R^2} + \frac{r^4}{2R^4} \right)}_{\text{second excited}}.
	\label{eq:first-functions}
\end{align}

The ground state $\phi_0$ is a Gaussian. The excited states have nodes.

\subsection{The Matrix Elements}
\label{sec:matrix-elements}

The radial Laplacian acting on the Gaussian-Laguerre functions gives
\begin{equation}
	\underbrace{\mathcal{L}_r \phi_n = \frac{1}{r}\partial_r
		(r \partial_r \phi_n)}_{\text{radial Laplacian}},
	\label{eq:radial-laplacian}
\end{equation}
with matrix elements
\begin{equation}
	\underbrace{\mathcal{L}_{mn} = \int_0^\infty \phi_m \mathcal{L}_r
		\phi_n \, r \, dr}_{\text{matrix elements}}.
	\label{eq:matrix-elements}
\end{equation}

Using the properties of Laguerre polynomials, the matrix elements are
\begin{equation}
	\boxed{
		\underbrace{\mathcal{L}_{mn} = -\frac{n+1}{2} \delta_{m,n+1}
			- \frac{2n+1}{2} \delta_{mn} - \frac{n}{2} \delta_{m,n-1}}
		_{\text{tridiagonal}}.
	}
	\label{eq:matrix-elements-explicit}
\end{equation}

This is the key result: the radial Laplacian is \textbf{tridiagonal} in
the Gaussian-Laguerre basis. This makes the spectral equations
tractable.

\subsection{The Dimensionless Variables}
\label{sec:dimensionless-variables}

It is convenient to introduce dimensionless variables:
\begin{equation}
	\underbrace{x = \frac{\xi}{L}}_{\text{dimensionless helical coordinate}},
	\label{eq:dimensionless-x}
\end{equation}
where $L$ is the packet length. All fields are then functions of
$(r, x)$.

\subsection{The Helical Symmetry}
\label{sec:helical-symmetry}

The helical ansatz has the following symmetries:

\begin{enumerate}
	\item \textbf{Translation:} the fields depend on $z$ and $\theta$
	only through $\xi$, so they are invariant under simultaneous
	translations in $z$ and $\theta$.
	
	\item \textbf{Reflection:} the fields are even in $\xi$ (for the
	symmetric breather).
	
	\item \textbf{Azimuthal:} the fields depend on $\theta$ only
	through $\xi$, so they are not invariant under azimuthal
	rotations alone.
\end{enumerate}

The helical symmetry reduces the problem from three dimensions to two:
the independent variables are $(r, \xi)$.

\subsection{The Physical Picture}
\label{sec:physical-picture}

The helical breather is a localized, oscillatory structure that winds
around a central axis. Its properties are:

\begin{enumerate}
	\item \textbf{Localization:} the amplitude decays exponentially in
	$r$ and $\xi$.
	
	\item \textbf{Oscillation:} the amplitude and orientation oscillate
	with frequency $\omega_{\mathrm{osc}}$.
	
	\item \textbf{Helicity:} the structure winds around the axis with
	pitch $l$.
	
	\item \textbf{Azimuthal spin:} the spin density is azimuthal.
\end{enumerate}

The breather is a \textbf{bound state} of the spin-vorticity
oscillation. The spin and vorticity oscillate around each other, and
the amplitude is modulated by this oscillation.

% Concise-edit: redundant recap removed; detailed derivation above retained.

\section{The Breathing Ansatz}
\label{sec:breathing}
% ------------------------------------------------------------

The helical ansatz of Section~\ref{sec:helical} describes a steady
helical structure. To obtain a localized breather, we add an
oscillatory time dependence to the amplitude and orientation. In this
section we derive the spectral equations for the breathing soliton.

\subsection{The Oscillatory Amplitude}
\label{sec:oscillatory-amplitude}

The amplitude is decomposed into a slow envelope and a fast oscillation:
\begin{equation}
	\underbrace{\cA(r,\xi,t) = \cA_0(r,\xi) + \delta \cA(r,\xi)
		e^{i \omega_{\mathrm{osc}} t}}_{\text{oscillatory amplitude}},
	\label{eq:oscillatory-amplitude-later-2}
\end{equation}
where $\cA_0$ is the slow (envelope) amplitude, $\delta \cA$ is the
oscillatory correction, and $\omega_{\mathrm{osc}}$ is the breathing
frequency.

For the helical breather, the amplitude is separable:
\begin{equation}
	\underbrace{\cA_0(r,\xi) = \sum_{n=0}^{N} A_n(\xi) \phi_n(r)}
	_{\text{radial expansion}},
	\label{eq:amplitude-expansion}
\end{equation}
where $\phi_n(r)$ are the Gaussian-Laguerre functions of
Section~\ref{sec:radial-expansion}.

The oscillatory correction $\delta \cA$ satisfies a similar expansion:
\begin{equation}
	\underbrace{\delta \cA(r,\xi) = \sum_{n=0}^{N} \delta A_n(\xi)
		\phi_n(r)}_{\text{oscillatory expansion}}.
	\label{eq:oscillatory-expansion}
\end{equation}

\subsection{The Oscillatory Orientation}
\label{sec:oscillatory-orientation}

The orientation quaternion is
\begin{equation}
	\underbrace{U(r,\xi,t) = U_0(r,\xi) e^{\omega_{\mathrm{osc}} t
			\hat{\phi}}}_{\text{oscillatory orientation}},
	\label{eq:oscillatory-orientation}
\end{equation}
where $U_0$ is the slow orientation and $\hat{\phi}$ is the azimuthal
axis.

The slow orientation is a phase rotation about the azimuthal axis:
\begin{equation}
	\underbrace{U_0(r,\xi) = e^{\theta_0(r,\xi) \hat{\phi}}}_{\text{slow orientation}},
	\label{eq:slow-orientation}
\end{equation}
where $\theta_0(r,\xi)$ is the slow phase.

The full orientation is therefore
\begin{equation}
	\underbrace{U(r,\xi,t) = e^{\theta_0(r,\xi) \hat{\phi}}
		e^{\omega_{\mathrm{osc}} t \hat{\phi}}}_{\text{full orientation}}
	= e^{(\theta_0(r,\xi) + \omega_{\mathrm{osc}} t) \hat{\phi}}.
	\label{eq:full-orientation}
\end{equation}

The orientation is a pure phase rotation about the azimuthal axis, with
a phase that is the sum of the slow spatial part and the fast temporal
part.

\subsection{The Spin Phase Gradient}
\label{sec:spin-phase-gradient-breathing}

The spin phase gradient is
\begin{equation}
	\underbrace{v_i = \mathrm{Im}(\bar{U} \partial_i U)}_{\text{spin phase gradient}}.
	\label{eq:spin-phase-gradient-later-2}
\end{equation}
With the orientation \eqref{eq:full-orientation}:
\begin{equation}
	\underbrace{\bar{U} \partial_i U = \partial_i \theta_0 \hat{\phi}
		+ (\theta_0 + \omega_{\mathrm{osc}} t) \partial_i \hat{\phi}}
	_{\text{Maurer-Cartan form}}.
	\label{eq:maurer-cartan-breathing}
\end{equation}
Since $\hat{\phi}$ is independent of $r$ and $\xi$ (it depends only on
$\theta$), and $\partial_\theta \hat{\phi} = -\hat{r}$, we have
\begin{equation}
	\underbrace{\partial_i \hat{\phi} = -\hat{r} \delta_{i\theta}}
	_{\text{azimuthal derivative}}.
	\label{eq:azimuthal-derivative}
\end{equation}

The spin phase gradient is therefore
\begin{equation}
	\underbrace{v_i = \partial_i \theta_0 \hat{\phi} \cdot \hat{\phi}
		- (\theta_0 + \omega_{\mathrm{osc}} t) \delta_{i\theta}
		\hat{r} \cdot \hat{\phi}}_{\text{spin phase gradient}}.
	\label{eq:spin-phase-gradient-explicit}
\end{equation}
Taking the imaginary part and projecting onto the helical direction:
\begin{equation}
	\underbrace{v_\xi = \partial_\xi \theta_0}_{\text{helical phase gradient}}.
	\label{eq:helical-phase-gradient}
\end{equation}

The azimuthal part vanishes because $\hat{r} \cdot \hat{\phi} = 0$.

\subsection{The Madelung Velocity}
\label{sec:madelung-velocity-breathing}

The Madelung velocity is
\begin{equation}
	\underbrace{\mathbf{v}_{\mathrm{Madelung}} = \frac{\hbar}{M}
		\mathrm{Im}(\bar{U} \nabla U)}_{\text{Madelung velocity}}.
	\label{eq:madelung-velocity-later-2}
\end{equation}
With the spin phase gradient \eqref{eq:helical-phase-gradient}:
\begin{equation}
	\underbrace{\mathbf{v}_{\mathrm{Madelung}} = \frac{\hbar}{M}
		\partial_\xi \theta_0 \hat{\xi}}_{\text{helical Madelung velocity}}.
	\label{eq:helical-madelung-velocity}
\end{equation}

The Madelung velocity is along the helical direction.

\subsection{The Dimensionless Spectral Equations}
\label{sec:spectral-equations}

Substituting the breathing ansatz into the amplitude equation
\eqref{eq:amplitude-equation-explicit} and projecting onto the
Gaussian-Laguerre mode $\phi_m$, we obtain the spectral equations:
\begin{equation}
	\boxed{
		\underbrace{(u_\xi - c) \partial_x A_m}_{\text{advection}}
		= \underbrace{\alpha \left( m A_{m-1} + (2m+1) A_m
			+ (m+1) A_{m+1} \right)}_{\text{radial diffusion}}
		\underbrace{- \beta A_m}_{\text{relaxation}}
		+ \underbrace{\gamma \, \mathrm{sech}(x) \delta_{m0}}_{\text{driving}},
	}
	\label{eq:spectral-equations}
\end{equation}
where the dimensionless parameters are
\begin{align}
	\underbrace{\tilde{u} = \frac{u_\xi - c}{c_{\mathrm{vortex}}}}
	_{\text{dimensionless advection}}, \\
	\underbrace{\alpha = \frac{L D_s}{2 c_{\mathrm{vortex}}}}
	_{\text{dimensionless diffusion}}, \\
	\underbrace{\beta = \frac{L}{\tau_s c_{\mathrm{vortex}}}}
	_{\text{dimensionless relaxation}}, \\
	\underbrace{\gamma = \frac{L \mathcal{F}_0}{c_{\mathrm{vortex}}}}
	_{\text{dimensionless driving}}.
	\label{eq:dimensionless-parameters}
\end{align}

This is the core system of spectral ODEs. It determines the helical
profile of the breather.

\subsection{The Structure of the Spectral Equations}
\label{sec:spectral-structure}

The spectral equations have the following structure:

\begin{enumerate}
	\item \textbf{Advection:} the left-hand side is the advection of
	the mode amplitude by the helical velocity.
	
	\item \textbf{Radial diffusion:} the first term on the right-hand
	side couples each mode to its neighbors. The coupling is
	tridiagonal.
	
	\item \textbf{Relaxation:} the second term is the relaxation of
	each mode to zero.
	
	\item \textbf{Driving:} the third term is the driving of the
	$m = 0$ mode by the sech profile.
\end{enumerate}

The equations are \textbf{linear} in the mode amplitudes and
\textbf{coupled} through the tridiagonal diffusion term.

\subsection{The Boundary Conditions}
\label{sec:boundary-conditions}

For a localized breather, the mode amplitudes must vanish at the
boundaries:
\begin{equation}
	\underbrace{A_m(\pm \infty) = 0}_{\text{localization}}.
	\label{eq:boundary-conditions}
\end{equation}

The boundary conditions ensure that the breather is localized in the
helical coordinate.

\subsection{The Dimensionless Parameters}
\label{sec:dimensionless-parameters-values}

For the illustrative physical parameter choice:

\begin{align}
	\underbrace{\tilde{u} = 1.67 \times 10^{-2}}_{\text{advection}}, \\
	\underbrace{\alpha = 9.65 \times 10^{-27}}_{\text{diffusion}}, \\
	\underbrace{\beta = 3.33 \times 10^{-1}}_{\text{relaxation}}, \\
	\underbrace{\gamma = 3.33 \times 10^{-2}}_{\text{driving}}.
	\label{eq:dimensionless-values}
\end{align}

\textbf{The diffusion parameter $\alpha$ is negligibly small.} This
means the radial diffusion term is negligible, and the modes decouple.

\subsection{The Decoupled Limit}
\label{sec:decoupled-limit}

In the limit $\alpha \to 0$, the spectral equations decouple:
\begin{align}
	\underbrace{\tilde{u} \partial_x A_0 = -\beta A_0 + \gamma
		\mathrm{sech}(x)}_{\text{driven mode}}, \\
	\underbrace{\tilde{u} \partial_x A_m = -\beta A_m}_{\text{decaying modes}},
	\quad m \geq 1.
	\label{eq:decoupled}
\end{align}

The $m = 0$ mode satisfies a driven equation, while the $m \geq 1$
modes satisfy homogeneous decay equations.

\subsection{The Analytic Solution for the Driven Mode}
\label{sec:analytic-solution}

The $m = 0$ equation is
\begin{equation}
	\underbrace{\tilde{u} \partial_x A_0 = -\beta A_0 + \gamma
		\mathrm{sech}(x)}_{\text{driven equation}}.
	\label{eq:driven-equation}
\end{equation}

This is a first-order linear ODE. The solution is
\begin{equation}
	\boxed{
		\underbrace{A_0(x) = \frac{\gamma}{\tilde{u}} \int_{-\infty}^{x}
			e^{-\beta(x - x')/\tilde{u}} \mathrm{sech}(x') \, dx'}
		_{\text{convolution integral}}.
	}
	\label{eq:analytic-solution}
\end{equation}

This is a convolution of the sech driving with an exponential kernel.

\subsection{The Localization Condition}
\label{sec:localization-condition}

For the solution \eqref{eq:analytic-solution} to be localized, we need
\begin{equation}
	\underbrace{\frac{\beta}{\tilde{u}} > 0}_{\text{localization}}.
	\label{eq:localization-condition}
\end{equation}

Since $\beta > 0$, this requires
\begin{equation}
	\underbrace{\tilde{u} > 0}_{\text{positive advection}},
	\label{eq:positive-advection}
\end{equation}
or equivalently
\begin{equation}
	\underbrace{u_\xi - c > 0}_{\text{breather speed}}.
	\label{eq:breather-speed-condition}
\end{equation}

\textbf{The breather speed must be less than the helical velocity.}

\subsection{The Mode Structure}
\label{sec:mode-structure}

The mode structure is:

\begin{enumerate}
	\item \textbf{Mode 0:} the dominant mode, driven by the sech term,
	localized with a sech-like profile.
	
	\item \textbf{Modes 1, 2, ...:} exponentially small, because the
	diffusion parameter $\alpha$ is negligibly small.
\end{enumerate}

\textbf{The breather is a single-mode solution.} The higher modes are
suppressed by $\alpha \sim 10^{-26}$.

% Concise-edit: redundant recap removed; detailed derivation above retained.

\section{Numerical Solution}
\label{sec:numerical}
% ------------------------------------------------------------

The spectral equations \eqref{eq:spectral-equations} are a coupled
system of linear ODEs for the mode amplitudes $A_m(x)$. In this section
we solve them numerically and characterize the breather solution.

\subsection{The Dimensionless Parameters}
\label{sec:dimensionless-parameters}

The dimensionless parameters are defined by
\begin{align}
	\underbrace{\tilde{u} = \frac{u_\xi - c}{c_{\mathrm{vortex}}}}
	_{\text{dimensionless advection}}, \\
	\underbrace{\alpha = \frac{L D_s}{2 c_{\mathrm{vortex}}}}
	_{\text{dimensionless diffusion}}, \\
	\underbrace{\beta = \frac{L}{\tau_s c_{\mathrm{vortex}}}}
	_{\text{dimensionless relaxation}}, \\
	\underbrace{\gamma = \frac{L \mathcal{F}_0}{c_{\mathrm{vortex}}}}
	_{\text{dimensionless driving}}.
	\label{eq:dimensionless-parameters-later-2}
\end{align}

For the illustrative physical parameter choice:

\begin{align}
	\underbrace{\tilde{u} = 1.67 \times 10^{-2}}_{\text{advection}}, \\
	\underbrace{\alpha = 9.65 \times 10^{-27}}_{\text{diffusion}}, \\
	\underbrace{\beta = 3.33 \times 10^{-1}}_{\text{relaxation}}, \\
	\underbrace{\gamma = 3.33 \times 10^{-2}}_{\text{driving}}.
	\label{eq:dimensionless-values-later-2}
\end{align}

The diffusion parameter $\alpha$ is negligibly small. This means the
radial diffusion term is negligible, and the modes decouple.

\subsection{The Sparse Linear System}
\label{sec:sparse-system}

The spectral equations \eqref{eq:spectral-equations} are discretized on
a uniform grid in $x$ with central differences:
\begin{equation}
	\underbrace{\partial_x A_m \approx \frac{A_{m,j+1} -
			A_{m,j-1}}{2 \Delta x}}_{\text{central difference}},
	\label{eq:central-difference}
\end{equation}
where $j$ indexes the grid points.

The resulting linear system is
\begin{equation}
	\underbrace{\mathbf{M} \mathbf{A} = \mathbf{F}}_{\text{linear system}},
	\label{eq:linear-system}
\end{equation}
where $\mathbf{M}$ is a sparse block-tridiagonal matrix, $\mathbf{A}$
is the vector of mode amplitudes, and $\mathbf{F}$ is the driving
vector.

The matrix $\mathbf{M}$ has the following structure:
\begin{enumerate}
	\item \textbf{Block-tridiagonal:} each block corresponds to a
	grid point, and the blocks are coupled through the advection and
	diffusion terms.
	
	\item \textbf{Sparse:} the number of non-zero entries is
	$\mathcal{O}(N \cdot J)$, where $N$ is the number of modes and
	$J$ is the number of grid points.
	
	\item \textbf{Non-singular:} the boundary conditions ensure that
	$\mathbf{M}$ is invertible.
\end{enumerate}

The linear system is solved using a sparse linear solver.
% ------------------------------------------------------------
% Repeated numerical setup consolidated into the preceding section.

\subsection{Numerical Results}
\label{sec:numerical-results}

The numerical solution of the spectral equations is shown in
Figure~\ref{fig:breather}. The four panels display:

\begin{enumerate}
	\item \textbf{Mode profiles} (top left): the amplitudes
	$A_m(x)$ for $m = 0, 1, 2, 3$. The $m = 0$ mode dominates,
	confirming the single-mode structure. The higher modes are
	suppressed by the small diffusion parameter $\alpha \sim
	10^{-26}$.
	
	\item \textbf{Radial amplitude} (top right): the amplitude
	$\cA_0(r, 0)$ at $\xi = 0$ as a function of $r/R$. The profile
	is Gaussian, consistent with the ground-state
	Gaussian-Laguerre function.
	
	\item \textbf{Associated fields} (bottom left): the normalized
	phase gradient $\theta'(\xi)$ and the vorticity
	$\omega_\phi(\xi)$. Both have the sech-like profile of the
	driven mode.
	
	\item \textbf{Two-dimensional amplitude} (bottom right): the
	full amplitude $|\cA_0(r, \xi)|$ in the $(r, \xi)$ plane. The
	breather is localized in both directions.
\end{enumerate}

\begin{figure}[htbp]
	\centering
	\manuscriptfigure{helical_breather.png}{1.0}
	\caption{Numerical solution of the spectral equations for the
		helical breathing soliton. Top left: mode profiles $A_m(x)$.
		Top right: radial amplitude $\cA_0(r, 0)$. Bottom left:
		associated fields $\theta'(\xi)$ and $\omega_\phi(\xi)$.
		Bottom right: two-dimensional amplitude
		$|\cA_0(r, \xi)|$.}
	\label{fig:breather}
\end{figure}

\subsection{The Physical Amplitude}
\label{sec:physical-amplitude}

The physical amplitude is reconstructed from the mode amplitudes:
\begin{equation}
	\underbrace{\cA_0(r,\xi) = \sum_{n=0}^{N} A_n(\xi)
		\phi_n(r)}_{\text{physical amplitude}},
	\label{eq:physical-amplitude}
\end{equation}
where $\phi_n(r)$ are the Gaussian-Laguerre functions.

The peak amplitude is
\begin{equation}
	\underbrace{\cA_0(0, 0) = 9.98 \times 10^{-2}}_{\text{peak amplitude}}.
	\label{eq:peak-amplitude}
\end{equation}

The amplitude is localized in both $r$ and $\xi$, with a Gaussian
radial profile and a sech-like helical profile.

\subsection{Parameter Sweep and Mass Spectrum}
\label{sec:parameter-sweep}

We perform a parameter sweep over the driving parameter $\gamma$ and
the density scale $\rho$. The mass is computed from the corrected
volume element:
\begin{equation}
	\underbrace{M_{\mathrm{eff}} = \frac{2\pi L}{c_{\mathrm{vortex}}^2}
		\int_0^\infty \int_{-\infty}^\infty \rho |\cA_0(r,x)|^2
		\, r \, dr \, dx}_{\text{effective mass}}.
	\label{eq:effective-mass}
\end{equation}

The mass scales as
\begin{equation}
	\underbrace{M_{\mathrm{eff}} \propto \gamma^2 \rho}_{\text{mass scaling}}.
	\label{eq:mass-scaling}
\end{equation}

Figure~\ref{fig:mass-spectrum} shows the mass spectrum and parameter
sweep. The mass depends on the driving parameter $\gamma$ and the
density scale $\rho$. The mass crosses the electron mass $m_e$ at a
specific combination of parameters.

\begin{figure}[htbp]
	\centering
	\manuscriptfigure{helical_breather_complete.png}{0.85}
	\caption{Mass spectrum and parameter sweep of the helical breather.
		The figure shows the dependence of the effective mass on the
		driving parameter $\gamma$ and the density scale $\rho$.}
	\label{fig:mass-spectrum}
\end{figure}

\subsection{The Mode Structure}
\label{sec:mode-structure-later-2}

The mode structure is characterized by:

\begin{enumerate}
	\item \textbf{Single-mode dominance:} the $m = 0$ mode carries
	essentially all of the amplitude. The higher modes are
	suppressed by $\alpha \sim 10^{-26}$.
	
	\item \textbf{Localization:} the mode amplitudes decay to zero
	at the boundaries, confirming the localized nature of the
	breather.
	
	\item \textbf{sech-like profile:} the $m = 0$ mode has a
	sech-like profile, consistent with the analytic solution
	\eqref{eq:analytic-solution}.
\end{enumerate}

\subsection{The Source-Sink Structure}
\label{sec:source-sink}

The breather has a source-sink structure in the helical coordinate.
The velocity field has a source region where the flow diverges and a
sink region where it converges. This structure is responsible for the
localization of the breather in the helical direction.

The source-sink structure is visible in the bottom-left panel of
Figure~\ref{fig:breather}, where the phase gradient $\theta'(\xi)$
changes sign across the core.

% Concise-edit: redundant recap removed; detailed derivation above retained.

\section{The Breather Family}
\label{sec:family}
% ------------------------------------------------------------

The numerical solution of Section~\ref{sec:numerical} yields a
one-parameter family of localized breather solutions. In this section
we characterize the family, identify its free parameters, and compute
its physical properties.

\subsection{The Four Free Parameters}
\label{sec:free-parameters}

The breather family is parameterized by four quantities:

\begin{enumerate}
	\item \textbf{Amplitude scale $A_0$:} the overall scale of the
	amplitude. Determined by the driving strength $\gamma$.
	
	\item \textbf{Radial width $R$:} the characteristic radius of the
	Gaussian radial profile. Determined by the balance of diffusion
	and relaxation.
	
	\item \textbf{Oscillation frequency $\omega_{\mathrm{osc}}$:} the
	frequency of the breathing oscillation. Determined by the
	breather condition.
	
	\item \textbf{Breather speed $c$:} the speed at which the
	breather translates. A free parameter, subject to the
	localization condition.
\end{enumerate}

The four parameters define a four-dimensional family of breathers.

\subsection{The Amplitude Scale}
\label{sec:amplitude-scale}

The amplitude scale $A_0$ is determined by the dimensionless driving
parameter $\gamma$:
\begin{equation}
	\underbrace{A_0 \propto \gamma}_{\text{amplitude scaling}}.
	\label{eq:amplitude-scaling}
\end{equation}

From the numerical solution:
\begin{equation}
	\underbrace{A_0 = 9.98 \times 10^{-2}}_{\text{peak amplitude}}
	\quad \text{for} \quad
	\underbrace{\gamma = 3.33 \times 10^{-2}}_{\text{base driving}}.
	\label{eq:amplitude-value}
\end{equation}

The amplitude scale determines the peak amplitude of the breather, and
hence its mass.

\subsection{The Radial Width}
\label{sec:radial-width}

The radial width $R$ is determined by the balance of the radial
diffusion and the relaxation:
\begin{equation}
	\underbrace{R \sim \sqrt{2 D_s \tau_s}}_{\text{radial width}}.
	\label{eq:radial-width}
\end{equation}

With $D_s \sim 10^{-5}$ m$^2$/s and $\tau_s \sim 10^{-21}$ s:
\begin{equation}
	\underbrace{R \sim 1.4 \times 10^{-13} \, \text{m}}_{\text{Compton scale}}.
	\label{eq:radial-width-value}
\end{equation}

This is the Compton wavelength of the electron. The radial width is
fixed by the calibration to the electron mass.

\subsection{The Oscillation Frequency}
\label{sec:oscillation-frequency}

The oscillation frequency $\omega_{\mathrm{osc}}$ is determined by the
breather condition
\begin{equation}
	\underbrace{\omega_{\mathrm{osc}} = c \theta'}_{\text{breather condition}}
	= \underbrace{c \frac{M}{\hbar} v_\xi}_{\text{Madelung condition}}.
	\label{eq:oscillation-frequency}
\end{equation}

With $v_\xi = c$ (Madelung condition):
\begin{equation}
	\underbrace{\omega_{\mathrm{osc}} = \frac{M c^2}{\hbar}}_{\text{oscillation frequency}}.
	\label{eq:oscillation-frequency-final}
\end{equation}

The oscillation frequency is proportional to the kinetic energy of the
breather.

\subsection{The Breather Speed}
\label{sec:breather-speed}

The breather speed $c$ is a free parameter. It is subject to the
localization condition
\begin{equation}
	\underbrace{u_\xi - c > 0}_{\text{localization}}.
	\label{eq:localization-condition-later-2}
\end{equation}

Since $u_\xi$ is the helical velocity, the breather speed must be less
than the helical velocity for the breather to be localized.

The breather speed determines the de Broglie wavelength:
\begin{equation}
	\underbrace{\lambda_{\mathrm{dB}} = \frac{h}{M c}}_{\text{de Broglie wavelength}}.
	\label{eq:de-broglie}
\end{equation}

\subsection{The Mode Structure}
\label{sec:family-mode-structure}

The mode structure is characterized by:

\begin{enumerate}
	\item \textbf{$m = 0$ dominance:} the $m = 0$ mode carries
	essentially all of the amplitude. The higher modes are
	suppressed by the small diffusion parameter.
	
	\item \textbf{sech-like profile:} the $m = 0$ mode has a
	sech-like profile in the helical coordinate.
	
	\item \textbf{Gaussian radial profile:} the radial profile is
	Gaussian, consistent with the ground-state Gaussian-Laguerre
	function.
\end{enumerate}

The mode structure is determined by the linear spectral equations and
is independent of the detailed nonlinear dynamics.

\subsection{The Amplitude Profile}
\label{sec:amplitude-profile}

The amplitude profile is
\begin{equation}
	\underbrace{\cA_0(r,\xi) = A_0 e^{-r^2/(2R^2)}
		\mathrm{sech}(\xi/L)}_{\text{amplitude profile}},
	\label{eq:amplitude-profile}
\end{equation}
where $A_0$ is the peak amplitude, $R$ is the radial width, and $L$ is
the packet length.

The profile is the product of a Gaussian in $r$ and a sech in $\xi$.
This is the natural profile for a localized breather on the helical
background.

\subsection{The Associated Fields}
\label{sec:associated-fields}

The other fields are determined by the amplitude and the closure
conditions:

\subsubsection{Density Deviation}

\begin{equation}
	\underbrace{\delta \rho(r,\xi) = \rho_{\mathrm{scale}}
		\cA_0(r,\xi)}_{\text{density deviation}}.
	\label{eq:density-deviation}
\end{equation}

\subsubsection{Spin Density}

\begin{equation}
	\underbrace{\mathcal{S}(r,\xi,t) = \cA_{\mathrm{scale}}
		\cA_0(r,\xi) \cos(\omega_{\mathrm{osc}} t) \hat{\phi}}
	_{\text{spin density}}.
	\label{eq:spin-density-family}
\end{equation}

\subsubsection{Vorticity}

\begin{equation}
	\underbrace{\omega_\phi(r,\xi) = \frac{\cA_{\mathrm{scale}}}
		{\rho \chi_\rho} \cA_0(r,\xi)}_{\text{vorticity}}.
	\label{eq:vorticity-family}
\end{equation}

\subsubsection{Velocity}

\begin{equation}
	\underbrace{u_\xi(r,\xi) = \frac{\hbar}{M} \theta'(r,\xi)}
	_{\text{velocity}},
	\label{eq:velocity-family}
\end{equation}
with $\theta'$ determined by the Madelung condition.

\subsubsection{Pressure}

\begin{equation}
	\underbrace{P(r,\xi) = c_s^2 \rho_{\mathrm{scale}}
		\cA_0(r,\xi)}_{\text{pressure}}.
	\label{eq:pressure-family}
\end{equation}

\subsubsection{Spin Pressure}

\begin{equation}
	\underbrace{\Pi^{(s)}(r,\xi) = \frac{\rho_{\mathrm{scale}}^3
			\cA_0^3}{\rho^2 \chi_\rho}}_{\text{spin pressure}}.
	\label{eq:spin-pressure-family}
\end{equation}

The spin pressure decays as $\cA_0^3$, which is much faster than the
amplitude.

\subsection{The Far-Field Structure}
\label{sec:family-far-field}

In the far field ($r \gg R$ or $|\xi| \gg L$), the dominant fields are:

\begin{enumerate}
	\item \textbf{Density:} $\delta \rho \sim e^{-r^2/(2R^2)}
	e^{-|\xi|/L}$.
	
	\item \textbf{Spin density:} $\mathcal{S} \sim e^{-r^2/(2R^2)}
	e^{-|\xi|/L}$.
	
	\item \textbf{Vorticity:} $\omega_\phi \sim e^{-r^2/(2R^2)}
	e^{-|\xi|/L}$.
	
	\item \textbf{Energy density:} $\delta e_{\mathrm{tot}} \sim
	e^{-r^2/(2R^2)} e^{-|\xi|/L}$.
\end{enumerate}

The spin pressure and other fields decay faster and are negligible in
the far field.

\subsection{The Multipole Moments}
\label{sec:family-multipoles}

The multipole moments are determined by the amplitude profile:

\subsubsection{Mass (Monopole)}

\begin{equation}
	\underbrace{M = 2\pi^2 \rho_{\mathrm{scale}} A_0 R^2 L}
	_{\text{mass}}.
	\label{eq:mass-family}
\end{equation}

\subsubsection{Charge (Monopole)}

\begin{equation}
	\underbrace{Q = \alpha_q \frac{\rho_{\mathrm{scale}}^3}{\rho^2
			\chi_\rho} \frac{\pi^2}{3} A_0^3 R^2 L}_{\text{charge}}.
	\label{eq:charge-family}
\end{equation}

\subsubsection{Magnetic Moment (Dipole)}

\begin{equation}
	\underbrace{\mathbf{m} = \cA_{\mathrm{scale}} \cdot 2\pi^2 A_0
		R^2 L \cos(\omega_{\mathrm{osc}} t) \hat{\phi}}
	_{\text{magnetic moment}}.
	\label{eq:magnetic-family}
\end{equation}

The magnetic moment oscillates with the breathing frequency.

\subsection{The Mass Scaling}
\label{sec:mass-scaling}

The mass scales as
\begin{equation}
	\underbrace{M \propto \gamma^2 \rho}_{\text{mass scaling}}.
	\label{eq:mass-scaling-family}
\end{equation}

This is the product of the amplitude squared (proportional to
$\gamma^2$) and the density scale. The mass is the dominant observable
of the breather family.

\subsection{The Parameter Space}
\label{sec:parameter-space}

The breather family lives in a four-dimensional parameter space:
\begin{equation}
	\underbrace{(A_0, R, \omega_{\mathrm{osc}}, c) \in \mathbb{R}^4}
	_{\text{parameter space}}.
	\label{eq:parameter-space}
\end{equation}

Each point in this space corresponds to a distinct breather solution
with a distinct mass, size, oscillation frequency, and speed.

\subsection{The Physical Interpretation}
\label{sec:family-interpretation}

The breather family is a continuum of localized, oscillatory, helical
solutions of the Madelung equations. Each member of the family is a
distinct physical object with:

\begin{enumerate}
	\item A distinct mass (determined by $A_0$ and $\rho$).
	\item A distinct size (determined by $R$).
	\item A distinct oscillation frequency (determined by
	$\omega_{\mathrm{osc}}$).
	\item A distinct speed (determined by $c$).
\end{enumerate}

The family is the natural generalization of the static soliton to the
oscillatory case.

% Concise-edit: redundant recap removed; detailed derivation above retained.

\part{Far-Field Structure}
% ============================================================

% ------------------------------------------------------------
\section{Multipole Moments}
\label{sec:farfield}
% ------------------------------------------------------------

The far-field structure of the breather is determined by its multipole
moments. In this section we compute the electric monopole, the magnetic
dipole, the gravitational monopole, and the gravitational quadrupole
from the breather solution.

\subsection{The Density Deviation}
\label{sec:density-deviation-farfield}

The density deviation is the source of the gravitational field:
\begin{equation}
	\underbrace{\delta \rho(r,\xi) = \rho_{\mathrm{scale}}
		\cA_0(r,\xi)}_{\text{density deviation}},
	\label{eq:density-deviation-later-2}
\end{equation}
where $\cA_0$ is the amplitude profile
\eqref{eq:amplitude-profile}.

With the amplitude profile:
\begin{equation}
	\underbrace{\delta \rho(r,\xi) = \rho_{\mathrm{scale}} A_0
		e^{-r^2/(2R^2)} \mathrm{sech}(\xi/L)}_{\text{density profile}}.
	\label{eq:density-profile}
\end{equation}

The density is localized in both $r$ and $\xi$.

\subsection{The Electric Monopole}
\label{sec:electric-monopole}

The electric charge is the integral of the spin pressure deviation:
\begin{equation}
	\underbrace{Q = \alpha_q \int \delta \Pi^{(s)} \, d^3x}
	_{\text{electric charge}},
	\label{eq:electric-charge}
\end{equation}
where $\alpha_q$ is the charge coupling and $\delta \Pi^{(s)}$ is the
spin pressure deviation.

From the spin pressure \eqref{eq:spin-pressure-family}:
\begin{equation}
	\underbrace{\delta \Pi^{(s)} = \frac{\rho_{\mathrm{scale}}^3
			\cA_0^3}{\rho^2 \chi_\rho}}_{\text{spin pressure}}.
	\label{eq:spin-pressure}
\end{equation}

The charge integral becomes
\begin{equation}
	\underbrace{Q = \alpha_q \frac{\rho_{\mathrm{scale}}^3}{\rho^2
			\chi_\rho} \int \cA_0^3 \, d^3x}_{\text{charge integral}}.
	\label{eq:charge-integral}
\end{equation}

The volume element in helical coordinates is
\begin{equation}
	\underbrace{d^3x = 2\pi L \, r \, dr \, dx}_{\text{volume element}},
	\label{eq:volume-element}
\end{equation}
where $x = \xi/L$.

The integral evaluates to
\begin{equation}
	\underbrace{\int \cA_0^3 \, d^3x = \frac{\pi^2}{3} A_0^3 R^2 L}
	_{\text{charge integral}}.
	\label{eq:charge-integral-value}
\end{equation}

The electric charge is
\begin{equation}
	\boxed{
		\underbrace{Q = \alpha_q \frac{\rho_{\mathrm{scale}}^3}{\rho^2
				\chi_\rho} \frac{\pi^2}{3} A_0^3 R^2 L}_{\text{electric monopole}}.
	}
	\label{eq:electric-monopole}
\end{equation}

This is the electric monopole moment. It is determined by the
amplitude scale, the density scale, and the charge coupling.

\subsection{The Magnetic Dipole}
\label{sec:magnetic-dipole}

The magnetic moment is the integral of the spin density:
\begin{equation}
	\underbrace{\mathbf{m} = \int \mathcal{S} \, d^3x}
	_{\text{magnetic moment}}.
	\label{eq:magnetic-moment}
\end{equation}

From the spin density \eqref{eq:spin-density-family}:
\begin{equation}
	\underbrace{\mathcal{S}(r,\xi,t) = \cA_{\mathrm{scale}}
		\cA_0(r,\xi) \cos(\omega_{\mathrm{osc}} t) \hat{\phi}}
	_{\text{spin density}}.
	\label{eq:spin-density-later-2}
\end{equation}

The magnetic moment becomes
\begin{equation}
	\underbrace{\mathbf{m} = \cA_{\mathrm{scale}} \cos(\omega_{\mathrm{osc}} t)
		\int \cA_0 \, d^3x \, \hat{\phi}}_{\text{magnetic integral}}.
	\label{eq:magnetic-integral}
\end{equation}

The integral evaluates to
\begin{equation}
	\underbrace{\int \cA_0 \, d^3x = 2\pi^2 A_0 R^2 L}
	_{\text{magnetic integral}}.
	\label{eq:magnetic-integral-value}
\end{equation}

The magnetic moment is
\begin{equation}
	\boxed{
		\underbrace{\mathbf{m} = 2\pi^2 \cA_{\mathrm{scale}} A_0
			R^2 L \cos(\omega_{\mathrm{osc}} t) \hat{\phi}}
		_{\text{magnetic dipole}}.
	}
	\label{eq:magnetic-dipole}
\end{equation}

The magnetic moment oscillates with the breathing frequency. This is a
prediction of the model: the magnetic dipole of the breather is
time-dependent.

\subsection{The Gravitational Monopole}
\label{sec:gravitational-monopole}

The gravitational mass is the integral of the total energy density:
\begin{equation}
	\underbrace{M = \frac{1}{c_{\mathrm{vortex}}^2} \int
		\delta e_{\mathrm{tot}} \, d^3x}_{\text{gravitational mass}}.
	\label{eq:gravitational-mass}
\end{equation}

The energy density is dominated by the rest mass:
\begin{equation}
	\underbrace{\delta e_{\mathrm{tot}} \approx \delta \rho
		c_{\mathrm{vortex}}^2}_{\text{rest mass}}.
	\label{eq:rest-mass}
\end{equation}

The mass integral becomes
\begin{equation}
	\underbrace{M = \rho_{\mathrm{scale}} \int \cA_0 \, d^3x}
	_{\text{mass integral}}.
	\label{eq:mass-integral}
\end{equation}

Using the same integral as before:
\begin{equation}
	\boxed{
		\underbrace{M = 2\pi^2 \rho_{\mathrm{scale}} A_0 R^2 L}
		_{\text{gravitational monopole}}.
	}
	\label{eq:gravitational-monopole}
\end{equation}

This is the gravitational monopole moment. It is determined by the
amplitude scale and the density scale.

\subsection{The Gravitational Quadrupole}
\label{sec:gravitational-quadrupole}

The gravitational quadrupole moment is
\begin{equation}
	\underbrace{Q_{ij} = \int (3 x_i x_j - r^2 \delta_{ij})
		\delta \rho \, d^3x}_{\text{quadrupole moment}}.
	\label{eq:quadrupole}
\end{equation}

For the breather, the density is elongated along the helical direction.
The dominant quadrupole component is
\begin{equation}
	\underbrace{Q_{\xi\xi} = \int (3 \xi^2 - r^2) \delta \rho
		\, d^3x}_{\text{helical quadrupole}}.
	\label{eq:helical-quadrupole}
\end{equation}

With the density profile:
\begin{equation}
	\underbrace{Q_{\xi\xi} = \rho_{\mathrm{scale}} A_0 \int (3 L^2 x^2
		- r^2) e^{-r^2/(2R^2)} \mathrm{sech}(x) \, d^3x}
	_{\text{quadrupole integral}}.
	\label{eq:quadrupole-integral}
\end{equation}

The integral evaluates to
\begin{equation}
	\boxed{
		\underbrace{Q_{\xi\xi} = \rho_{\mathrm{scale}} A_0 \left(
			3 \pi^2 L^2 - 2 \pi^2 R^2 \right) R^2 L}
		_{\text{gravitational quadrupole}}.
	}
	\label{eq:gravitational-quadrupole}
\end{equation}

The quadrupole moment is non-zero because the breather is elongated
along the helical direction.

\subsection{The Multipole Hierarchy}
\label{sec:multipole-hierarchy}

The multipole moments are summarized in Table~\ref{tab:multipoles}.

\begin{table}[htbp]
	\centering
	\begin{tabular}{@{}lll@{}}
		\toprule
		Moment & Expression & Source \\
		\midrule
		Electric monopole & $Q = \alpha_q \frac{\rho_{\mathrm{scale}}^3}
		{\rho^2 \chi_\rho} \frac{\pi^2}{3} A_0^3 R^2 L$ & Spin pressure \\
		Magnetic dipole & $\mathbf{m} = 2\pi^2 \cA_{\mathrm{scale}}
		A_0 R^2 L \cos(\omega_{\mathrm{osc}} t) \hat{\phi}$ & Spin density \\
		Gravitational monopole & $M = 2\pi^2 \rho_{\mathrm{scale}}
		A_0 R^2 L$ & Energy density \\
		Gravitational quadrupole & $Q_{\xi\xi} = \rho_{\mathrm{scale}}
		A_0 (3\pi^2 L^2 - 2\pi^2 R^2) R^2 L$ & Elongation \\
		\bottomrule
	\end{tabular}
	\caption{The multipole moments of the helical breather.}
	\label{tab:multipoles}
\end{table}

The far field is dominated by the monopole moments (electric and
gravitational) at large distances. The magnetic dipole is the leading
time-dependent moment.

\subsection{The Far-Field Potentials}
\label{sec:far-field-potentials}

The far-field potentials are determined by the multipole moments:

\subsubsection{Electric Potential}

\begin{equation}
	\underbrace{\Phi_e(r) = \frac{Q}{4\pi\epsilon_0 r}}
	_{\text{electric potential}}.
	\label{eq:electric-potential}
\end{equation}

\subsubsection{Magnetic Field}

\begin{equation}
	\underbrace{\mathbf{B}(r) = \frac{\mu_0}{4\pi} \frac{3(\mathbf{m}
			\cdot \hat{\mathbf{r}})\hat{\mathbf{r}} - \mathbf{m}}{r^3}}
	_{\text{magnetic field}}.
	\label{eq:magnetic-field}
\end{equation}

\subsubsection{Gravitational Potential}

\begin{equation}
	\underbrace{\Phi_g(r) = -\frac{GM}{r}}_{\text{gravitational potential}}.
	\label{eq:gravitational-potential}
\end{equation}

The gravitational potential has the standard $1/r$ form. There is no
Yukawa screening.

\subsection{The Gyromagnetic Ratio}
\label{sec:gyromagnetic-ratio}

The gyromagnetic ratio is the ratio of the magnetic moment to the
spin angular momentum:
\begin{equation}
	\underbrace{\frac{|\mathbf{m}|}{S} = \frac{2\pi^2
			\cA_{\mathrm{scale}} A_0 R^2 L}{\hbar/2}}_{\text{gyromagnetic ratio}}.
	\label{eq:gyromagnetic-ratio}
\end{equation}

With the spin calibration $S = \hbar/2$:
\begin{equation}
	\underbrace{\frac{|\mathbf{m}|}{S} = \frac{4\pi^2
			\cA_{\mathrm{scale}} A_0 R^2 L}{\hbar}}_{\text{gyromagnetic ratio}}.
	\label{eq:gyromagnetic-ratio-value}
\end{equation}

For the electron, the gyromagnetic ratio is $\mu_B/(\hbar/2) = 2e/m_e$.
The model reproduces this by construction, through the magnetic moment
constraint.

\subsection{The Charge-to-Mass Ratio}
\label{sec:charge-to-mass}

The charge-to-mass ratio is
\begin{equation}
	\underbrace{\frac{Q}{M} = \alpha_q \frac{\rho_{\mathrm{scale}}^2}
		{\rho^2 \chi_\rho} \frac{A_0^2}{3}}_{\text{charge-to-mass ratio}}.
	\label{eq:charge-to-mass}
\end{equation}

For the electron, the charge-to-mass ratio is $e/m_e$. The model
reproduces this by construction, through the charge and mass
constraints.

% Concise-edit: redundant recap removed; detailed derivation above retained.

\section{The Gravitational Sector}
\label{sec:gravity}
% ------------------------------------------------------------

The gravitational sector of the model is described by the dictionary
that associates the total energy density with the gravitational
potential. In this section we examine the Poisson equation, the
far-field potential, and the role of the gravitational constant $G$.

\subsection{The Poisson Equation}
\label{sec:poisson-equation}

The gravitational potential is sourced by the total energy density:
\begin{equation}
	\underbrace{\nabla^2 \Phi_g = \frac{4\pi G}{c_{\mathrm{vortex}}^2}
		e_{\mathrm{tot}}}_{\text{Poisson equation}},
	\label{eq:poisson}
\end{equation}
where $G$ is the gravitational constant and $e_{\mathrm{tot}}$ is the
total energy density.

In the far field, the energy density is dominated by the rest mass:
\begin{equation}
	\underbrace{e_{\mathrm{tot}} \approx \rho_{\mathrm{scale}}
		\cA_0 c_{\mathrm{vortex}}^2}_{\text{rest mass}}.
	\label{eq:energy-density-rest}
\end{equation}

The Poisson equation becomes
\begin{equation}
	\underbrace{\nabla^2 \Phi_g = 4\pi G \rho_{\mathrm{scale}}
		\cA_0}_{\text{simplified Poisson}}.
	\label{eq:poisson-simplified}
\end{equation}

This is the standard Poisson equation with the density deviation as
the source.

\subsection{The Far-Field Potential}
\label{sec:far-field-potential}

For a localized source, the far-field potential is
\begin{equation}
	\underbrace{\Phi_g(r) = -\frac{GM}{r}}_{\text{far-field potential}},
	\label{eq:far-field-potential}
\end{equation}
where $M$ is the gravitational monopole moment
\eqref{eq:gravitational-monopole}:
\begin{equation}
	\underbrace{M = 2\pi^2 \rho_{\mathrm{scale}} A_0 R^2 L}
	_{\text{gravitational mass}}.
	\label{eq:gravitational-mass-later-2}
\end{equation}

The far-field potential is the standard Newtonian potential. It has the
$1/r$ behavior characteristic of a point mass.

\subsection{The Multipole Expansion}
\label{sec:multipole-expansion}

The full gravitational potential is
\begin{equation}
	\underbrace{\Phi_g(\mathbf{r}) = -G \int \frac{\rho_{\mathrm{scale}}
			\cA_0(\mathbf{r}')}{|\mathbf{r} - \mathbf{r}'|} \, d^3x'}
	_{\text{Newtonian potential}}.
	\label{eq:newtonian-potential}
\end{equation}

For $r \gg R$ and $r \gg L$, the multipole expansion is
\begin{equation}
	\underbrace{\Phi_g(\mathbf{r}) = -G \left[ \frac{M}{r} +
		\frac{\mathbf{d} \cdot \hat{\mathbf{r}}}{r^2} + \frac{1}{2}
		\frac{Q_{ij} \hat{r}_i \hat{r}_j}{r^3} + \ldots \right]}
	_{\text{multipole expansion}},
	\label{eq:multipole-expansion}
\end{equation}
where $M$ is the total mass (monopole), $\mathbf{d}$ is the dipole
moment, and $Q_{ij}$ is the quadrupole moment.

\subsection{The Vanishing Dipole}
\label{sec:vanishing-dipole}

The dipole moment is
\begin{equation}
	\underbrace{\mathbf{d} = \int \mathbf{x} \, \delta \rho \, d^3x}
	_{\text{dipole moment}}.
	\label{eq:dipole-moment}
\end{equation}

For the symmetric breather, the density is even in $\xi$ and even in
$r$. The integral of $\mathbf{x} \, \delta \rho$ vanishes by symmetry:
\begin{equation}
	\boxed{
		\underbrace{\mathbf{d} = 0}_{\text{vanishing dipole}}.
	}
	\label{eq:vanishing-dipole}
\end{equation}

The far field is therefore dominated by the monopole and quadrupole
terms.

\subsection{The Gravitational Quadrupole}
\label{sec:gravitational-quadrupole-later-2}

The quadrupole moment is
\begin{equation}
	\underbrace{Q_{ij} = \int (3 x_i x_j - r^2 \delta_{ij})
		\delta \rho \, d^3x}_{\text{quadrupole moment}}.
	\label{eq:quadrupole-moment}
\end{equation}

For the breather, the density is elongated along the helical direction.
The dominant component is
\begin{equation}
	\underbrace{Q_{\xi\xi} = \int (3 \xi^2 - r^2) \delta \rho
		\, d^3x}_{\text{helical quadrupole}}.
	\label{eq:helical-quadrupole-later-2}
\end{equation}

From the numerical solution:
\begin{equation}
	\underbrace{Q_{\xi\xi} = \rho_{\mathrm{scale}} A_0 \left(
		3 \pi^2 L^2 - 2 \pi^2 R^2 \right) R^2 L}
	_{\text{quadrupole value}}.
	\label{eq:quadrupole-value}
\end{equation}

The quadrupole moment is non-zero and provides the leading correction
to the monopole potential at large distances.

\subsection{No Yukawa Screening}
\label{sec:no-yukawa}

The gravitational potential is not screened. There is no Yukawa
suppression at large distances.

The reason is that the association
\begin{equation}
	\underbrace{\Phi_g \leftrightarrow e_{\mathrm{tot}}}
	_{\text{dictionary}},
	\label{eq:dictionary}
\end{equation}
is not a literal identification. The gravitational potential and the
energy density are distinct fields. The energy density sources the
potential; the potential is not the energy density.

If the association were a literal identification, the Poisson equation
would become
\begin{equation}
	\underbrace{\nabla^2 \Phi_g = \kappa^2 \Phi_g}_{\text{Yukawa equation}},
	\label{eq:yukawa-equation}
\end{equation}
with the Yukawa solution
\begin{equation}
	\underbrace{\Phi_g(r) = -\frac{GM}{r} e^{-\kappa r}}
	_{\text{Yukawa potential}}.
	\label{eq:yukawa-potential}
\end{equation}

This is not the correct interpretation. The gravitational potential is
not screened.

\subsection{The Gravitational Constant}
\label{sec:gravitational-constant}

The gravitational constant $G$ is a fundamental input of the model. It
is not derived from the gyron parameters.

The calibration condition is simply
\begin{equation}
	\boxed{
		\underbrace{G = G_{\mathrm{observed}} = 6.674 \times 10^{-11}
			\, \mathrm{m}^3 \mathrm{kg}^{-1} \mathrm{s}^{-2}}
		_{\text{gravitational constant}}.
	}
	\label{eq:gravitational-constant}
\end{equation}

This is not a constraint on the gyron parameters. It is a statement
that the model reproduces the observed gravitational field, given the
observed value of $G$.

\subsection{The Electric-Gravitational Hierarchy}
\label{sec:electric-gravitational-hierarchy}

The ratio of the electric to gravitational force between two electrons
is
\begin{equation}
	\underbrace{\frac{F_e}{F_g} = \frac{e^2}{4\pi\epsilon_0 G m_e^2}
		\approx 4.17 \times 10^{42}}_{\text{electric-gravitational hierarchy}}.
	\label{eq:hierarchy}
\end{equation}

This hierarchy is enforced by the charge and mass constraints. It is
not derived from the spin sector. The model reproduces the hierarchy
but does not explain it.

\subsection{Comparison with the Standard Model}
\label{sec:gravity-comparison}

In the Standard Model, $G$ is a fundamental constant that is not
derived from other parameters. The Standard Model does not explain why
gravity is so weak compared to the other forces.

The gyron model is in the same position. It reproduces the
electric-gravitational hierarchy but does not explain it.

% Concise-edit: redundant recap removed; detailed derivation above retained.

\section{Coupling Equations}
\label{sec:coupling-equations}
% ------------------------------------------------------------

The model parameters are constrained by the requirement that the
breather reproduces the observed properties of the electron. In this
section we state the coupling equations and determine the combinations
of parameters that they fix.

\subsection{Speed of Light}
\label{sec:c-speed}

The emergent speed of light is the vortex-spin wave speed:
\begin{equation}
	\boxed{
		\underbrace{c_{\mathrm{vortex}} = \frac{\sqrt{\chi_\rho}}{\tau_s}}
		_{\text{speed of light}}.
	}
	\label{eq:c-constraint}
\end{equation}

With the measured value $c = 2.998 \times 10^8$ m/s, this fixes
\begin{equation}
	\underbrace{\frac{\chi_\rho}{\tau_s^2} = c^2}_{\text{fixed combination}}.
	\label{eq:c-fixed}
\end{equation}

\textbf{This fixes one combination: $\chi_\rho/\tau_s^2$.}

\subsection{Mass}
\label{sec:mass-constraint}

The gravitational mass is the integral of the total energy density:
\begin{equation}
	\underbrace{M = 2\pi^2 \rho_{\mathrm{scale}} A_0 R^2 L}
	_{\text{gravitational mass}}.
	\label{eq:mass-constraint}
\end{equation}

Setting $M = m_e$ gives
\begin{equation}
	\boxed{
		\underbrace{\rho_{\mathrm{scale}} A_0 R^2 L =
			\frac{m_e}{2\pi^2}}_{\text{mass constraint}}.
	}
	\label{eq:mass-constraint-explicit}
\end{equation}

\textbf{This fixes one combination: $\rho_{\mathrm{scale}} A_0 R^2 L$.}

\subsection{Charge}
\label{sec:charge-constraint}

The electric charge is
\begin{equation}
	\underbrace{Q = \alpha_q \frac{\rho_{\mathrm{scale}}^3}{\rho^2
			\chi_\rho} \frac{\pi^2}{3} A_0^3 R^2 L}_{\text{electric charge}}.
	\label{eq:charge-constraint}
\end{equation}

Setting $Q = e$ gives
\begin{equation}
	\boxed{
		\underbrace{\alpha_q \frac{\rho_{\mathrm{scale}}^3}{\rho^2
				\chi_\rho} A_0^3 R^2 L =
			\frac{3e}{\pi^2}}_{\text{charge constraint}}.
	}
	\label{eq:charge-constraint-explicit}
\end{equation}

Using the mass constraint to eliminate $\rho_{\mathrm{scale}} A_0 R^2 L$:
\begin{equation}
	\underbrace{\alpha_q \frac{\rho_{\mathrm{scale}}^2}{\rho^2
			\chi_\rho} A_0^2 = \frac{3e}{\pi^2} \cdot
		\frac{2\pi^2}{m_e}}_{\text{simplified charge}}.
	\label{eq:charge-simplified}
\end{equation}

\textbf{This fixes one combination: $\alpha_q
	\rho_{\mathrm{scale}}^2 A_0^2/(\rho^2 \chi_\rho)$.}

\subsection{Spin}
\label{sec:spin-constraint}

The spin angular momentum is
\begin{equation}
	\underbrace{S = 4 L \frac{\rho \chi_\rho}{\lambda \tau_a} I_S}
	_{\text{spin angular momentum}},
	\label{eq:spin-constraint}
\end{equation}
where
\begin{equation}
	\underbrace{I_S = \int_0^\infty \int_{-\infty}^\infty \cA_0(r,x)
		\, r \, dr \, dx}_{\text{spin integral}}.
	\label{eq:spin-integral}
\end{equation}

Setting $S = \hbar/2$ gives
\begin{equation}
	\boxed{
		\underbrace{\frac{\rho \chi_\rho}{\lambda \tau_a} =
			\frac{\hbar}{8 L I_S}}_{\text{spin constraint}}.
	}
	\label{eq:spin-constraint-explicit}
\end{equation}

\textbf{This fixes one combination:
	$\rho \chi_\rho/(\lambda \tau_a)$.}

\subsection{Planck's Constant}
\label{sec:planck-constraint}

The effective diffusion coefficient is calibrated to Planck's constant:
\begin{equation}
	\underbrace{D_{\mathrm{eff}} = D_s + c_2 D_a \mathcal{Q}_{\Hh}
		= -\frac{\hbar}{2M}}_{\text{Planck constraint}}.
	\label{eq:planck-constraint}
\end{equation}

With $M = m_e$:
\begin{equation}
	\boxed{
		\underbrace{D_s + c_2 D_a \mathcal{Q}_{\Hh} =
			-\frac{\hbar}{2m_e}}_{\text{Planck constraint}}.
	}
	\label{eq:planck-constraint-explicit}
\end{equation}

\textbf{This fixes one combination: $D_s + c_2 D_a
	\mathcal{Q}_{\Hh}$.}

\subsection{Magnetic Moment}
\label{sec:magnetic-constraint}

The magnetic moment is
\begin{equation}
	\underbrace{|\mathbf{m}| = 2\pi^2 \cA_{\mathrm{scale}} A_0 R^2 L
		\cos(\omega_{\mathrm{osc}} t)}_{\text{magnetic moment}}.
	\label{eq:magnetic-constraint}
\end{equation}

Time-averaging over one oscillation period:
\begin{equation}
	\underbrace{\langle |\mathbf{m}| \rangle = \frac{4}{\pi}
		\cA_{\mathrm{scale}} \pi^2 A_0 R^2 L}_{\text{time-averaged}}.
	\label{eq:magnetic-averaged}
\end{equation}

Setting $\langle |\mathbf{m}| \rangle = \mu_B = e\hbar/(2m_e)$ gives
\begin{equation}
	\boxed{
		\underbrace{\cA_{\mathrm{scale}} A_0 R^2 L =
			\frac{\mu_B \pi}{4 \pi^2}}_{\text{magnetic constraint}}.
	}
	\label{eq:magnetic-constraint-explicit}
\end{equation}

Simplifying:
\begin{equation}
	\underbrace{\cA_{\mathrm{scale}} A_0 R^2 L =
		\frac{e \hbar}{8 \pi m_e}}_{\text{magnetic constraint}}.
	\label{eq:magnetic-constraint-simplified}
\end{equation}

\textbf{This fixes one combination: $\cA_{\mathrm{scale}} A_0 R^2 L$.}

\subsection{Summary of Coupling Equations}
\label{sec:coupling-summary}

The coupling equations are summarized in
Table~\ref{tab:coupling-equations}.

\begin{table}[htbp]
	\centering
	\begin{tabular}{@{}lll@{}}
		\toprule
		Constraint & Equation & Fixes \\
		\midrule
		Speed of light & Eq.~\eqref{eq:c-constraint} &
		$\chi_\rho/\tau_s^2$ \\
		Mass & Eq.~\eqref{eq:mass-constraint-explicit} &
		$\rho_{\mathrm{scale}} A_0 R^2 L$ \\
		Charge & Eq.~\eqref{eq:charge-constraint-explicit} &
		$\alpha_q \rho_{\mathrm{scale}}^2 A_0^2/(\rho^2
		\chi_\rho)$ \\
		Spin & Eq.~\eqref{eq:spin-constraint-explicit} &
		$\rho \chi_\rho/(\lambda \tau_a)$ \\
		Planck & Eq.~\eqref{eq:planck-constraint-explicit} &
		$D_s + c_2 D_a \mathcal{Q}_{\Hh}$ \\
		Magnetic moment & Eq.~\eqref{eq:magnetic-constraint-simplified}
		& $\cA_{\mathrm{scale}} A_0 R^2 L$ \\
		\bottomrule
	\end{tabular}
	\caption{The coupling equations of the model.}
	\label{tab:coupling-equations}
\end{table}

Six coupling equations fix six combinations of parameters.

\subsection{The Gauge-Invariant Combinations}
\label{sec:gauge-invariant-combinations}

Combining the coupling equations with the gauge freedom, the
gauge-invariant combinations are:

\begin{align}
	\underbrace{\frac{\chi_\rho}{c^2 \tau_s^2} = 1}
	_{\text{speed of light}}, \\
	\underbrace{\frac{\rho_{\mathrm{scale}} A_0 R^2 L}{m_e} =
		\frac{1}{2\pi^2}}_{\text{mass}}, \\
	\underbrace{\frac{\alpha_q \rho_{\mathrm{scale}}^2 A_0^2}{\rho^2
			\chi_\rho} = \frac{3e}{\pi^2 m_e}}_{\text{charge}}, \\
	\underbrace{\frac{\rho \chi_\rho}{\lambda \tau_a} \frac{8 L
			I_S}{\hbar} = 1}_{\text{spin}}, \\
	\underbrace{\frac{D_s + c_2 D_a \mathcal{Q}_{\Hh}}{\hbar/(2m_e)}
		= -1}_{\text{Planck}}, \\
	\underbrace{\frac{\cA_{\mathrm{scale}} A_0 R^2 L}{\mu_B \pi /
			(4\pi^2)} = 1}_{\text{magnetic}}.
	\label{eq:gauge-invariant-combinations}
\end{align}

These are the constraints that the model must satisfy.

\subsection{The Residual Freedom}
\label{sec:residual-freedom}

After applying the six coupling equations and the gauge freedom, the
model retains a residual freedom. The residual parameters are:

\begin{enumerate}
	\item \textbf{The ratio $\tau_a/\tau_s$:} the nematic and spin
	relaxation times.
	
	\item \textbf{The ratio $c_4/c_2$:} the higher cross-section
	moment.
	
	\item \textbf{The ratio $D_a/D_s$:} the nematic and isotropic
	diffusion coefficients.
	
	\item \textbf{The overall scale:} the gauge freedom.
\end{enumerate}

The residual freedom is four-dimensional.

\subsection{The Predictions}
\label{sec:predictions}

Given the coupling equations, the model predicts:

\begin{table}[htbp]
	\centering
	\begin{tabular}{@{}lll@{}}
		\toprule
		Observable & Predicted Value & Status \\
		\midrule
		Electron mass & $m_e$ & By construction \\
		Electron charge & $e$ & By construction \\
		Electron spin & $\hbar/2$ & By construction \\
		Magnetic moment & $\mu_B$ & By construction \\
		Speed of light & $c$ & By construction \\
		Planck's constant & $\hbar$ & By construction \\
		Gravitational constant & $G$ & Input \\
		\bottomrule
	\end{tabular}
	\caption{Predictions of the model.}
	\label{tab:predictions}
\end{table}

The model reproduces the observed values of the electron properties
and the speed of light by construction. The gravitational constant $G$
remains a fundamental input.

\subsection{The Residual Freedom and Predictions}
\label{sec:residual-predictions}

The residual freedom corresponds to the following physical predictions:

\begin{enumerate}
	\item \textbf{Mass spectrum:} the quantized masses $M_n = n M_1$
	depend on the winding number $n$.
	
	\item \textbf{Magnetic moment:} the anomalous moment $a_e$ depends
	on the ratio $c_4/c_2$.
	
	\item \textbf{Gravitational wave emission:} the quadrupole
	radiation depends on the ratio $\tau_a/\tau_s$.
	
	\item \textbf{Charge quantization:} the charge is continuous unless
	a topological invariant is imposed.
\end{enumerate}

% Concise-edit: redundant recap removed; detailed derivation above retained.

\section{Residual Freedom}
\label{sec:residual-freedom-section}
% ------------------------------------------------------------

After applying the six coupling equations and the gauge freedom, the
model retains a residual freedom. In this section we identify the
residual parameters, determine their physical significance, and
discuss the predictions that depend on them.

\subsection{The Determined Combinations}
\label{sec:determined-combinations}
The six calibrated combinations are collected in
Table~\ref{tab:coupling-equations}; they are not independent predictions.

\subsection{The Free Parameters}
\label{sec:free-parameters-residual}

The model has twelve gyron parameters. Six combinations are fixed,
leaving six effective free parameters. One is removed by the gauge
freedom, leaving \textbf{five effective free parameters}.

The residual parameters are:

\begin{enumerate}
	\item $\tau_a/\tau_s$ (relaxation ratio):
	the ratio of the nematic and spin relaxation times.
	
	\item $D_a/D_s$ (diffusion ratio):
	the ratio of the nematic and isotropic diffusion coefficients.
	
	\item $c_4/c_2$ (shape ratio):
	the ratio of the higher cross-section moments.
	
	\item $\lambda$ (Jeffery parameter):
	the aspect ratio of the gyron.
	
	\item $\mathcal{Q}_{\Hh}$ (alignment scalar):
	the projection of the nematic onto the spin direction.
\end{enumerate}

These five parameters are not fixed by the electron properties. They
determine the predictions of the model for other observables.

\subsection{The Gauge Freedom}
\label{sec:gauge-freedom-residual}

The gauge freedom is the residual scale invariance:
\begin{equation}
	\underbrace{m \to \lambda m, \quad I_\rho \to \lambda I_\rho,
		\quad \tau_s \to \lambda \tau_s, \quad \chi_\rho \to \lambda
		\chi_\rho}_{\text{gauge transformation}}.
	\label{eq:gauge-transformation-residual}
\end{equation}

Under this transformation, the physical predictions are invariant:
\begin{equation}
	\underbrace{m_e, \, e, \, \mu_e, \, a_e, \, G, \, c, \,
		M_{\mathrm{active}} \to \text{unchanged}}_{\text{invariance}}.
	\label{eq:gauge-invariance-residual}
\end{equation}

The gauge freedom reflects the fact that the microscopic parameters
are not directly observable; only the macroscopic soliton properties
are.

\subsection{The Gauge-Invariant Combinations}
\label{sec:gauge-invariant-residual}
The six invariant combinations are given explicitly in
Eq.~\eqref{eq:gauge-invariant-combinations}.

\subsection{The Physical Predictions}
\label{sec:physical-predictions}

The residual freedom determines the predictions of the model for
observables beyond the electron properties:

\begin{enumerate}
	\item \textbf{Mass spectrum:} the quantized masses $M_n = n M_1$
	depend on the winding number $n$. The quantum of mass is
	$M_1 = 2\pi \hbar/(v_0 W)$.
	
	\item \textbf{Magnetic moment anomaly:} the anomalous moment
	$a_e = (g-2)/2$ depends on the ratio $c_4/c_2$.
	
	\item \textbf{Gravitational wave emission:} the quadrupole
	radiation depends on the ratio $\tau_a/\tau_s$.
	
	\item \textbf{Charge quantization:} the charge is continuous
	unless a topological invariant is imposed.
	
	\item \textbf{Hadron spectrum:} the higher breathers may
	correspond to hadrons.
\end{enumerate}

\subsection{The Mass Spectrum}
\label{sec:mass-spectrum-residual}

The mass spectrum is quantized by the vacuum-matching boundary
condition:
\begin{equation}
	\underbrace{M_n = n M_1}_{\text{mass spectrum}},
	\qquad
	\underbrace{M_1 = \frac{2\pi \hbar}{v_0 W}}_{\text{mass quantum}}.
	\label{eq:mass-spectrum-residual}
\end{equation}

With the numerical values:
\begin{equation}
	\underbrace{M_1 = 3.86 m_e}_{\text{mass quantum}}.
	\label{eq:mass-quantum-value}
\end{equation}

The quantized masses are:
\begin{table}[htbp]
	\centering
	\begin{tabular}{@{}ll@{}}
		\toprule
		$n$ & $M_n/m_e$ \\
		\midrule
		1 & 3.86 \\
		2 & 7.72 \\
		3 & 11.58 \\
		4 & 15.44 \\
		5 & 19.30 \\
		\bottomrule
	\end{tabular}
	\caption{Quantized masses.}
	\label{tab:quantized-masses}
\end{table}

The mass spectrum is linear in the winding number $n$.

\subsection{The Anomalous Moment}
\label{sec:anomalous-moment}

The anomalous moment $a_e = (g-2)/2$ depends on the ratio $c_4/c_2$.
The leading-order contribution is
\begin{equation}
	\underbrace{a_e \sim \frac{c_4}{c_2}}_{\text{anomalous moment}}.
	\label{eq:anomalous-moment}
\end{equation}

The observed value is $a_e = 1.16 \times 10^{-3}$. This fixes
\begin{equation}
	\underbrace{\frac{c_4}{c_2} \sim 10^{-3}}_{\text{shape ratio}}.
	\label{eq:shape-ratio}
\end{equation}

The shape ratio is one of the residual parameters.

\subsection{The Gravitational Wave Emission}
\label{sec:gravitational-wave}

The breather emits gravitational waves through its oscillating
quadrupole moment. The power emitted is
\begin{equation}
	\underbrace{P_{\mathrm{GW}} = \frac{G}{5 c^5} \left\langle
		\dddot{Q}_{ij} \dddot{Q}_{ij} \right\rangle}
	_{\text{gravitational wave power}}.
	\label{eq:gravitational-wave-power}
\end{equation}

The quadrupole oscillates with frequency $2\omega_{\mathrm{osc}}$. The
emission depends on the ratio $\tau_a/\tau_s$.

\subsection{Charge Quantization}
\label{sec:charge-quantization-residual}

The charge is continuous in the current model:
\begin{equation}
	\underbrace{Q = \alpha_q \frac{\rho_{\mathrm{scale}}^3}{\rho^2
			\chi_\rho} \frac{\pi^2}{3} A_0^3 R^2 L}_{\text{continuous charge}}.
	\label{eq:continuous-charge}
\end{equation}

Quantization would require a topological invariant, such as quantized
helicity or circulation. This is an open problem.

\subsection{The Hadron Spectrum}
\label{sec:hadron-spectrum}

The higher breathers ($n > 1$) may correspond to hadrons. The mass
quantum $M_1 = 3.86 m_e \approx 1.97$ MeV is much smaller than the
pion mass (140 MeV). The hadron spectrum would require a different
calibration.

Alternatively, the hadrons may correspond to breathers on different
wave modes (sound, nematic, hybrid). This is a direction for future
work.

\subsection{What \texorpdfstring{$G$}{G} Remains}
\label{sec:g-remains}

The gravitational constant $G$ remains a fundamental input:
\begin{equation}
	\underbrace{G = G_{\mathrm{observed}} = 6.674 \times 10^{-11}
		\, \mathrm{m}^3 \mathrm{kg}^{-1} \mathrm{s}^{-2}}
	_{\text{fundamental input}}.
	\label{eq:g-input}
\end{equation}

It is not derived from the gyron parameters. The model reproduces the
observed gravitational field, given the observed value of $G$.

% Concise-edit: redundant recap removed; detailed derivation above retained.

\part{Discussion}
% ============================================================

% ------------------------------------------------------------
\section{What the Model Explains}
\label{sec:explains}
% ------------------------------------------------------------

The quaternionic gyron spin fluid is a classical kinetic theory with
no quantum mechanics, no relativity, and no gauge structure built in.
Yet it reproduces, at the level of equation form, several features of
the Standard Model and its gravitational extension. We list here what
the model explains, with the appropriate caveats.

\subsection{Emergent Relativistic Kinematics}
\label{sec:explains-relativity}

In the marginal regime
\begin{equation}
	\underbrace{v_{\mathrm{th}} = c_{\mathrm{vortex}}}
	_{\text{marginal condition}},
	\label{eq:explains-marginal}
\end{equation}
the transverse vortex-spin wave satisfies a Lorentz-invariant wave
equation:
\begin{equation}
	\underbrace{\frac{\partial^2 \cS_\perp}{\partial t^2}
		- c_{\mathrm{vortex}}^2 \nabla^2 \cS_\perp = 0}
	_{\text{Lorentz-invariant wave equation}}.
	\label{eq:explains-wave}
\end{equation}

The emergent structure includes:

\begin{enumerate}
	\item \textbf{A Minkowski metric:}
	\begin{equation}
		\underbrace{ds^2 = -c_{\mathrm{vortex}}^2 dt^2 + dx^2 + dy^2
			+ dz^2}_{\text{emergent metric}}.
		\label{eq:explains-metric}
	\end{equation}
	
	\item \textbf{Lorentz transformations:} for a soliton moving with
	velocity $U$ along the $z$-axis, the energy and momentum
	transform as
	\begin{equation}
		\underbrace{E' = \gamma(E - UP), \qquad P' = \gamma\left(P -
			\frac{UE}{c_{\mathrm{vortex}}^2}\right)}_{\text{Lorentz
				transformation}}.
		\label{eq:explains-lorentz}
	\end{equation}
	
	\item \textbf{The mass-shell condition:}
	\begin{equation}
		\underbrace{E^2 - P^2 c_{\mathrm{vortex}}^2 = M^2
			c_{\mathrm{vortex}}^4}_{\text{mass-shell}}.
		\label{eq:explains-mass-shell}
	\end{equation}
	
	\item \textbf{Length contraction and time dilation:}
	\begin{equation}
		\underbrace{L_z' = \frac{L_z}{\gamma}, \qquad \Delta t' =
			\gamma \Delta t}_{\text{contraction and dilation}}.
		\label{eq:explains-contraction}
	\end{equation}
\end{enumerate}

This is a genuine structural emergence: the Galilean microscopic
theory produces a Lorentz-invariant effective description in a
specific limit. The model explains \emph{why} there is a speed limit
(it is the vortex-spin wave speed) but not \emph{why} that speed has
the measured value $3 \times 10^8$ m/s.

\subsection{Maxwell-Form Electromagnetic Sector}
\label{sec:explains-maxwell}

The vorticity equation has the same form as the magnetic induction
equation:
\begin{equation}
	\underbrace{\partial_t \omega_i + u_j \partial_j \omega_i =
		\omega_j \partial_j u_i + D_\omega \nabla^2 \omega_i}
	_{\text{magnetic induction form}},
	\label{eq:explains-induction}
\end{equation}
where $D_\omega$ is an effective magnetic diffusivity.

The antisymmetric stress divergence has the same form as the Lorentz
force:
\begin{equation}
	\underbrace{-\frac{1}{\rho} \partial_j P_{ij}^{\mathrm{anti}}
		= -\frac{\eta_r}{\rho} \partial_j \omega_{ij}
		+ \frac{1}{2\rho\tau_s} (\nabla \times \cS)_i + \ldots}
	_{\text{Lorentz force form}}.
	\label{eq:explains-lorentz-force}
\end{equation}

This gives an equation-level analogue of electromagnetism. The model
explains why the electromagnetic equations have their specific form
(they are the transport equations of the spin-vorticity sector) but
not the magnitude of the electric charge or the structure of the gauge
group.

\subsection{Charge Conjugation and Particle-Antiparticle Structure}
\label{sec:explains-charge-conjugation}

The discrete $\mathbb{Z}_2$ symmetry
\begin{equation}
	\underbrace{\cS_\phi \to -\cS_\phi, \qquad \omega_\phi \to
		-\omega_\phi, \qquad \delta \Pi^{(s)} \to -\delta \Pi^{(s)}}
	_{\text{charge conjugation}},
	\label{eq:explains-c}
\end{equation}
gives two branches of the helical soliton with:

\begin{enumerate}
	\item \textbf{Opposite charge-like signature:}
	\begin{equation}
		\underbrace{Q \to -Q}_{\text{charge flip}}.
		\label{eq:explains-charge-flip}
	\end{equation}
	
	\item \textbf{Opposite magnetic moment:}
	\begin{equation}
		\underbrace{\mathbf{m} \to -\mathbf{m}}_{\text{moment flip}}.
		\label{eq:explains-moment-flip}
	\end{equation}
	
	\item \textbf{Equal gravitational mass:}
	\begin{equation}
		\underbrace{M_{e^+} = M_{e^-}}_{\text{mass equality}}.
		\label{eq:explains-mass-equality}
	\end{equation}
\end{enumerate}

The mass equality is enforced by the Bernoulli compensation condition
along helical streamlines. This is a classical, dynamical analogue of
charge conjugation, and it provides a mechanism by which particle and
antiparticle can have equal mass without being imposed by hand.

\subsection{The Quantum Potential Term}
\label{sec:explains-quantum-potential}

The spin amplitude $\cA = |\cQ|$ satisfies a conservation law with a
Laplacian term:
\begin{equation}
	\underbrace{\partial_t \cA + \partial_j (u_j \cA) =
		D_{\mathrm{eff}} \nabla^2 \cA + \ldots}
	_{\text{amplitude equation}},
	\label{eq:explains-amplitude}
\end{equation}
where the effective diffusion coefficient is
\begin{equation}
	\underbrace{D_{\mathrm{eff}} = D_s + c_2 D_a \mathcal{Q}_{\Hh}}
	_{\text{effective diffusion}}.
	\label{eq:explains-d-eff}
\end{equation}

When $D_{\mathrm{eff}} < 0$, the Laplacian term has the same sign as
the dispersive term appearing in the Madelung hydrodynamic
representation of quantum mechanics:
\begin{equation}
	\underbrace{D_{\mathrm{eff}} \nabla^2 \cA \leftrightarrow
		-\frac{\hbar^2}{2m^2} \nabla \left( \frac{\nabla^2 \sqrt{\cA}}
		{\sqrt{\cA}} \right)}_{\text{Madelung form}}.
	\label{eq:explains-madelung}
\end{equation}

The sign condition forces a strongly prolate gyron. This is a
structural emergence of a quantum-potential-like term from a classical
kinetic theory.

\subsection{Soliton Taxonomy}
\label{sec:explains-taxonomy}

The balance equations admit a one-parameter family of helical vortex
solitons parametrized by the pitch $l$. The solution set is partitioned
into islands by the sign changes of the $l$-dependent corrections. Each
island corresponds to a distinct family of solitons with different mass
and moment.

The taxonomy is a qualitative prediction of the model, independent of
the detailed soliton profile.

\subsection{The Hadron-Like Mass Spectrum}
\label{sec:explains-mass-spectrum}

The mass spectrum is quantized by the vacuum-matching boundary
condition:
\begin{equation}
	\underbrace{M_n = n M_1}_{\text{mass spectrum}},
	\qquad
	\underbrace{M_1 = \frac{2\pi \hbar}{v_0 W}}_{\text{mass quantum}}.
	\label{eq:explains-mass-spectrum}
\end{equation}

The quantized masses are:
\begin{table}[htbp]
	\centering
	\begin{tabular}{@{}ll@{}}
		\toprule
		$n$ & $M_n/m_e$ \\
		\midrule
		1 & 3.86 \\
		2 & 7.72 \\
		3 & 11.58 \\
		4 & 15.44 \\
		5 & 19.30 \\
		\bottomrule
	\end{tabular}
	\caption{Quantized masses of the helical breather.}
	\label{tab:explains-quantized-masses}
\end{table}

The mass spectrum is linear in the winding number $n$. This matches
the qualitative structure of the hadron spectrum, which is
approximately linear in the radial quantum number.

\subsection{The Electric-Gravitational Hierarchy}
\label{sec:explains-hierarchy}

The model reproduces the ratio
\begin{equation}
	\underbrace{\frac{F_e}{F_g} = \frac{e^2}{4\pi\epsilon_0 G m_e^2}
		\approx 4.17 \times 10^{42}}_{\text{electric-gravitational
			hierarchy}}
	\label{eq:explains-hierarchy}
\end{equation}
through the charge and mass constraints. The hierarchy is enforced by
the charge constraint, not derived from the spin sector. The model
reproduces the hierarchy but does not explain why the electric and
gravitational couplings have such different strengths.

\subsection{The Far-Field Multipole Structure}
\label{sec:explains-multipoles}

The breather has the far-field structure of:

\begin{enumerate}
	\item \textbf{Electric monopole:}
	\begin{equation}
		\underbrace{Q = \alpha_q \frac{\rho_{\mathrm{scale}}^3}{\rho^2
				\chi_\rho} \frac{\pi^2}{3} A_0^3 R^2 L}
		_{\text{electric monopole}}.
		\label{eq:explains-electric}
	\end{equation}
	
	\item \textbf{Magnetic dipole:}
	\begin{equation}
		\underbrace{\mathbf{m} = 2\pi^2 \cA_{\mathrm{scale}} A_0 R^2 L
			\cos(\omega_{\mathrm{osc}} t) \hat{\phi}}
		_{\text{magnetic dipole}}.
		\label{eq:explains-magnetic}
	\end{equation}
	
	\item \textbf{Gravitational monopole:}
	\begin{equation}
		\underbrace{M = 2\pi^2 \rho_{\mathrm{scale}} A_0 R^2 L}
		_{\text{gravitational monopole}}.
		\label{eq:explains-gravitational}
	\end{equation}
	
	\item \textbf{Gravitational quadrupole:}
	\begin{equation}
		\underbrace{Q_{\xi\xi} = \rho_{\mathrm{scale}} A_0 \left(
			3 \pi^2 L^2 - 2 \pi^2 R^2 \right) R^2 L}
		_{\text{gravitational quadrupole}}.
		\label{eq:explains-quadrupole}
	\end{equation}
\end{enumerate}

This is the standard far-field multipole structure of a localized,
spinning, charged object.

\subsection{Summary}
\label{sec:explains-summary}

The model explains:

\begin{enumerate}
	\item \textbf{Emergent relativistic kinematics:} the transverse
	vortex-spin wave satisfies a Lorentz-invariant wave equation, with
	the speed of light identified as the vortex-spin wave speed.
	
	\item \textbf{Maxwell-form electromagnetic sector:} the vorticity
	equation and the antisymmetric stress divergence have the same
	form as the magnetic induction equation and the Lorentz force.
	
	\item \textbf{Charge conjugation and particle-antiparticle
		structure:} the discrete $\mathbb{Z}_2$ symmetry gives two
	branches with opposite charge, opposite magnetic moment, and
	equal gravitational mass.
	
	\item \textbf{The quantum potential term:} the spin amplitude
	equation contains a Laplacian term with the same functional form
	as the Madelung quantum potential.
	
	\item \textbf{Soliton taxonomy:} the balance equations admit a
	one-parameter family of helical vortex solitons.
	
	\item \textbf{Hadron-like mass spectrum:} the mass spectrum is
	linear in the winding number, matching the qualitative structure
	of the hadron spectrum.
	
	\item \textbf{Electric-gravitational hierarchy:} the ratio
	$F_e/F_g \approx 4.17 \times 10^{42}$ is reproduced by the
	charge and mass constraints.
	
	\item \textbf{Far-field multipole structure:} the breather has
	the standard multipole structure of a localized, spinning,
	charged object.
\end{enumerate}

What the model does not explain is the subject of the next section.

% ------------------------------------------------------------
\section{What the Model Does Not Explain}
\label{sec:not-explains}
% ------------------------------------------------------------

Equally important is what the model fails to explain. We list these
here to delimit the scope of the claims.

\subsection{The Values of Fundamental Constants}
\label{sec:not-values}

The model does not predict the numerical values of:

\begin{itemize}
	\item the speed of light $c$,
	\item the elementary charge $e$,
	\item the electron mass $m_e$,
	\item the gravitational constant $G$,
	\item Planck's constant $\hbar$.
\end{itemize}

Each of these is used as a calibration of the kinetic parameters, not
derived from first principles. The model is a predictive machine: once
the kinetic parameters are fixed by the measured anchors, the model
predicts other observables as functions of the residual freedom. It
does not predict the anchors themselves.

\subsection{Charge Quantization}
\label{sec:not-charge}

Charge is continuous in the current model. The charge-like signature
\begin{equation}
	\underbrace{Q = \alpha_q \int \delta \Pi^{(s)} \, d^3x}
	_{\text{continuous charge}}
	\label{eq:not-charge}
\end{equation}
is a real number that varies continuously with the soliton parameters.

Quantization would require a topological invariant, such as quantized
helicity or circulation. This is not present in the model as
constructed.

\subsection{Quantum Mechanics}
\label{sec:not-quantum}

The model does not derive quantum mechanics. It produces a term of the
same functional form as the Madelung quantum potential in the spin
amplitude equation, but the term is classical in origin: it comes from
the spin flux closure, not from a complex wavefunction.

The following quantum features are absent:

\begin{itemize}
	\item the Schr\"odinger equation,
	\item the superposition principle,
	\item quantum interference,
	\item the Born rule,
	\item spin-statistics.
\end{itemize}

The model is a classical analogue of quantum kinematics, not a
derivation of quantum mechanics.

\subsection{General Relativity}
\label{sec:not-gr}

The gravitational sector is treated in the Newtonian limit, with the
Poisson equation for the gravitational potential. The model does not
derive:

\begin{itemize}
	\item the Einstein field equations,
	\item the equivalence principle,
	\item the geometric structure of general relativity.
\end{itemize}

A full relativistic treatment of gravity is beyond the scope of the
current model.

\subsection{The Standard Model Gauge Structure}
\label{sec:not-sm}

The model reproduces the equation form of electromagnetism but not the
gauge structure. It does not contain:

\begin{itemize}
	\item $SU(2) \times U(1)$ symmetry,
	\item the weak interaction,
	\item the strong interaction,
	\item the three generations of fermions.
\end{itemize}

The structural analogy with the Standard Model (charge conjugation,
families, emergent gauge structure) is suggestive but not a
derivation.

\subsection{The Value of \texorpdfstring{$\hbar$}{hbar}}
\label{sec:not-hbar}

The effective quantum potential coefficient $|D_{\mathrm{eff}}|$
depends on the kinetic parameters through the combination
\begin{equation}
	\underbrace{D_{\mathrm{eff}} = D_s + c_2 D_a \mathcal{Q}_{\Hh}}
	_{\text{effective diffusion}}.
	\label{eq:not-d-eff}
\end{equation}

The model does not predict this combination; it uses the measured
$\hbar$ to fix it. The origin of $\hbar$ remains unexplained.

\subsection{The Lepton Mass Hierarchy}
\label{sec:not-lepton}

The mass spectrum of the breather is linear in the winding number:
\begin{equation}
	\underbrace{M_n = n M_1}_{\text{linear spectrum}},
	\label{eq:not-linear}
\end{equation}
whereas the observed lepton hierarchy is exponential:
\begin{equation}
	\underbrace{\frac{m_\mu}{m_e} \approx 207, \qquad
		\frac{m_\tau}{m_e} \approx 3477}_{\text{exponential hierarchy}}.
	\label{eq:not-exponential}
\end{equation}

The linear spectrum does not match the lepton hierarchy. The breather
spectrum may correspond instead to the hadron spectrum, which is
approximately linear in the radial quantum number.

% Concise-edit: redundant recap removed; detailed derivation above retained.

\section{Open Problems}
\label{sec:open}
% ------------------------------------------------------------

Several problems remain open and would sharpen the model if resolved.

\subsection{Charge Quantization}
\label{sec:open-charge}

A topological invariant that quantizes the charge would require either:

\begin{enumerate}
	\item a quantized helicity, with the charge proportional to the
	helicity of the soliton:
	\begin{equation}
		\underbrace{Q \propto \mathcal{H} = \int \mathbf{u} \cdot
			(\nabla \times \mathbf{u}) \, d^3x}
		_{\text{quantized helicity}},
		\label{eq:open-helicity}
	\end{equation}
	or
	\item a quantized circulation, with the charge proportional to
	the winding number of the spin field:
	\begin{equation}
		\underbrace{Q \propto \oint \cS \cdot d\boldsymbol{\ell}}
		_{\text{quantized circulation}}.
		\label{eq:open-circulation}
	\end{equation}
\end{enumerate}

Neither is present in the current model. A topological extension would
require additional structure in the collision integral or the boundary
conditions.

\subsection{Origin of \texorpdfstring{$\hbar$}{hbar}}
\label{sec:open-hbar}

The effective quantum potential coefficient
\begin{equation}
	\underbrace{|D_{\mathrm{eff}}| = \frac{\hbar}{2M}}
	_{\text{calibrated value}}
	\label{eq:open-d-eff}
\end{equation}
is a combination of transport coefficients. Whether it can be derived
from a more fundamental principle (e.g., a minimum action, a
quantization condition, or a topological constraint) is an open
question. The model currently uses $\hbar$ as an input.

\subsection{Stability Spectrum}
\label{sec:open-stability}

The full stability analysis of the helical soliton requires:

\begin{enumerate}
	\item the reduced ODEs for the amplitudes $(R, L, \Gamma, U, S_0,
	\hat{a}_0)$ including the nematic terms,
	\item the fixed point of the reduced system,
	\item the eigenvalues of the Jacobian at the fixed point.
\end{enumerate}

This is a finite calculation but has not been carried out. The
stability spectrum would determine whether the soliton is
self-stabilizing or requires fine-tuning.

\subsection{Island Structure and Family Number}
\label{sec:open-islands}

The number and size of the islands of existence require the explicit
$l$-dependent corrections to the balance equations. Computing these
would determine whether the model predicts a single soliton family, or
multiple families analogous to the generations of the Standard Model.

\subsection{Maxwell Consistency}
\label{sec:open-maxwell}

The equation-level analogy with Maxwell's equations is at the level of
the induction equation and the Lorentz force. Whether the full Maxwell
system, including the displacement current and the gauge structure,
emerges from the gyron model at higher order is an open question.

\subsection{Cherenkov-Like Instability}
\label{sec:open-cherenkov}

A soliton moving faster than the local average gyron velocity can
radiate energy into the gyron population through a Cherenkov-like
mechanism. The radiative decay rate and its effect on the soliton
lifetime have not been computed.

\subsection{The Residual Gauge Freedom}
\label{sec:open-gauge}

After calibration by the measured anchors, the model retains about
three residual degrees of freedom: the relaxation time ratio
$\tau_a/\tau_s$, the transport ratio $D_a/D_s$, and the shape focusing
ratio $c_4/c_2$. These parametrize the family of solitons that
reproduce the anchors. Identifying observables that are sensitive to
each residual parameter, and checking whether the predictions are
consistent with measurement, is the natural next step for the model.

\subsection{Connection to Existing Literature}
\label{sec:open-literature}

The model draws on several existing frameworks:

\begin{itemize}
	\item Cosserat and Eringen micropolar theory for the couple-stress
	structure of the antisymmetric stress,
	\item Relativistic spin hydrodynamics for the covariant
	formulation of spin transport,
	\item Doi-Edwards and Hess orientational kinetic theory for the
	nematic order parameter and anisotropic closure,
	\item Jeffery orbits for the orientation dynamics of axisymmetric
	bodies in a flow,
	\item Burgers vortex for the radial structure of the core,
	\item Helical vortex and travelling-helix solutions for the
	background geometry and existence precedents.
\end{itemize}

A more thorough comparison with these frameworks, and an
identification of which features are new and which are inherited,
would clarify the model's place in the literature.

% Concise-edit: redundant recap removed; detailed derivation above retained.

\section{Conclusion}
\label{sec:conclusion}
% ------------------------------------------------------------

We have developed a classical kinetic theory for a spin fluid of
axisymmetric rigid bodies with quaternionic orientation, with a closed
moment hierarchy that retains explicit dependence on the shape
parameters $c_2, c_4$ and the Jeffery parameter $\lambda$.

The model produces:

\begin{enumerate}
	\item A dictionary mapping the fluid fields onto electromagnetic
	and gravitational analogues, at the level of equation form.
	
	\item An emergent speed of light identified with the vortex-spin
	wave speed, in a marginal regime of thermal synchronization.
	
	\item An emergent Lorentz invariance for the transverse spin
	sector, with the full structure of Minkowski spacetime and the
	mass-shell condition.
	
	\item A term of the same functional form as the Madelung quantum
	potential in the spin amplitude equation, with a sign condition
	that forces a strongly prolate gyron.
	
	\item A family of helical vortex solitons, classified by charge
	conjugation, parity, and family number.
	
	\item A discrete $\mathbb{Z}_2$ symmetry relating particle and
	antiparticle branches, with opposite charge and equal mass
	enforced by the Bernoulli compensation condition.
\end{enumerate}

The model is classical throughout. It does not derive quantum
mechanics, general relativity, or the Standard Model. It is a toy
model in the analogue-gravity tradition: a classical system whose
effective description reproduces the kinematics and structural features
of a target physical field, for the purpose of understanding the
emergence of that field from a simpler substrate.

The predictions of the model are specific and falsifiable. The sign
condition for the quantum potential term forces a strongly prolate
gyron. The discrete symmetry relating particle and antiparticle
requires the collision integral to be chirally symmetric. The mass
equality between the two branches requires the odd corrections to the
Bernoulli relation to be below the observed $10^{-9}$ bound. Each of
these is a checkable condition, and each could falsify the model.

The most pressing open problems are:

\begin{enumerate}
	\item the stability analysis of the helical soliton,
	\item the computation of the island structure,
	\item the quantitative estimate of the mass difference between
	the two branches,
	\item the origin of charge quantization,
	\item the derivation of $\hbar$ from a more fundamental principle.
\end{enumerate}

These would sharpen the model from a structural description to a
quantitative one, and would allow direct comparison with the observed
particle spectrum.

\subsection{Future Directions: Solitons on Other Wave Modes}
\label{sec:conclusion-future}

The helical vortex solitons analyzed in this paper are constructed on
the transverse vortex-spin wave of the spin sector. The linear spectrum
contains additional branches, however: the longitudinal sound mode,
the spin relaxation mode, the nematic relaxation mode, and the hybrid
mode that appears only on the soliton background.

Each of these branches is a candidate carrier for a nonlinear localized
structure. If the nonlinear terms in the transport equations balance
the dispersion of a given linear branch, a soliton can form around
that branch: a localized, travelling, particle-like structure whose
mass, spin, and stability are determined by the nonlinear balance
rather than put in by hand.

The gyron model therefore suggests a general program: for each branch
of the linear spectrum, ask whether a nonlinear soliton exists, and if
so, what its quantum numbers are.

Different solitons on different branches would have different masses,
spins, and coupling strengths. A sound soliton, for example, would
carry no electromagnetic charge (the sound mode decouples from the
spin sector at leading order) but would carry energy and hence a
gravitational mass; a soliton on the hybrid mode would carry both spin
and nematic character and would couple to the electromagnetic sector
through the spin channel.

The spectrum of solitons across the branches could therefore provide a
structural analogue of a zoo of elementary particles, distinguished by
which wave mode they are built on, rather than by being introduced as
separate fields.

We emphasize that this is a program, not a result. The nonlinear
soliton solutions have not been constructed, and whether they exist,
what their masses are, and whether they match any observed particle
spectrum are open questions. But the structure of the gyron model
suggests that the same kinetic framework that produces the helical
vortex solitons of this paper may also produce other particle-like
structures on the remaining branches of the spectrum, and that the
resulting catalogue could be compared systematically with the observed
particle content.

\subsection{Closing Remarks}
\label{sec:conclusion-closing}

The gyron spin fluid is a toy model. It is not a theory of everything.
It does not derive the Standard Model. But it does show that a
classical kinetic theory with quaternionic orientation can reproduce
the kinematics of relativistic spinning particles, the structure of
electromagnetism, the charge conjugation symmetry, and a
quantum-potential-like term, all within a single framework. The
quaternionic formulation is the minimal extension of the real
three-vector spin fluid that carries the phase, the $SU(2)$ structure,
and the chirality needed for the Standard Model.

The model is classical throughout. It uses the measured values of
$\hbar$, $c$, $G$, $m_e$, and $e$ as inputs, and predicts the electron
mass, charge, spin, and magnetic moment by construction. The
gravitational constant $G$ remains a fundamental input.

The most important lesson of the model is that the Madelung surface
--- the constraint surface on which the gyron dynamics reduces to the
Madelung form --- is a genuine sector of the classical kinetic theory.
On this surface, the spin amplitude and spin phase gradient satisfy
the Madelung equations, including a quantum-potential term whose
coefficient is calibrated to the physical value of $\hbar$. The
helical breathing soliton is the natural localized solution on this
surface.

Whether this is a step toward a deeper understanding of the emergence
of quantum mechanics from classical kinetics, or simply an elegant
analogy, remains to be seen. The model is offered as a tool for
exploring the question.

	% ============================================================
	\clearpage
\appendix
	% ============================================================
	% ============================================================
	\appendix
	% ============================================================
	
	% ------------------------------------------------------------
	\section{Glossary of Terms}
	\label{app:glossary}
	% ------------------------------------------------------------
	
	This appendix defines the fields, coefficients, and parameters of
	the quaternionic gyron model in words. It is intended as a conceptual
	map for readers who want to understand what each quantity means, not
	just how it is denoted.
	
	\subsection{Fluid Fields}
	\label{app:glossary-fields}
	
	\begin{description}
		\item[$\rho$ (mass density).] The mass per unit volume of the
		gyron fluid. Conserved by the collision operator.
		
		\item[$\mathbf{u}$ (fluid velocity).] The mass-weighted mean
		velocity of the gyrons.
		
		\item[$\cQ$ (quaternionic spin density).] The mass-weighted
		intrinsic angular momentum per unit volume, expressed as a
		quaternion. Its norm $\cA = |\cQ|$ is the spin amplitude; its
		argument $U = \cQ/|\cQ|$ is the orientation.
		
		\item[$\cP$ (conjugate momentum density).] The conjugate momentum
		to the orientation quaternion.
		
		\item[$P_{ij}$ (translational pressure tensor).] The second moment
		of the velocity fluctuations about the mean.
		
		\item[$\Pi_{ab}$ (quaternionic spin pressure).] The second moment
		of the spin fluctuations about the mean spin.
		
		\item[$Q_{ia}$ (spin flux).] The correlation between velocity
		fluctuations and spin fluctuations.
	\end{description}
	
	\subsection{Transport Coefficients}
	\label{app:glossary-transport}
	
	\begin{description}
		\item[$D_s$ (spin diffusion coefficient).] The coefficient of the
		isotropic part of the spin flux.
		
		\item[$D_a$ (nematic spin transport coefficient).] The coefficient
		of the anisotropic part of the spin flux.
		
		\item[$\eta$ (shear viscosity).] The coefficient of the traceless
		symmetric part of the viscous stress.
		
		\item[$\eta_r$ (rotational viscosity).] The coefficient of the
		antisymmetric part of the viscous stress.
		
		\item[$\kappa_2$ (spin-nematic coupling).] The coefficient of the
		leading nematic correction to the antisymmetric stress.
	\end{description}
	
	\subsection{Relaxation Times}
	\label{app:glossary-relaxation}
	
	\begin{description}
		\item[$\tau$ (translational relaxation time).] The timescale on
		which the translational pressure tensor relaxes to its isotropic
		value.
		
		\item[$\tau_s$ (spin relaxation time).] The timescale on which
		the spin density relaxes to its slaved value.
		
		\item[$\tau_a$ (orientational relaxation time).] The timescale on
		which the nematic order parameter relaxes to isotropy.
	\end{description}
	
	\subsection{Shape Parameters}
	\label{app:glossary-shape}
	
	\begin{description}
		\item[$a, b$ (gyron dimensions).] The length and transverse scale
		of the gyron.
		
		\item[$\lambda$ (Jeffery parameter).] The dimensionless parameter
		$\lambda = (a^2 - b^2)/(a^2 + b^2)$ that controls the alignment
		of an axisymmetric body in a shear flow.
		
		\item[$c_2$ (quadrupole cross-section moment).] The coefficient of
		$P_2$ in the Legendre expansion of the hard-body cross section.
		
		\item[$\mathcal{Q}_{\Hh}$ (alignment scalar).] The projection of
		the nematic order parameter onto the spin direction.
	\end{description}
	
	\subsection{Soliton Parameters}
	\label{app:glossary-soliton}
	
	\begin{description}
		\item[$R$ (core radius).] The characteristic radial scale of the
		soliton.
		
		\item[$L$ (helical packet length).] The characteristic length of
		the soliton along the helical direction.
		
		\item[$\Gamma$ (circulation).] The line integral of the azimuthal
		velocity around the core.
		
		\item[$l$ (helical pitch).] The pitch of the helical flow.
		
		\item[$\xi$ (helical coordinate).] The coordinate $\xi = z -
		l\theta$ that is constant along helical streamlines.
	\end{description}
	
	\subsection{Derived Quantities}
	\label{app:glossary-derived}
	
	\begin{description}
		\item[$c_{\mathrm{vortex}}$ (vortex-spin wave speed).] The speed
		of the transverse spin wave, $c_{\mathrm{vortex}} =
		\sqrt{\chi_\rho}/\tau_s$. Identified with the emergent speed of
		light.
		
		\item[$D_{\mathrm{eff}}$ (effective Laplacian coefficient).] The
		coefficient of the $\nabla^2 \cA$ term in the spin amplitude
		equation.
		
		\item[$v_{\mathrm{th}}$ (thermal speed).] The typical gyron
		thermal speed.
		
		\item[$\chi_\rho$ (spin susceptibility).] The response of the
		spin density to an applied vorticity.
		
		\item[$e_{\mathrm{tot}}$ (total energy density).] The sum of
		translational, spin, and interaction energy densities.
	\end{description}
	
	\subsection{Anchors}
	\label{app:glossary-anchors}
	
	The measured quantities used to calibrate the kinetic parameters:
	
	\begin{table}[htbp]
		\centering
		\begin{tabular}{@{}lll@{}}
			\toprule
			Quantity & Symbol & Value \\
			\midrule
			Electron mass & $m_e$ & $9.109 \times 10^{-31}$ kg \\
			Elementary charge & $e$ & $1.602 \times 10^{-19}$ C \\
			Electron magnetic moment & $\mu_e$ & $-1.00116 \mu_B$ \\
			Anomalous moment & $a_e$ & $1.15965 \times 10^{-3}$ \\
			Gravitational constant & $G$ & $6.674 \times 10^{-11}$
			m$^3$kg$^{-1}$s$^{-2}$ \\
			Speed of light & $c$ & $2.998 \times 10^8$ m/s \\
			Planck's constant & $\hbar$ & $1.055 \times 10^{-34}$ J$\cdot$s \\
			\bottomrule
		\end{tabular}
		\caption{Measured anchors.}
		\label{tab:anchors}
	\end{table}
	
	% ------------------------------------------------------------
	\section{Quaternion Algebra}
	\label{app:quaternion}
	% ------------------------------------------------------------
	
	\subsection{Definition}
	\label{app:quaternion-definition}
	
	A quaternion is written as
	\begin{equation}
		q = q_0 + q_1 i + q_2 j + q_3 k,
		\label{eq:quaternion-def}
	\end{equation}
	with multiplication rules
	\begin{equation}
		i^2 = j^2 = k^2 = ijk = -1,
		\label{eq:quaternion-squares}
	\end{equation}
	\begin{equation}
		ij = k, \quad jk = i, \quad ki = j, \quad
		ji = -k, \quad kj = -i, \quad ik = -j.
		\label{eq:quaternion-products}
	\end{equation}
	
	\subsection{Conjugation and Norm}
	\label{app:quaternion-conjugation}
	
	The conjugate is
	\begin{equation}
		\bar{q} = q_0 - q_1 i - q_2 j - q_3 k,
		\label{eq:quaternion-conjugate}
	\end{equation}
	and the norm is
	\begin{equation}
		|q|^2 = q\bar{q} = q_0^2 + q_1^2 + q_2^2 + q_3^2.
		\label{eq:quaternion-norm}
	\end{equation}
	For unit quaternions, $|q| = 1$, so $q \in S^3$.
	
	\subsection{Rotation of Vectors}
	\label{app:quaternion-rotation}
	
	A vector $\mathbf{a} \in \Rr^3$ is rotated by $q$ to
	\begin{equation}
		\mathbf{a}' = q \, \mathbf{a} \, \bar{q},
		\label{eq:quaternion-rotation}
	\end{equation}
	where $\mathbf{a}$ is embedded as the pure imaginary quaternion
	$a_1 i + a_2 j + a_3 k$.
	
	\subsection{The Maurer-Cartan Form}
	\label{app:quaternion-maurer-cartan}
	
	The Maurer-Cartan form is
	\begin{equation}
		\omega_i = \bar{U} \partial_i U \in \mathfrak{su}(2).
		\label{eq:maurer-cartan-app}
	\end{equation}
	Its curvature is
	\begin{equation}
		F_{ij} = \partial_i \omega_j - \partial_j \omega_i
		+ [\omega_i, \omega_j].
		\label{eq:curvature-app}
	\end{equation}
	For a pure gauge configuration, $F_{ij} = 0$.
	
	\subsection{The Exponential Map}
	\label{app:quaternion-exponential}
	
	For a unit imaginary quaternion $\hat{n}$ (with $\hat{n}^2 = -1$),
	the exponential map is
	\begin{equation}
		e^{\theta \hat{n}} = \cos\theta + \hat{n} \sin\theta.
		\label{eq:exponential-map}
	\end{equation}
	This is the quaternionic analogue of the complex exponential.
	
	\subsection{The Spin Angular Momentum}
	\label{app:quaternion-spin}
	
	The spin angular momentum in quaternion form is
	\begin{equation}
		\sigma_i(q) = \mathrm{Im}\left( q \, I_\rho \boldsymbol{\omega} \,
		\bar{q} \right).
		\label{eq:spin-angular-momentum-app}
	\end{equation}
	
	% ------------------------------------------------------------
	\section{The Quaternionic Jeffery Equation}
	\label{app:jeffery}
	% ------------------------------------------------------------
	
	\subsection{The Jeffery Equation}
	\label{app:jeffery-equation}
	
	For a rigid axisymmetric body in a flow, the orientation evolves by
	\begin{equation}
		\dot{q} = \frac{1}{2} \Omega \, q + \lambda \mathcal{S} \, q,
		\label{eq:jeffery-app}
	\end{equation}
	where
	\begin{equation}
		\Omega = \Omega_{ij} e_i e_j, \qquad
		\mathcal{S} = S_{ij} e_i e_j,
		\label{eq:vorticity-strain-app}
	\end{equation}
	are the vorticity and strain-rate tensors embedded as pure imaginary
	quaternions, and
	\begin{equation}
		\lambda = \frac{(a/b)^2 - 1}{(a/b)^2 + 1}
		\label{eq:jeffery-parameter-app}
	\end{equation}
	is the Jeffery shape parameter.
	
	\subsection{The Vorticity and Strain Tensors}
	\label{app:jeffery-tensors}
	
	The vorticity and strain-rate tensors are
	\begin{equation}
		\Omega_{ij} = \frac{1}{2}(\partial_i u_j - \partial_j u_i), \qquad
		S_{ij} = \frac{1}{2}(\partial_i u_j + \partial_j u_i).
		\label{eq:vorticity-strain-def}
	\end{equation}
	
	\subsection{The Jeffery Coupling}
	\label{app:jeffery-coupling}
	
	The Jeffery coupling in the nematic balance equation is
	\begin{equation}
		\mathcal{J}_{ab} = \hat{\mathcal{A}}_{ak} \mathcal{S}_{kb}
		+ \hat{\mathcal{A}}_{bk} \mathcal{S}_{ka}
		- \tfrac{1}{2} \delta_{ab} \hat{\mathcal{A}}_{kl} \mathcal{S}_{kl}.
		\label{eq:jeffery-coupling-app}
	\end{equation}
	
	\subsection{The Vector Limit}
	\label{app:jeffery-vector}
	
	In the vector limit, the Jeffery equation reduces to
	\begin{equation}
		\dot{\hat{n}}_i = \Omega_{ij} \hat{n}_j + \lambda (\delta_{ij}
		- \hat{n}_i \hat{n}_j) S_{jk} \hat{n}_k.
		\label{eq:jeffery-vector}
	\end{equation}
	
	% ------------------------------------------------------------
	\section{Spectral ODEs and Matrix Elements}
	\label{app:spectral}
	% ------------------------------------------------------------
	
	\subsection{The Spectral Equations}
	\label{app:spectral-equations}
	
	The spectral equations are
	\begin{equation}
		(u_\xi - c) \partial_x A_m = \alpha \left( m A_{m-1}
		+ (2m+1) A_m + (m+1) A_{m+1} \right)
		- \beta A_m + \gamma \mathrm{sech}(x) \delta_{m0},
		\label{eq:spectral-app}
	\end{equation}
	where the dimensionless parameters are
	\begin{align}
		\tilde{u} &= \frac{u_\xi - c}{c_{\mathrm{vortex}}}, \\
		\alpha &= \frac{L D_s}{2 c_{\mathrm{vortex}}}, \\
		\beta &= \frac{L}{\tau_s c_{\mathrm{vortex}}}, \\
		\gamma &= \frac{L \mathcal{F}_0}{c_{\mathrm{vortex}}}.
		\label{eq:dimensionless-parameters-app}
	\end{align}
	
	\subsection{The Gaussian-Laguerre Functions}
	\label{app:spectral-gaussian-laguerre}
	
	The Gaussian-Laguerre functions are
	\begin{equation}
		\phi_n(r) = e^{-r^2/(2R^2)} L_n(r^2/R^2),
		\label{eq:gaussian-laguerre-app}
	\end{equation}
	where $L_n$ are Laguerre polynomials.
	
	\subsection{The Matrix Elements}
	\label{app:spectral-matrix}
	
	The matrix elements of the radial Laplacian are
	\begin{equation}
		\mathcal{L}_{mn} = -\frac{n+1}{2} \delta_{m,n+1}
		- \frac{2n+1}{2} \delta_{mn} - \frac{n}{2} \delta_{m,n-1}.
		\label{eq:matrix-elements-app}
	\end{equation}
	This is a tridiagonal matrix.
	
	\subsection{The Analytic Solution}
	\label{app:spectral-analytic}
	
	For the $m = 0$ mode, the analytic solution is
	\begin{equation}
		A_0(x) = \frac{\gamma}{\tilde{u}} \int_{-\infty}^{x}
		e^{-\beta(x - x')/\tilde{u}} \mathrm{sech}(x') \, dx'.
		\label{eq:analytic-solution-app}
	\end{equation}
	
	% ------------------------------------------------------------
	\section{Numerical Solution Details}
	\label{app:numerical}
	% ------------------------------------------------------------
	
	\subsection{The Sparse Linear System}
	\label{app:numerical-sparse}
	
	The spectral equations are discretized on a uniform grid in $x$ with
	central differences:
	\begin{equation}
		\partial_x A_m \approx \frac{A_{m,j+1} - A_{m,j-1}}{2 \Delta x},
		\label{eq:central-difference-app}
	\end{equation}
	where $j$ indexes the grid points.
	
	The resulting linear system is
	\begin{equation}
		\mathbf{M} \mathbf{A} = \mathbf{F},
		\label{eq:linear-system-app}
	\end{equation}
	where $\mathbf{M}$ is a sparse block-tridiagonal matrix, $\mathbf{A}$
	is the vector of mode amplitudes, and $\mathbf{F}$ is the driving
	vector.
	
	\subsection{The Dimensionless Parameters}
	\label{app:numerical-dimensionless}
	
	The dimensionless parameters are
	\begin{align}
		\tilde{u} &= 1.67 \times 10^{-2}, \\
		\alpha &= 9.65 \times 10^{-27}, \\
		\beta &= 3.33 \times 10^{-1}, \\
		\gamma &= 3.33 \times 10^{-2}.
		\label{eq:dimensionless-values-app}
	\end{align}
	
	\subsection{The Numerical Method}
	\label{app:numerical-method}
	
	The linear system is solved using a sparse linear solver. The grid
	parameters are:
	
	\begin{table}[htbp]
		\centering
		\begin{tabular}{@{}ll@{}}
			\toprule
			Parameter & Value \\
			\midrule
			Number of modes $N$ & 10 \\
			Number of grid points $J$ & 2000 \\
			Domain half-width $x_{\max}$ & 20.0 \\
			\bottomrule
		\end{tabular}
		\caption{Numerical parameters.}
		\label{tab:numerical-parameters}
	\end{table}
	
	\subsection{The Mode Structure}
	\label{app:numerical-mode}
	
	The mode structure is characterized by:
	
	\begin{enumerate}
		\item \textbf{Single-mode dominance:} the $m = 0$ mode carries
		essentially all of the amplitude.
		
		\item \textbf{Localization:} the mode amplitudes decay to zero
		at the boundaries.
		
		\item \textbf{sech-like profile:} the $m = 0$ mode has a
		sech-like profile.
	\end{enumerate}
	
	% ------------------------------------------------------------
	\section{Far-Field Multipole Expansion}
	\label{app:multipole}
	% ------------------------------------------------------------
	
	\subsection{The Multipole Expansion}
	\label{app:multipole-expansion}
	
	For a localized source, the far-field potential is
	\begin{equation}
		\Phi(\mathbf{r}) = -G \left[ \frac{M}{r} + \frac{\mathbf{d}
			\cdot \hat{\mathbf{r}}}{r^2} + \frac{1}{2}
		\frac{Q_{ij} \hat{r}_i \hat{r}_j}{r^3} + \ldots \right],
		\label{eq:multipole-expansion-app}
	\end{equation}
	where $M$ is the monopole, $\mathbf{d}$ is the dipole, and $Q_{ij}$ is
	the quadrupole.
	
	\subsection{The Monopole Moments}
	\label{app:multipole-monopole}
	
	The monopole moments are:
	
	\begin{itemize}
		\item \textbf{Gravitational mass:} $M = 2\pi^2
		\rho_{\mathrm{scale}} A_0 R^2 L$.
		\item \textbf{Electric charge:} $Q = \alpha_q
		\frac{\rho_{\mathrm{scale}}^3}{\rho^2 \chi_\rho}
		\frac{\pi^2}{3} A_0^3 R^2 L$.
	\end{itemize}
	
	\subsection{The Dipole Moment}
	\label{app:multipole-dipole}
	
	The dipole moment is
	\begin{equation}
		\mathbf{d} = \int \mathbf{x} \, \delta \rho \, d^3x.
		\label{eq:dipole-moment-app}
	\end{equation}
	For a symmetric breather, $\mathbf{d} = 0$.
	
	\subsection{The Quadrupole Moment}
	\label{app:multipole-quadrupole}
	
	The quadrupole moment is
	\begin{equation}
		Q_{ij} = \int (3 x_i x_j - r^2 \delta_{ij}) \delta \rho \, d^3x.
		\label{eq:quadrupole-moment-app}
	\end{equation}
	For the breather, the dominant component is
	\begin{equation}
		Q_{\xi\xi} = \rho_{\mathrm{scale}} A_0 \left( 3 \pi^2 L^2
		- 2 \pi^2 R^2 \right) R^2 L.
		\label{eq:quadrupole-value-app}
	\end{equation}
	
	% ------------------------------------------------------------
	\section{Cherenkov-Like Instability}
	\label{app:cherenkov}
	% ------------------------------------------------------------
	
	\subsection{The Threshold}
	\label{app:cherenkov-threshold}
	
	The threshold for radiative decay is
	\begin{equation}
		U_{\mathrm{crit}} = v_{\mathrm{th}} = \sqrt{\frac{k_B T}{m}}.
		\label{eq:cherenkov-threshold}
	\end{equation}
	
	Below this speed, the gyrons can "keep up" with the soliton. Above
	this speed, the soliton outruns the bulk of the population, creating
	a wake that carries energy away.
	
	\subsection{The Decay Rate}
	\label{app:cherenkov-decay}
	
	The radiative decay rate scales as
	\begin{equation}
		\Gamma_{\mathrm{rad}} \sim n \sigma v_{\mathrm{th}}
		\left( \frac{U^2}{v_{\mathrm{th}}^2} - 1 \right)^{3/2},
		\qquad U > v_{\mathrm{th}},
		\label{eq:cherenkov-decay}
	\end{equation}
	where $n$ is the gyron density and $\sigma$ is the collision cross
	section.
	
	\subsection{The Total Decay Rate}
	\label{app:cherenkov-total}
	
	The total decay rate is the sum of the intrinsic dissipation and the
	radiative contribution:
	\begin{equation}
		\Gamma_{\mathrm{total}} = \Gamma_{\mathrm{intrinsic}}
		+ \Gamma_{\mathrm{rad}}(U).
		\label{eq:cherenkov-total}
	\end{equation}
	
	\subsection{The Relation to the Emergent Speed Limit}
	\label{app:cherenkov-emergent}
	
	In the marginal regime, the Cherenkov threshold coincides with the
	vortex-spin wave speed:
	\begin{equation}
		U_{\mathrm{crit}} = v_{\mathrm{th}} = c_{\mathrm{vortex}}.
		\label{eq:cherenkov-emergent}
	\end{equation}
	
	The emergent Lorentz invariance therefore comes with a dynamical
	speed limit whose origin is the interaction between the soliton and
	the medium.
	% ============================================================
\begin{thebibliography}{99}

\bibitem{meno1991}
F. Meno, \emph{Physics Essays} \textbf{4}, 94 (1991).
% Full article title and page range require verification.

\bibitem{meno2000}
F. Meno, \emph{Cats, Atoms, Gyrons, Aether, and the Universe}
(AEtherpress, Pittsburgh, 2000).

\bibitem{saul2003}
L. Saul, ``Spin Waves as Metric in a Kinetic Space-Time,''
\emph{Physics Letters A} \textbf{316}, 385 (2003);
arXiv:gr-qc/0304063.

\bibitem{saul2015}
L. A. Saul, ``Statistical Models of Light with an Example in
Relativistic Dynamics,'' \emph{Electronic Journal of Theoretical Physics}
\textbf{12}, 205--210 (2015).

\bibitem{kyriazi2013}
H. Kyriazi, ``A Viable, Atomistic, Kinetic, Unified Aether Theory,''
viXra:1311.0060v4 (2013).

\bibitem{weng2025}
Z.-H. Weng, ``Gauge fields and four interactions in the
trigintaduonion spaces,'' \emph{Mathematical Methods in the Applied
Sciences} \textbf{48}, 590--604 (2025), doi:10.1002/mma.10345.

		
		\bibitem{gyron}
		[Gyron spin fluid paper.]
		
		\bibitem{cosserat1909}
		E. Cosserat and F. Cosserat,
		\emph{Th\'eorie des corps d\'eformables},
		Hermann, Paris (1909).
		
		\bibitem{curtiss1956}
		C. F. Curtiss,
		``Kinetic theory of nonspherical molecules,''
		\emph{J. Chem. Phys.} \textbf{24}, 225--241 (1956).
		
		\bibitem{eringen1966}
		A. C. Eringen,
		``Theory of micropolar fluids,''
		\emph{J. Math. Mech.} \textbf{16}, 1--18 (1966).
		
		\bibitem{florkowski2018}
		W. Florkowski, B. Friman, A. Jaiswal, and E. Speranza,
		``Relativistic fluid dynamics with spin,''
		\emph{Phys. Rev. C} \textbf{97}, 041901(R) (2018).
		
		\bibitem{barcelo2011}
		C. Barcelo, S. Liberati, and M. Visser,
		``Analogue gravity,''
		\emph{Living Rev. Relativ.} \textbf{14}, 3 (2011).
		
		\bibitem{unruh1981}
	W. G. Unruh,
	``Experimental black-hole evaporation?,''
	\emph{Phys. Rev. Lett.} \textbf{46}, 1351--1353 (1981).
	\newblock doi:10.1103/PhysRevLett.46.1351.

	\bibitem{visser1998}
	M. Visser,
	``Acoustic black holes: Horizons, ergospheres, and Hawking radiation,''
	\emph{Class. Quantum Grav.} \textbf{15}, 1767--1791 (1998).
	\newblock doi:10.1088/0264-9381/15/7/006.

	\bibitem{vissermolina2010}
	M. Visser and C. Molina-Par\'is,
	``Acoustic geometry for general relativistic barotropic irrotational fluid flow,''
	\emph{New J. Phys.} \textbf{12}, 095014 (2010).
	\newblock doi:10.1088/1367-2630/12/9/095014.

	\bibitem{jeffery1922}
		G. B. Jeffery,
		``The motion of ellipsoidal particles immersed in a viscous fluid,''
		\emph{Proc. R. Soc. Lond. A} \textbf{102}, 161--179 (1922).
		
		\bibitem{bgk1954}
		P. L. Bhatnagar, E. P. Gross, and M. Krook,
		``A model for collision processes in gases,''
		\emph{Phys. Rev.} \textbf{94}, 511--525 (1954).
		
		\bibitem{hess1976}
		S. Hess,
		``Fokker-Planck-equation approach to flow alignment in liquid crystals,''
		\emph{Z. Naturforsch. A} \textbf{31}, 1034--1037 (1976).
		
		\bibitem{doi1986}
		M. Doi and S. F. Edwards,
		\emph{The Theory of Polymer Dynamics},
		Clarendon Press, Oxford (1986).
		
		\bibitem{burgers1948}
		J. M. Burgers,
		``A mathematical model illustrating the theory of turbulence,''
		\emph{Adv. Appl. Mech.} \textbf{1}, 171--199 (1948).
		
		\bibitem{madelung1927}
		E. Madelung,
		``Quantentheorie in hydrodynamischer Form,''
		\emph{Z. Phys.} \textbf{40}, 322--326 (1927).
		
		\bibitem{takabayasi1952}
		T. Takabayasi,
		``On the formulation of quantum mechanics associated with classical pictures,''
		\emph{Prog. Theor. Phys.} \textbf{8}, 143--182 (1952).
		
		\bibitem{xu2022}
		Y. Xu, I. Delbende, and M. Rossi,
		``Quasi-equilibrium states for helical vortices with swirl,''
		\emph{J. Fluid Mech.} \textbf{944}, A24 (2022).
		
		\bibitem{davila2022}
		J. Davila, M. del Pino, M. Musso, and J. Wei,
		``Travelling helices and the vortex filament conjecture
		in the incompressible Euler equations,''
		\emph{Calc. Var. Partial Differential Equations} \textbf{61}, 119 (2022).
		
		\bibitem{furey2018}
		C. Furey,
		``Standard model physics from an algebra,''
		[reference].
		
		\bibitem{baez2002}
		J. Baez,
		``The octonions,''
		\emph{Bull. Amer. Math. Soc.} \textbf{39}, 145--205 (2002).
		
	\end{thebibliography}
	
\end{document}